% concatenated from Exchangeable_Cycles_ArXiv.tex
\documentclass[fleqn,11pt]{article}
\usepackage[margin=1in]{geometry}                % See geometry.pdf to learn the layout options. There are lots.
\geometry{letterpaper}
\usepackage{authblk}

\usepackage{easyReview}
\usepackage[dvipsnames]{xcolor}
%\renewcommand{\addColor}{\textcolor{magenta}}
%\setcitestyle{round}
%\usepackage[sort&compress,numbers,round]{natbib} % can use parenthesis citations w/ acs



% See https://creativecommons.org/licenses/ for a full explanation of the 
% Creative Commons licenses.
%
% We strongly recommend CC BY-ND or CC BY. Both these licenses allow others
% to freely distribute your work while giving credit to you. The difference
% between them is that CC BY-ND only allows distribution unchanged and in
% whole, while CC BY also allows remixing, tweaking and building upon your work.  

% TITLE
% If needed, include a line break (\\) at an appropriate place in the title.
% Footnotes on titles are not allowed.
%\title{Sorting out Clusters of Cycles in Graphs:\\
%Concepts and Algorithms}
%\title{Theory and algorithms for clusters of cycles in graphs for material networks}
%\date{\today}

% AUTHORS
%\author{Perrin E. Ruth}
%\email{pruth@umd.edu}
%\author{Maria K. Cameron}
%\affiliation[University of Maryland, College Park]
%{Department of Mathematics, University of Maryland, College Park}



% PACKAGES
\usepackage{amsmath,amsthm,amsfonts,amssymb}
\usepackage{graphicx}
\usepackage{hyperref}
% format of link color and inline citations
\hypersetup{colorlinks,linkcolor=blue,citecolor=blue}

% THEOREM ENVIRONMENTS
\newtheorem{definition}{Definition}
% share numbering between theorem, lemma, and corollary
\newtheorem{theorem}{Theorem}
\newtheorem{lemma}[theorem]{Lemma}
\newtheorem{corollary}[theorem]{Corollary}
\newtheorem{prop}{Proposition}

\usepackage{xcolor}
\usepackage{bm,soul}
\usepackage{cancel}
\usepackage{mathtools}
% Tikz!
%\usepackage{tikz}
%\usetikzlibrary{shapes,decorations,decorations.markings,arrows.meta,
%	backgrounds,calligraphy,
%	patterns,patterns.meta}
%\tikzset{
%	namednode/.style={
%		draw,
%		circle,
%		inner sep=0pt,
%		minimum size=.5cm},
%	regnode/.style={circle, draw, fill=black,
%		inner sep=0pt, minimum size = 5pt}
%	}
% algorithms
\usepackage[ruled,linesnumbered]{algorithm2e}
\SetKwComment{Comment}{/* }{ */}
% enumerate with less spaces and labels controllable
\usepackage{enumitem}
\newenvironment{myenumerate}
{ \begin{enumerate}
	\setlength{\itemsep}{0pt}
	\setlength{\parskip}{0pt}
	\setlength{\parsep}{0pt}       }
{ \end{enumerate}                  } 


% shorthands for annoying symbols to repeatedly write
% cycle bases
\newcommand\calB{\mathcal{B}}
\newcommand\calF{\mathcal{F}}
\newcommand\calM{\mathcal{M}}
% other cycle sets
\newcommand\calP{\mathcal{P}}     % polyhedron
\newcommand\calT{\mathcal{T}}     % trivial-interchangeability class - depricated
\newcommand\calS{\mathcal{S}}     % sli class
\newcommand\calX{\mathcal{X}}     % arbitrary
\newcommand\calY{\mathcal{Y}}
\newcommand\calZ{\mathcal{Z}}
\newcommand\CycSp{\mathcal{C}}    % sets of cycles
\newcommand\RelCyc{\mathcal{C_R}} % relevant cycles
% some operators
% expansion of cycle (arg. 2) w.r.t. basis (arg. 1)
\newcommand\expand[2]{\mathcal{E}_{#1}\left(#2\right)}   % symmetric difference for sets of cycles
\newcommand\expandd[2]{\mathcal{E}_{#1}(#2)}             % less horizontal spacing by not using left, right
\newcommand\symdiff{\mathop{\triangle}}            
\usepackage{scalerel}
\DeclareMathOperator*{\symdiffseries}{\scalerel*{\triangle}{\sum}}   
% polyhedron/short loop-interchangeability
\newcommand\pInt{ \stackrel{\tt pi}{\longleftrightarrow} }
\newcommand\slInt{ \stackrel{\tt sli}{\longleftrightarrow} }             
%% direct/trivial-interchangeability (depricated)
%\newcommand\DirEx{\sim_d}
%\newcommand\TrEx{\sim_t}%\overset{<}{\longleftrightarrow}}
% smaller cycles
\newcommand\smCyc{\textnormal{(smaller cycles)}}
% inner product, to avoid repeated langle typing
\newcommand\iprod[1]{\langle#1\rangle} 
\DeclareMathOperator\Span{span}                         % span in terms of a vector space

% Remove with ToC
% Note without this \ref{*section*} fails.
%\SectionNumbersOn

% footnotes to * and dagger, etc. so as not to conflict with references
\renewcommand*{\thefootnote}{\fnsymbol{footnote}}

\begin{document}
	
	{\title{Theory and algorithms for clusters of cycles in graphs for material networks}% Force line breaks with \\
		% \thanks{A footnote to the article title}%
		
		
		\author[1]{Perrin E. Ruth\thanks{pruth@umd.edu}}
		\author[1]{Maria Cameron\thanks{mariakc@umd.edu}}
		\affil[1]{\small{Department of Mathematics, University of Maryland, College Park, MD 20742, USA}}}
		\maketitle

%\pagecolor{lightgray} % night time reading mode


\begin{abstract}
	Analysis of complex networks, particularly material networks such as the carbon skeleton of hydrocarbons generated in hydrocarbon pyrolysis in carbon-rich systems, is essential for effectively describing, modeling, and predicting their features.  An important and the most challenging part of this analysis is the extraction and effective description of cycles, when many of them coalesce into complex clusters. A deterministic minimum cycle basis (MCB) is generally non-unique and biased to the vertex enumeration. The union of all MCBs, called the set of relevant cycles, is unique, but may grow exponentially with the graph size. To resolve these issues, we propose a method to sample an MCB uniformly at random. The output MCB is statistically well-defined, and its size is proportional to the number of edges. We review and advance the theory of graph cycles from previous works of Vismara, Gleiss et al., and Kolodzik et al. In particular, we utilize the polyhedron-interchangeability ({\tt pi}) and short loop-interchangeability ({\tt sli}) classes to partition the relevant cycles. We introduce a postprocessing step forcing pairwise intersections of relevant cycles to consist of a single path. This permits the definition of a dual graph whose nodes are cycles and edges connect pairs of intersecting cycles. The {\tt pi} classes identify building blocks for crystalline structures. The {\tt sli} classes group together sets of large redundant cycles. We present the application to an amorphous hydrocarbon network, where we (i) theorize how the number of relevant cycles may explode with system size and (ii) observe small polyhedral structures related to diamond.
%	We propose a method to statistically sample cycles from a chemical network: 
%	to sample a minimum cycle basis (MCB)~\cite{kavithaCycle2009,plotkinMathematical1971} from a graph uniformly at random.
%	A deterministic MCB is limited because it is not uniquely defined. A random MCB is statistically well-defined. 
%	Larger classes of cycles (e.g., the relevant cycles composed of the union of the MCBs of a graph~\cite{vismaraUnion1997}) are not polynomial with respect to graph size.
%	A random MCB grows linearly with the number of edges. 
%	We also contribute theory for how pairs of MCB cycles can intersect, and a postprocessing step to to force these pairwise intersections into paths.  
%	This permits the definition of a dual graph whose nodes are cycles in a random MCB and edges are their intersection paths. % break?
%	The theoretical backbone for this approach are the polyhedron-interchangeability ({\tt pi}) and short loop-interchangeability ({\tt sli}) classes to partition the relevant cycles. 
%	This theory advances the works~\cite{vismaraUnion1997,gleissInterchangeability2000,kolodzikUnique2012}.
%	The {\tt pi} classes identify building blocks for crystalline structures. The {\tt sli} classes group together sets of large redundant cycles. % break?
%	We investigate large amorphous hydrocarbon networks where we (i) theorize how the number of relevant cycles may explode with system size and 
%	(ii) observe small polyhedral structures that relate to diamond.
	
%	We investigate the problem of statistically sampling rings from complex chemical networks. 
%	The conventional approach of taking cycles from a minimum cycle basis (MCB)~\cite{kavithaCycle2009}, also called a smallest set of smallest rings~\cite{plotkinMathematical1971}, is not unique.
%	Unfortunately, an MCB of a graph is not in general unique and depends on node labels. 
%	Another common approach is to consider a larger set of cycles that is unique, e.g., the relevant cycles composed of the union of the MCBs of a graph.
%	These larger sets often are not polynomial with respect to graph size. 
%	We propose an alternate approach: to sample the MCBs of a graph uniformly at random. 
%	A random MCB is not unique, but it is statistically well-defined. 
%	While it is not unique, a random MCB is statistically well-defined.
%	We create theory for how pairs of MCB cycles can intersect, and we introduce a postprocessing step to force these pairwise intersections into paths. 
%	This permits the definition of a dual graph whose nodes are cycles in a random MCB and edges are their intersection paths. 
%	The theoretical backbone for this approach uses polyhedron-interchangeability (\tt pi) and short loop-
\end{abstract}

% remove with submission
\tableofcontents


%\add{Some notes:

\section{Introduction}

%{\color{red} use MCB, but repeat it a couple times\\
%{\tt pi}: polyhedral-interchangeability\\
%{\tt sli}: short loop interchangeability\\
%by 7pm}

%\add{
%Potential abbreviations:
%\begin{itemize}
%	\item Minimum cycle basis (MCB). This is standard in the literature. Maybe just full version in Figures?
%	\item Molecular Dynamics (MD). Also standard.
%	\item Interchangeability
%	\begin{itemize}
%		\item Letter I. Not my favorite because it is short and conflicting: I as in myself, $i$ the imaginary number, {\tt i} as an index, etc.
%		\item Three letters ICE. Easier to distinguish and direct- and trivial-interchangeability as DICE and TICE can be read as words. Con: the name ICE also matches a government agency.
%		\item Other? 
%	\end{itemize}
%%	 as \emph{i}. Not my favorite, very short has some ambiguity with I as in me, $i$ as in the imaginary number, {\tt i} as an index, etc.
%%	\item Interchangeability as \emph{ICE}. This has a more consistent format and it is easy to read. DICE and TICE also sound nice (accidental pun). Italics or verbatim may be better because this format matches the government agency.
%	\item V families for ``modified Vismara cycle families''.
%\end{itemize}
%}

\begin{figure*}
	\centering
	\includegraphics{Intro_Figs.pdf}
%	\includegraphics[height=1in]{Figures/Adamantane_Print.jpg}
%	Placeholder 1, Placeholder 2
	\caption{\textbf{(a)} Snapshot of the giant biconnected component in a simulation initialized to adamantane ($\rm C_{10}H_{16}$) at 4000K and 40.5GPa. 
	We highlight (i) \emph{essential}~\cite{gleissInterchangeability2000} pentagonal and hexagonal rings belonging to every MCB in blue, (ii) a {\tt sli} class composed of two 11-membered rings separated by a hexagonal ring in orange, and (iii) two {\tt pi} classes with the same shape as barallene ($\rm C_8H_8$) in green.
%	{\tt sli} classes containing multiple cycles of length $|C|\le 20$ are shown in orange. 
%	{\tt pi} classes containing multiple {\tt sli} classes are shown in green. 
	\textbf{(b)} A graph made of diamond-shaped cycles chained into a loop like a bracelet. An MCB of this graph includes all of the small diamond-shaped loops and one large loop. \textbf{(c)} Carbon skeleton of adamantane, the smallest diamondoid, i.e., diamond is formed by repeatedly tiling adamantane. An MCB of adamantane includes any 3 of the 4 hexagonal faces. \textbf{(d)} Depiction of the process of modifying a cycle $C_2\to C_2'$ such that the cycles $C_1$ and $C_2$ overlap over a single path.}
	\label{fig:intro_graphs}
\end{figure*}

\subsection{Motivation}


%Highly cyclic chemical networks appear in countless and diverse contexts ranging from amorphous materials such as polymers~\cite{torres-knoopModeling2018}, soot~\cite{villa-alemanRaman2022}, and silicate melts~\cite{clarkShockramp2023} to structured materials such as metals~\cite{???}, crystal~\cite{???}, and diamond~\cite{???}.
Amorphous material networks appear in numerous contexts such as polymers~\cite{torres-knoopModeling2018}, energetic carbon~\cite{dernovMapping2025,deringerMachine2017}, soot~\cite{villa-alemanRaman2022}, and silicate melts~\cite{clarkShockramp2023}. 
Understanding the network topology of these systems is critical because it influences their material properties.
% as it is tied to their physical structure and material properties.
The local atomistic structures within these networks can be leveraged to make predictions about their thermodynamic limits and to reduce combinatorial complexity.
%To better understand the thermodynamic limits of these systems, and to reduce combinatorial complexity, it is often desirable to characterize the local atomistic structures of macroscale materials. 
Here, we review, discuss, and extend theory and algorithms for the seemingly simple task of sampling rings from material networks. 
The assumptions under which these algorithms are efficient are consistent with the carbon skeletons of large hydrocarbons arising in hydrocarbon pyrolysis, 
%We will discuss these results in the context of hydrocarbon pyrolysis, 
a complex chemical reaction system composed of carbon and hydrogen at extreme temperatures and pressures.
%To guide our discussion, we consider hydrocarbon pyrolysis, a complex chemical reaction system composed of carbon and hydrogen at extreme temperatures and pressures.
%Understanding the network topology of these materials is important as it is tied to their physical structure and material properties.
%Of particular interest to the authors is the characterization of large molecules -- as in some fraction of the system size -- formed by the pyrolysis of hydrocarbons at extreme temperatures and pressures. 
%Initially, these large hydrocarbons can be described as tree-like structures with sparse, mostly disjoint loops, and their size distribution can be effectively predicted using random graphs~\cite{ruthCyclic2025}.
%As this system becomes more carbon-rich, the loops collect into clusters, e.g., Figure~\ref{fig:intro_graphs} (a). 
%In many cases, these clusters are amorphous, and robust definitions are needed to describe their structure.
%In this paper, we review methods for sampling loops in chemical graphs and provide a novel approach for addressing their limitations. 
%Of particular interest to the authors is the study of hydrocarbon pyrolysis at extreme temperatures and pressures. ... old ...

%{\color{red} 
%	1. MD is really valuable, but post-processing steps are needed. (1) They are expensive and extrapolating these results to other sets of initial conditions is desirable. (2) Analysis is needed to get a mechanistic understanding of these issues.
%	
%	Kinetic Monte Carlo involves modeling a chemical system using a set of reactions. These reaction and their rates can be learned from molecular dynamics and quantum simulations~\cite{}... Exponential growth, cite Cayley/Sloane...
%	
%	This is distracting from the message... MD alone is not mechanistic. 
%	
%For KMC, reactions need creative solutions not to be limited by combinatorial growth Cayley 1875}

Conventional approaches for studying hydrocarbon pyrolysis include molecular dynamics {(MD)}~\cite{chengThermodynamics2023,chenTransferable2019,dufour-decieuxTemperature2022,dufour-decieuxPredicting2023} simulations and kinetic Monte Carlo~\cite{dufour-decieuxAtomicLevel2021,dufour-decieuxTemperature2022}. 
MD is flexible and can be used to simulate various systems and thermodynamic conditions. 
However, MD simulations are expensive. They may take days or even weeks to achieve nanosecond timescales. Furthermore, they need to be performed anew for every set of initial conditions. 
%Molecular dynamics simulations are expensive and need to be performed once per set of initial conditions. 
Kinetic Monte Carlo~\cite{gillespieGeneral1976,highamModeling2008} simulates a chemical system using a set of known chemical reactions. 
These reactions and their reaction rates can be learned from a set of molecular dynamics simulations. This approach is limited by the combinatorial growth in the number of reactions, e.g., a simulation with $\sim10^3$ atoms results in $\sim10^4$ potential reactions~\cite{yangLearning2017,yangDataDriven2019,chenTransferable2019}.
%\add{Maybe, Cayley here for fun}
The combinatorial growth of hydrocarbons has been studied as early as 1875, where Cayley used generating functions to estimate the number of alkanes of a given length~\cite{cayley1875ueber,rains1999cayley}.
Due to this growth, only small molecules counts can be predicted accurately by this approach.
%Also, these works did not provide insights on the size of the largest molecule. 
To address this limitation, Refs.~\cite{dufour-decieuxAtomicLevel2021,dufour-decieuxTemperature2022} propose kinetic models with only local reactions accounting for chemical changes near the reaction site. 
%In the works~\cite{dufour-decieuxAtomicLevel2021,dufour-decieuxTemperature2022}, local reactions that only consider chemical changes near the reaction site were considered. 
This approach reduces the number of reactions and can be used to make predictions about the size of the largest molecule. 
However, like MD simulations, the local kinetic Monte Carlo requires a new run for each set of initial conditions. 
%Still, they require the simulation of full trajectories using kinetic Monte Carlo.
%The results from these simulations can be extrapolated to other thermodynamic conditions using kinetic Monte Carlo with a set of learned reactions and reaction rates. 
%Unfortunately, the presence of large molecules
%Kinetic Monte Carlo is limited by combinatorial complexity, requiring accurate estimates for upwards of xxx molecular reactions~\cite{vincent group}. 
%Furthermore, kinetic Monte Carlo methods poorly reproduce the size of the largest molecule. These issues can be improved by considering reactions that are localized to the reaction site~\cite{vincent paper}

Our research group utilized random graphs as equilibrium models for hydrocarbon pyrolysis~\cite{dufour-decieuxPredicting2023,ruthCyclic2025}. 
%In recent works~\cite{dufour-decieuxPredicting2023,ruthCyclic2025}, we used random graphs as equilibrium models for hydrocarbon pyrolysis.
For hydrogen-rich systems, the configuration model~\cite{newmanRandom2001,molloySize1998} -- a classical random graph model -- was used to estimate the molecule size distribution~\cite{dufour-decieuxPredicting2023}. 
This model is particularly effective for small alkanes ($\le20$ carbon atoms) or fully saturated acyclic hydrocarbons.
However, as the hydrogen-to-carbon ratio decreases, the configuration model significantly exaggerates the size of the largest molecule, mostly because it does not promote the formation of small carbon rings abundant in hydrocarbon molecules.
%mostly because of small carbon rings. 
Using Refs.~\cite{karrerRandom2010,newmanAssortative2002} as building blocks, we introduced a random graph model with disjoint loops that accurately predicted the size of the largest molecule~\cite{ruthCyclic2025}. This disjoint loop model also predicted the phase transition between the states with and without a giant molecule, i.e., a molecule that contains a macroscopic fraction of the total number of atoms. 
%Following this we were able to predict the size distribution of the largest molecule and the phase transition of the giant molecule -- a macroscopic
%This approach is effective for hydrogen-rich systems where all molecules are small ($\le 20$ carbon atoms) and form trees. 
%However, in the presence of a giant component, i.e., a macroscopic molecule that contains some fraction of the total number of atoms, the configuration model overestimates the size of the largest molecule by a factor of 2. Initially, this giant component consists of small loops joined by tree-like structures. The size of this giant component is accurately predicted by a random graphs model with disjoint loops~\cite{ruthCyclic2025}. 
In carbon-rich systems, the rings cluster together, and the disjoint loop assumption fails. See Figure~\ref{fig:intro_graphs}~(a). In this setting, proper statistics for carbon rings are required.

%\begin{enumerate}
%	\color{red}
%	\item Of particular interest to me..
%	\item This can be studied using molecular dynamics (cheng,vincent,the really large lammps md simulation) and through kinetic monte carlo.
%	\item KMC works well when molecules are small (hydrogen-rich)
%	\item Local reactions help (Vincent). Random graphs reduce number of parameters and help to describe how emergent molecules appear as a result of local structures (2023 paper)
%	\item In intermediate phases (up to H/C ratio 2.25) can be modeled with disjoint loops. Can capture the formation of a large molecule (our 2025 paper).
%	\item For carbon rich systems, a loop cluster forms. Here disjoint loops are not true, and do not capture the important structures in the system.
%	\item Common approaches for sampling loops have flaws that make it hard to study this system. We review and modify these definitions so we can effectively coarse grain the loop cluster into loops instead of atoms.
%\end{enumerate}

\subsection{An overview}

At a macroscopic scale, a ring cluster can be identified as a biconnected component or a collection of nodes where each pair of nodes is connected by two or more disjoint paths~\cite{diestelGraph2017,bollobasModern1998}. A method for computing biconnected components is given in the 1971 paper by Hopcroft and Tarjan~\cite{hopcroftAlgorithm1973}. At a microscopic scale, we want to obtain the building blocks of these clusters, i.e., their loops. In many cases, the na\"ive approach of computing all cycles is computationally intractable. The total number of cycles grows exponentially with graph size, many of which are redundant. Instead, most methods for sampling loops utilize the cycle space or the vector space of cycles.

Classically, the loops of a graph are sampled from a \emph{minimum cycle basis} (MCB) or a basis of the cycle space that minimizes the sum of cycle lengths~\cite{kavithaCycle2009}. In chemical systems, a minimum cycle basis is also called a smallest set of smallest rings, as introduced by Plotkin~\cite{plotkinMathematical1971}. {MCBs} are not without their flaws. Most notably, they are not unique, so the loops sampled via an MCB depend on node labels and algorithmic implementation. This issue has led to several competing definitions for sampling cycles in chemistry literature~\cite{bergerCounterexamples2004}.

Here, we focus on the approach introduced by Vismara, 1997~\cite{vismaraUnion1997}, where the set of \emph{relevant cycles} defined as the union of all MCBs is considered. 
Equivalently, the relevant cycles are the unique cycles that cannot be decomposed into smaller cycles. 
%Furthermore, Vismara outlines an approach for grouping the relevant cycles into families. 
In the worst case, the number of relevant cycles grows exponentially with respect to the number of nodes, but in practice the number of relevant cycles is typically much smaller~\cite{mayEfficient2014}. 
Nevertheless, this is an undesirable property of the relevant cycles. 
To guarantee properties such as the number of relevant cycles and their lengths can be computed in polynomial time, Vismara partitions the relevant cycles into families of related cycles~\cite{vismaraUnion1997}.
%To guarantee properties of the relevant cycles (e.g., the number of relevant cycles and cycle lengths) can be computed in polynomial time, Vismara partitions the relevant cycles into families of related cycles~\cite{vismaraUnion1997}.
%To handle this growth, Vismara partitions the relevant cycles into families of related cycles~\cite{vismaraUnion1997}. 
%In practice, we find that Vismara's cycle families are more useful as they group
%In the worst case, the number of relevant cycles can grow exponentially with the number of nodes, but in practice the number of relevant cycles grows cubically with the number of nodes~\cite{check}. Regardless, many of the relevant cycles are redundant, and often the relevant cycle families are more useful. 
%Unfortunately, Vismara's cycle families cannot are not unique.
%We review these definitions in Section~\ref{sub:MCB+Vis_Fam}.


Another approach would be to investigate Vismara's cycle families directly. However,  these families are not unique, thus returning us to our original issue. 
To resolve this issue, Gleiss~{et al.}, 2000,~\cite{gleissInterchangeability2000} introduce the \emph{interchangeability relation}, an equivalence relation that partitions the relevant cycles into unique interchangeability classes. 
This relation helps to characterize the lack of uniqueness in MCBs.
Two relevant cycles are interchangeable if one cycle can be swapped for the other in a cycle basis, and if the two cycles only differ by equal or shorter length cycles. 
Often, this relation groups together large cycles that differ by smaller cycles. 
For example, the cycles around the perimeter of the bracelet graph in Figure~\ref{fig:intro_graphs} (b) form a single interchangeability class. 
In other cases, this relation groups together cycles that belong to three-dimensional structures. 
For instance, the four hexagonal faces of adamantane in Figure~\ref{fig:intro_graphs} (c) form a single interchangeability class.


Kolodzik {et al.}, 2012~\cite{kolodzikUnique2012}, argue these interchangeability classes are too coarse and introduce the \emph{unique ring family (URF)-pair-relation}. This relation is more restrictive, where two relevant cycles are URF-pair-related if they share edges and if they differ by \emph{strictly shorter} cycles. This separates the cases above. The large cycles in the bracelet graph are URF-pair-related because they differ by smaller diamond cycles. The hexagonal faces in adamantane are not URF-pair-related because they differ by equal length hexagonal faces. The URF-pair-relation is not an equivalence relation. Instead, the unique ring families are given by the transitive closure of the URF-pair-relation.

\subsection{The goal and a summary of main results}

%{\color{red}Following paragraph mostly fixed. edit for readability and continuity}
{The goal of this work} is to review and extend theory and algorithms that describe and extract complex amorphous ring clusters in carbon-rich simulation data of hydrocarbon pyrolysis. To achieve this, we adapt and advance concepts and methods from Refs.~\cite{gleissInterchangeability2000,kolodzikUnique2012}. 
Our theoretical developments are the following:
\begin{itemize}
	\item 
	At a course level, we introduce the \emph{polyhedron-interchangeability} ({\tt pi}) relation, a modification of the interchangeability relation from Ref.~\cite{gleissInterchangeability2000}.
	Unlike interchangeability, {\tt pi} relates pairs of cycles that can be swapped between MCBs. We prove {\tt pi} is an equivalence relation (Theorem~\ref{Thm:DirInt_Eq_Rel}) and the resulting equivalence classes can be computed using a given MCB (Theorem~\ref{Thm:DirInt_Comp}). The theory of {\tt pi} also characterizes how we can modify MCBs by swapping cycles.
%	At a coarse level, we introduce \emph{direct-interchangeability}, a modification of the interchangeability relation from Ref.~\cite{gleissInterchangeability2000}. This definition is defined in terms of swapping cycles between minimum cycle bases. Thus, direct-interchangeability gives a more precise characterization of lack of uniqueness in minimum cycle bases.\add{Advertise key theorems. Make contributions more clear.}
	\item At a granular level, we introduce the \emph{short loop-interchangeability} ({\tt sli}) relation, adapted from the URF-pair-relation from Ref.~\cite{kolodzikUnique2012}. {\tt sli} removes the requirement that cycles share edges from the URF-pair-relation. As a result, {\tt sli} is an equivalence relation. 
	{\tt sli} was first introduced in the dissertation of Gleiss~\cite{gleissShort2001} as stronger interchangeability, where it was shown to be an equivalence relation. 
	We prove the {\tt sli} classes can be computed using a given MCB. 
	We suggest to use the {\tt pi} and {\tt sli} relations together to analyze ring clusters.
	%Furthermore, this trivial-interchangeability relation is an equivalence relation that naturally partitions the relevant cycles. This matches more closely matches the remaining theory in this work.\add{What was proven and what we prove.}
\end{itemize}


%In this work, we consider a two-layer partition of the relevant cycles. At a high level, we introduce the direct-interchangeability relation. This definition is defined by swapping cycles between minimum cycle bases. Thus, direct-interchangeability more accurately characterizes lack of uniqueness in minimum cycle bases. At a more granular level, we introduce the trivial-interchangeability relation inspired by the URF pair relation. In particular, trivial-interchangeability removes the requirement that pairs of cycles need to share nodes, which makes it into a equivalence relation.

%To resolve this issue Gleiss {et al.}~\cite{gleissInterchangeability2000} introduce the interchangeability relation, where a pair of relevant cycles are interchangeable if one can be swapped for the other in a cycle basis. This interchangeability relation gives insights about lack of uniqueness in minimum cycle bases. Furthermore, interchangeability is an equivalence relation, so it naturally partitions the relevant cycles into equivalence classes. These interchangeability classes are coarse, e.g., the hexagonal faces of adamantane in Figure~\ref{fig:intro_graphs} (x) form a single interchangeability class. Kolodzik {et al.}

%To the uninitiated, this approach is likely unappealing. Families of cycles are abstract and they lack obvious structure. Instead, we argue this approach is valuable as a theoretical tool for understanding the nature of the relevant cycles. To justify the practical value of this approach, we consider three consequences of our theory.
%Unfortunately, it is hard to apply this theory directly to our problem of hydrocarbon pyrolysis. Families of cycles are abstract and they lack uniform structure. This makes them hard to use for modeling loop clusters. Instead, we highlight three novel consequences of this theory that enhance upon the classical approach of sampling minimum cycle bases. 
%{\color{red} Positive...}
Using this theory we develop the following methods for sampling cycles in a graph.

%This approach is continued by Gleiss {et al.}, where pairs of relevant cycles are interchangeable if they can be swapped using the basis exchange property of cycle bases. We provide background for the basis exchange property in Section~\ref{sub:basis_ex}. Interchangeability is an equivalence relation so it naturally partitions the relevant cycles into equivalence classes. These interchangeability classes are unique and help to characterize lack of uniqueness in minimum cycle bases. Kolodzik {et al.}~\cite{kolodzikUnique2012} critique the interchangeability classes, arguing they are too coarse. They highlight this issue with the cube graph, where all of the square faces belong to the same interchangeability class. To address this, they introduce the unique ring families a more granular partition of the relevant cycles.
%
%Our approach is as follows. We consider a two-layer partition of the relevant cycles. At a course level, we define direct-interchangeability classes, a modification of the interchangeability classes from Ref.~\cite{gleissInterchangeability2000}. We make this modification because the direct-interchangeability relation more closely aligns to exchanging cycles between minimum cycle bases, see Definition~\ref{def:dir_intchg}, and the direct-interchangeability classes are easy to compute given Vismara's cycle families, see Section~\ref{sub:DICE_TICE_Comp}. At a granular level, we define trivial-interchangeability classes a modification of the unique ring families from Ref.~\cite{kolodzikUnique2012} that is defined using an equivalence relation.



%{\color{red} Our approach is a two layer partition... At a coarse level, we introduce direct-interchangeability classes, a refinement of the interchangeability classes from Ref.~\cite{gleissInterchangeability2000} that characterizes the lack of uniqueness in minimum cycle bases... Something about why we improve Kolodzik {et al.}}

%To justify the complex definitions introduced in this manuscript, we introduce three theoretical applications to highlight the utility of this approach:
\begin{itemize}
	\item \emph{Random sampling of minimum cycle bases.} We introduce a method for sampling MCBs uniformly at random. %The algorithm for this is outlined in Section~\ref{sub:rand_MCB}. 
	Of course, a random MCB is not deterministic, but it does have a concrete definition as a probabilistic object. 
	Thus, the practitioner can confidently sample loops from a random MCB knowing it is statistically well-defined.
%	This allows the practitioner to sample loops from a minimum cycle basis with the added confidence that these loops are at least statistically well-defined. See Section~\ref{sub:rand_MCB}.
	\item \emph{Pairwise intersection of cycles.} 
	A graph is defined by nodes and by edges denoting connections between nodes. 
	Similarly, we are interested in loops and their intersections that resulting in the formation of loop cluster. 
%	When studying a graph, we need to understand both the nodes and the edges that connect them. 
%	Similarly, we are not simply interested in sampling loops; we are also interested in how these loops overlap and how they coalesce into clusters. 
	In this work, we fully characterize the way pair of cycles in MCBs can intersect. 
	Let the intersection of two cycles contain more than one path as in Fig.~\ref{fig:intro_graphs}~(d). 
	We introduce a method for modifying one of these cycles to make the intersection consist of a single path. 
%	We also introduce a method for modifying a cycle $C_2\to C_2'$ such that the pair of cycles $C_1,C_2'$ intersect over a single path. See Figure~\ref{fig:intro_graphs}~(d). 
%	{\color{red} In Section~\ref{sub:pair_intersect}, we fully characterize the format for how pairs of cycles in a minimum cycle basis may intersect. 
%	This informs how these pairwise loop intersections should be stored in a computer.
%	Furthermore, in Section~\ref{sub:pair_intersect}, we use our theory to modify a minimum cycle basis $\calM\to\calM'$ to obtain a new minimum cycle basis where each pair of overlapping cycles intersect over a single path\textbf{***this is conjecture, missing an inductive argument***}. }
	This property is desirable as individual paths are conceptually simple and they are easy to represent computationally as node lists.
	
	Using this, we define a dual graph, which can be interpreted as a graph of cycles. The notion of dual graph has a long history in planar graphs~\cite{diestelGraph2017}, and graphs of cycles have been considered for modeling chemical systems as well~\cite{torres-knoopModeling2018,nouleho2019improving}. Our contribution is defining the structure of these dual graphs for MCBs.
\end{itemize}

We apply this theory to the hydrocarbon pyrolysis system. In most cases, we find that lack of uniqueness of MCBs is largely caused by rings at varying lengths as in Figure~\ref{fig:intro_graphs}~(b). We find ring statistics vary greatly between the relevant cycles and an MCB. In particular, the relevant cycles introduce redundant large rings. 
Furthermore, our analysis of the relevant cycles naturally uncovers three-dimensional structures. 
These structures take the form of polyhedra with non-flat faces as in Figure~\ref{fig:intro_graphs} (c). 
These polyhedra are useful for describing crystalline materials. 
Proper sampling of these polyhedra will require further theoretical advancements.
%Also, using our theory we observe some polyhedral structures like adamantane in Figure~\ref{fig:intro_graphs}~(c). 

%\begin{enumerate}
%	\color{red}
%	\item Later we apply this to hydrocarbon pyrolysis
%	\item Mostly observe non-uniqueness as a consequence of varying cycle lengths (fig. 1b).
%	\item Sometimes find 3d structures, e.g., adamantane. These are important for the study of dense carbon materials in the formation of diamond.
%\end{enumerate}

%Ultimately, the usefulness of our approach will be determined by how easily and effectively it characterizes material networks. In Section~\ref{sec:results}, we highlight how this theory decomposes large, amorphous loop clusters obtained by ReaxFF~\cite{vanduinReaxFF2001,chenowethReaxFF2008,srinivasanDevelopment2015,ashrafExtension2017} simulations of hydrocarbon pyrolysis. In Section~\ref{sec:algorithm}, we introduce algorithms for computing the theoretical structures discussed in this manuscript. In the GitHub repository \url{https://github.com/perrineruth/enhanced-minimum-cycle-bases}, we provide a Python implementation of these codes which take a simple graph from the NetworkX package~\cite{hagbergExploring2008} as input. 

This manuscript is organized as follows. In Section~\ref{sec:background}, we provide the necessary background for the cycle space including the relevant cycles from Ref.~\cite{vismaraUnion1997} and the partitions of the relevant cycles from Refs.~\cite{gleissInterchangeability2000,kolodzikUnique2012}. In Section~\ref{sec:Theory}, we introduce the {\tt pi} and {\tt sli} relations used to partition the relevant cycles. In Section~\ref{sec:algorithm}, we provide algorithms for computing the {\tt pi} and {\tt sli} classes. Our Python codes are available on Github~\cite{RuthGithub}. These codes take a simple graph from the NetworkX package~\cite{hagbergExploring2008} as input.
% \url{https://github.com/perrineruth/enhanced-minimum-cycle-bases} which takes a simple graph from NetworkX package~\cite{hagbergExploring2008} as input. 
In Section~\ref{sec:sampling}, we provide an algorithm for sampling MCBs uniformly at random, and we characterize the intersection of cycles in MCBs. In Section~\ref{sec:results}, we show how this approach can be used to characterize ring clusters from hydrocarbon pyrolysis simulations. 
These molecular dynamics simulations are obtained via LAMMPS~\cite{plimptonFast1995,thompsonLAMMPS2022} using the ReaxFF~\cite{vanduinReaxFF2001,chenowethReaxFF2008,srinivasanDevelopment2015,ashrafExtension2017} force field permitting chemical reactions.

%Regardless of context, the importance of characterizing the loop clusters that form these clusters is clear. 
%
%At a macroscopic scale, a loop cluster can be identified as a biconnected component, which can be computed using an algorithm developed by Hopcroft and Tarjan in 1971(). At a microscopic scale, we want to characterize the smallest components of these clusters, i.e., the loops. Unfortunately, sampling loops is not straightforward

%\begin{itemize}
%	\item Study hydrocarbon pyrolysis
%	\item Cannot use disjoint loops
%	\item describing cycle statistics of biconnected components
%	\item People find this issue often in chemical applications
%	\item Sort out concepts, useful, but not completely sufficient, and upgrade 
%	
%	\item Table of concepts: brief definitions
%	Columns: name, definition, pros cons Add with figure 9 to late Section 3
%\end{itemize}
%
%Some additional notes not yet incorporated into the paper:
%\begin{itemize}
%	\item Polyhedra motivation
%	\begin{itemize}
%		\item Paper by Yokoyama et al. ``Crystal Structure Generation Based on Polyhedra Using Dual Periodic Graphs''. Also computes polyhedra, but requires a periodic structure, ours does not. This work may help to motivate the importance of polyhedra in crystals.
%		\item Book ``Shaping Space
%		Exploring Polyhedra in Nature, Art, and the Geometrical Imagination''. Has three chapters on polyhedra in chemical spaces. Additional motivation.
%		\item Hamon and Taylor ``A synthesis of tetracyclo[5,3,1,12,6,04,9]dodecane (iceane)''. They mention iceane as a unit of the londsdaleite crystalline structure. Mention historically bicyclo[2,2,2]octane (Fig.~\ref{fig:MCB_Unique}) by our definition should share this role with iceane. \textit{Mention in discussion section perhaps?}
%	\end{itemize}
%\end{itemize}

%Cycle bases and minimum cycle bases of graphs have applications in many fields, see e.g.~\cite{kavithaCycle2009} and references therein. Of particular interest here, is minimum cycle bases as a tool for chemical ring perception, see e.g.~\cite{bergerCounterexamples2004}. Often, minimum cycle bases are avoided in chemical ring perception because they are computationally expensive and not necessarily unique. For instance, the OpenEye(trademark or something) package explicitly does not implement a minimum cycle basis -- or SSSR (smallest set of smallest rings) -- algorithm for this reason~\cite{websiteforOpeneye}. However, minimum cycle bases are still informative justifying their implementation in many other packages, e.g. in the Chemistry Development Toolkit~\cite{mayEfficient2014}. This conflict between the uniqueness of minimum cycle bases and their application is further discussed in the work by Bergers \textit{et al.}~\cite{bergerCounterexamples2004}.\textbf{*** so far more of a thesis, split here ***} Here, we modify the definition of Unique Ring Families by Kolodzik \textit{et al.}~\cite{kolodzikUnique2012} which gives a perspective of the cycle space that is more granular than a minimum cycle basis, and the definition of exchangeable cycles by Gleiss \textit{et al.}~\cite{gleissInterchangeability2000} to obtain a perspective of the cycle space that is less granular than a minimum cycle basis. 



\section{Background}
\label{sec:background}

\subsection{Cycle space and fundamental cycle bases}

\label{sub:cyc_space}


%We start with a brief discussion of cycle spaces and a method for computing their bases. 
Cycle spaces have a rich history in electronic circuit theory, chemistry, and other contexts. We start with a brief discussion of classical results in this field. We also introduce notation, which is merged from Refs.~\cite{gleissInterchangeability2000,kavithaCycle2009,vismaraUnion1997}. For a more detailed history and discussion of these topics we recommend the books on graph theory~\cite{bollobasModern1998,diestelGraph2017}, review of cycle bases~\cite{kavithaCycle2009}, review of cycles in chemical systems~\cite{bergerCounterexamples2004}, and references therein. 

\begin{figure}[t]
	\centering
	\includegraphics{CycSp.pdf}
	%	\begin{tikzpicture}[scale=.86*.707]
		%		\node () at (-.5,4) {\textbf{(a)}};
		%		{\footnotesize
			%		\node[namednode] (1) at (0,1) {1};
			%		\node[namednode] (2) at (1,2) {2};
			%		\node[namednode] (3) at (2,3) {3};
			%		\node[namednode] (4) at (3,2) {4};
			%		\node[namednode] (5) at (4,3) {5};
			%		\node[namednode] (6) at (5,2) {6};
			%		\node[namednode] (7) at (4,1) {7};
			%		\node[namednode] (8) at (3,0) {8};
			%		\node[namednode] (9) at (2,1) {9};
			%		\node[namednode](10) at (1,0){10};
			%		}
		%		\draw[very thick] (1)--(2)--(3)--(4)--(5)--(6)--(7)--(8)--(9)--(10)--(1);
		%		\draw[very thick] (2)--(9)--(4)--(7);
		%		\node () at (1,1) {$B_1$};
		%		\node () at (2,2) {$B_2$};
		%		\node () at (3,1) {$B_3$};
		%		\node () at (4,2) {$B_4$};
		%		
		%		\node () at (6.5,4) {\textbf{(b)}};
		%		
		%		{\footnotesize
			%		\node[namednode,opacity=.2] (1a) at (7, 1) {1};
			%		\node[namednode] (2a) at (8, 2) {2};
			%		\node[namednode] (3a) at (9, 3) {3};
			%		\node[namednode] (4a) at (10,2) {4};
			%		\node[namednode,opacity=.2] (5a) at (11,3) {5};
			%		\node[namednode,opacity=.2] (6a) at (12,2) {6};
			%		\node[namednode] (7a) at (11,1) {7};
			%		\node[namednode] (8a) at (10,0) {8};
			%		\node[namednode] (9a) at (9, 1) {9};
			%		\node[namednode,opacity=.2](10a) at (8, 0) {10};
			%		}
		%		\draw[very thick] (2a)--(3a)--(4a)--(7a)--(8a)--(9a)--(2a);
		%		\draw[thick,opacity=.2] (2a)--(1a)--(10a)--(9a)(4a)--(5a)--(6a)--(7a);
		%		\draw[thick,opacity=.1](4a)--(9a);
		%		\node() at (9.5,4) {$C_1=B_2\oplus B_3$};
		%		
		%		{\footnotesize
			%		\node[namednode] (1b) at (13,1) {1};
			%		\node[namednode] (2b) at (14,2) {2};
			%		\node[namednode,opacity=.2] (3b) at (15,3) {3};
			%		\node[namednode] (4b) at (16,2) {4};
			%		\node[namednode,opacity=.2] (5b) at (17,3) {5};
			%		\node[namednode,opacity=.2] (6b) at (18,2) {6};
			%		\node[namednode] (7b) at (17,1) {7};
			%		\node[namednode] (8b) at (16,0) {8};
			%		\node[namednode] (9b) at (15,1) {9};
			%		\node[namednode](10b) at (14,0) {10};
			%		}
		%		\draw[very thick] (1b)--(2b)--(9b)--(4b)--(7b)--(8b)--(9b)--(10b)--(1b);
		%		\draw[thick,opacity=.2] (2b)--(3b)--(4b)--(5b)--(6b)--(7b);
		%		\node() at (15.5,4) {$C_2=B_1\oplus B_3$};
		%		
		%		{\footnotesize
			%		\node[namednode] (1c) at (19,1) {1};
			%		\node[namednode] (2c) at (20,2) {2};
			%		\node[namednode,opacity=.2] (3c) at (21,3) {3};
			%		\node[namednode] (4c) at (22,2) {4};
			%		\node[namednode] (5c) at (23,3) {5};
			%		\node[namednode] (6c) at (24,2) {6};
			%		\node[namednode] (7c) at (23,1) {7};
			%		\node[namednode,opacity=.2] (8c) at (22,0) {8};
			%		\node[namednode] (9c) at (21,1) {9};
			%		\node[namednode](10c) at (20,0) {10};
			%		}
		%		\draw[very thick] (1c)--(2c)--(9c)--(10c)--(1c);
		%		\draw[very thick] (4c)--(5c)--(6c)--(7c)--(4c);
		%		\draw[thick,opacity=.2] (2c)--(3c)--(4c)--(9c)--(8c)--(7c);
		%		\node() at (21.5,4) {$C_3=B_1\oplus B_4$};
		%	\end{tikzpicture}
	\caption{A graph with four face cycles $B_1$, $B_2$, $B_3$, and $B_4$ forming a basis. Three additional cycles $C_1$, $C_2$, and $C_3$ are shown. $C_1$ is a single loop while $C_2$ and $C_3$ are not.}
	% remove node labels and make single column?
	\label{fig:cyc_space} 
\end{figure}

Let $G=(V,E)$ be a graph with $n\coloneq|V|$ nodes and $m\coloneq|E|$ edges. Formally, a \textit{cycle} is an edge-induced subgraph of $G$ where all nodes have even degree. More intuitively, a cycle is a collection of loops that do not share edges. See Fig.~\ref{fig:cyc_space}. We note, nearly all cycles we consider are simple loops, including all cycles introduced later in this section in fundamental and minimum cycle bases.
%Sometimes, single loop subgraphs are referred to as ``circuits'', but this notation is not universally
%A \textit{circuit}, is a single loop subgraph. Therefore, a circuit is a specific type of cycle, and a cycle is the union of circuits. The term ``cycle'' is used more universally than ``circuit'', and in practice most cycles that are important to us are circuits -- including the set of relevant cycles which we review in the next subsection. 
%Thus, when we refer to a cycle we are typically referring to a circuit.
%Outside of this subsection, we will avoid the term circuit
%The term circuit is less universal. 
%So, outside of this subsection we will avoid this term, and we note when we refer to a cycle we are typically referring to a circuit.

The utility of this definition is that the set of cycles -- denoted by $\CycSp$ -- forms a vector space called the \textit{cycle space}. %, denoted by $\CycSp$. 
The operation we apply to cycles is the symmetric difference $C_1\oplus C_2=(C_1\setminus C_2)\cup(C_2\setminus C_1)$. % -- where the cycle space is the closure of circuits under $\oplus$. 
The scalars of this vector space, 
0 and 1, comprise the field $GF(2)$ with addition and multiplication modulo 2.
% is the field $GF(2)$ consisting of binary values $\{0,1\}$ where addition {and multiplication} are computed modulo 2. 
Note, a cycle is necessarily its own additive inverse, $C\oplus C = \emptyset$.  
%Importantly, the cycle space is the closure of circuits under the symmetric difference.

Since a cycle $C\in\CycSp$ is a set of edges, it may be represented as an \textit{incidence vector}: a vector in $\{0,1\}^m$ where element $i$ indicates whether edge $i$ is contained in $C$. 
{Not all incidence vectors correspond to cycles.} 
In vector form, the symmetric difference of cycles can be computed via an element-wise XOR operation (or by vector addition modulo 2). We distinguish that the cycle space is only a subspace of the set of incidence vectors, and ultimately we want a lower-dimensional representation of cycles.

\begin{figure*}[t]
	\centering
	\includegraphics[width=6.5in]{FCB_Example.pdf}
%	\begin{tikzpicture}[scale=.95]
%		\node[] () at (0,2.75) {\textbf{(a)}};
%		
%		% nodes
%		{\footnotesize
%			\node[namednode] (1) at (1.866, 2)   {1};
%			\node[namednode] (2) at (1,     1.5) {2};
%			\node[namednode] (3) at (1,     0.5) {3};
%			\node[namednode] (4) at (1.866, 0)   {4};
%			\node[namednode] (5) at (2.732, 0.5) {5};
%			\node[namednode] (6) at (2.732, 1.5) {6};
%			\node[namednode] (7) at (3.598, 2)   {7};
%			\node[namednode] (8) at (4.464, 1.5) {8};
%			\node[namednode] (9) at (4.464, 0.5) {9};
%			\node[namednode](10) at (3.598, 0)   {10}; 
%		}
%		\node () at (1.866,1) {$B_1$};
%		\node () at (3.598,1) {$B_2$};
%		% edges
%		\draw[very thick, blue, <-] (1)edge(2)edge(6)
%		(2)edge(3)
%		(3)edge(4)
%		(6)edge(5)edge(7)
%		(7)edge(8)
%		(5)edge(10)(10)edge(9);
%		\draw[very thick, red] 	(5) -- node[auto,midway] (){$e_1$} (4)
%		(8) -- node[auto,midway] (){$e_2$} (9);
%		% labels
%		\node[anchor=north] () at (2.732,-.27) {$\begin{aligned}
%				{e_1}\to B_1&=(1,2,3,4,5,6)\\
%				{e_2}\to B_2&=(5,6,7,8,9,10)\\
%				\calF&=\{B_1,B_2\}
%			\end{aligned}$};%$\calF=\{C_1,C_2\}$};
%	
%	\node[] at (5.5,2.75) {\textbf{(b)}};
%	\node[anchor=west] at (6,2.75) {Cycle};
%	\node[anchor=west] at (6,1) {Visual};
%	\node[anchor=west] at (6,-.75) {$\calF$-vector};
%	{\footnotesize
%		\node[namednode] (1a) at (8.866, 2)   {1};
%		\node[namednode] (2a) at (8,     1.5) {2};
%		\node[namednode] (3a) at (8,     0.5) {3};
%		\node[namednode] (4a) at (8.866, 0)   {4};
%		\node[namednode] (5a) at (9.732, 0.5) {5};
%		\node[namednode] (6a) at (9.732, 1.5) {6};
%		\node[namednode,opacity=.2] (7a) at (10.598,2)   {7};
%		\node[namednode,opacity=.2] (8a) at (11.464,1.5) {8};
%		\node[namednode,opacity=.2] (9a) at (11.464,0.5) {9};
%		\node[namednode,opacity=.2](10a) at (10.598,0)   {10};
%	}
%	\draw[very thick] (4a)--(3a)--(2a)--(1a)--(6a)--(5a);
%	\draw[very thick, red] (4a)--(5a);
%	\draw[thick, opacity=.2] (6a)--(7a)--(8a)(9a)--(10a)--(5a);
%	\draw[thick, red, opacity=.2] (8a)--(9a);
%	\node[] () at (9.732,2.75) {$B_1$};
%	\node[] () at (9.732,-.75) {$(1,0)$};
%	\node[] (e1) at (8.75,-1.75) {$e_1\in C$};
%	\draw[thick,->] (e1) -- (9.5,-1);
%	\node[] (e2) at (10.75,-1.75) {$e_2\in C$};
%	\draw[thick,->] (e2) -- (10,-1);
%	
%	{\footnotesize
%		\node[namednode] (1b) at (8.866+4.5, 2)   {1};
%		\node[namednode] (2b) at (8+4.5,     1.5) {2};
%		\node[namednode] (3b) at (8+4.5,     0.5) {3};
%		\node[namednode] (4b) at (8.866+4.5, 0)   {4};
%		\node[namednode] (5b) at (9.732+4.5, 0.5) {5};
%		\node[namednode] (6b) at (9.732+4.5, 1.5) {6};
%		\node[namednode] (7b) at (10.598+4.5,2)   {7};
%		\node[namednode] (8b) at (11.464+4.5,1.5) {8};
%		\node[namednode] (9b) at (11.464+4.5,0.5) {9};
%		\node[namednode](10b) at (10.598+4.5,0)   {10};
%	}
%	\draw[very thick]   (4b)--(3b)--(2b)--(1b)--(6b)--(7b)--(8b)
%	(9b)--(10b)--(5b);
%	\draw[thick, opacity=.15] (5b)--(6b);
%	\draw[very thick, red] (4b)--(5b)(8b)--(9b);
%	\node[] () at (9.732+4.5,2.75) {$B_1\oplus B_2$};
%	\node[] () at (9.732+4.5,-.75) {$(1,1)$};
%	
%	\def\dy{5.5}
%	\node () at (0,2.5-\dy) {\textbf{(c)}};	
%	{\footnotesize
%		\node[namednode] (1c) at (1.866, 2-\dy)   {1};
%		\node[namednode] (2c) at (1,     1.5-\dy) {2};
%		\node[namednode] (3c) at (1,     0.5-\dy) {3};
%		\node[namednode] (4c) at (1.866, 0-\dy)   {4};
%		\node[namednode] (5c) at (2.732, 0.5-\dy) {5};
%		\node[namednode] (6c) at (2.732, 1.5-\dy) {6};
%		\node[namednode] (7c) at (3.598, 2-\dy)   {7};
%		\node[namednode] (8c) at (4.464, 1.5-\dy) {8};
%		\node[namednode] (9c) at (4.464, 0.5-\dy) {9};
%		\node[namednode](10c) at (3.598, 0-\dy)   {10}; 
%	}
%	\draw[very thick, blue, <-] (2c)edge(3c)(3c)edge(4c)(4c)edge(5c)(5c)edge(10c)(10c)edge(9c)
%	(2c)edge(1c)(1c)edge(6c)(6c)edge(7c)(7c)edge(8c);
%	\draw[very thick, red] (6c)--node[auto, midway]{$e_1'$}(5c)
%	(8c)--node[auto, midway](e2p){$e_2'$}(9c);
%	
%	{\footnotesize
%		\node[namednode] (1e) at (8.866, 2-\dy)   {1};
%		\node[namednode] (2e) at (8,     1.5-\dy) {2};
%		\node[namednode] (3e) at (8,     0.5-\dy) {3};
%		\node[namednode] (4e) at (8.866, 0-\dy)   {4};
%		\node[namednode] (5e) at (9.732, 0.5-\dy) {5};
%		\node[namednode] (6e) at (9.732, 1.5-\dy) {6};
%		\node[namednode] (7e) at (10.598,2-\dy)   {7};
%		\node[namednode] (8e) at (11.464,1.5-\dy) {8};
%		\node[namednode] (9e) at (11.464,0.5-\dy) {9};
%		\node[namednode](10e) at (10.598,0-\dy)   {10};
%	}
%	\node (B2p) at (9.732,2.75-\dy) {$B_2'=B_1\oplus B_2$};
%	\draw[thick,->] (e2p) edge[in=180] (B2p);
%	\draw[very thick] (8e)--(7e)--(6e)--(1e)--(2e)--(3e)--(4e)--(5e)--(10e)--(9e);
%	\draw[thick,red,opacity=0.2] (5e)--(6e);
%	\draw[very thick, red] (8e)--(9e);
%	
%	%		\node[] () at (9.732+4.5,1-\dy) {$\begin{aligned}
%			%				\calF' &= \{B_1',B_2'\}\\ &= \{B_1,B_1\oplus B_2\}
%			%		\end{aligned}$};
%\end{tikzpicture}
%	\includegraphics[width=.75\columnwidth]{FCB_Ex.png}
	\caption{\textbf{(a)} A graph $G$ with edges colored according to whether they are inside (blue) or outside (red) the spanning tree $T$. 
%		There are two cycles forming a fundamental cycle basis $\mathcal{F}$ each composed of an edge in $E(G)\setminus E(T)$ joined by a path in $T$. 
		The face cycles $B_1$ and $B_2$ form a fundamental cycle basis $\calF$, where they are composed of an edge in $E(G){\setminus} E(T)$ joined by a path in $T$.
		\textbf{(b)} A depiction of the cycles $B_1$ and $B_1\,{\oplus}\, B_2$ alongside their binary representation as fundamental cycle basis $\calF$ vectors. \textbf{(c)} Another spanning tree $T'$ resulting in a different fundamental cycle basis that contains $B_2'\,{=}\,B_1\,{\oplus}\, B_2$.}
	\label{fig:FCB}
\end{figure*}



%Here, the cycle space is the closure of circuits under the operation $\oplus$. 
This notion of a vector space leads to a natural definition of a basis. Specifically, a \textit{cycle basis} is a set of linearly independent cycles $\calB=\{B_1,\hdots,B_\nu\}$ that span the cycle space, where $\nu$ is the dimension of the cycle space. Then, we may expand a cycle $C$ as follows,
\begin{equation}
	C = b_1 B_1\oplus\hdots\oplus b_\nu B_\nu,\quad b_1,\hdots,b_\nu\in\{0,1\}.
	\label{eq:cyc_exp_format_full}
\end{equation}
Note, lower case letters are used for scalars, upper case for cycles, and math italics (e.g. $\calB$) for sets of cycles. Computationally, it is easier to store cycles as binary vectors $(b_1,\hdots,b_\nu)$ where the addition operator $\oplus$ is still interpreted as the element-wise XOR operation. 




Now that we know how to use a cycle basis, the next question is ``how do we construct such a basis?''
A classical approach is to utilize a spanning tree.



%Computationally, we can interpret this as a sparse representation of the $(b_1,\hdots,b_\nu)$ vector where we instead store the list of indices $i$ where $b_i=1$.
%{\color{red} some text on expanding a cycle w.r.t. $\calB$. End with:} 

\begin{definition}
	{Let $T$ be a spanning tree of $G$, or a forest if $G$ is disconnected, and $\{e_i\}$ the edges in $G$ outside of $T$, i.e., $e_i\in E(G)\setminus E(T)$. The \emph{fundamental cycle basis} induced by $T$ is the set of cycles $\calF=\{B_i\}$, where $B_i$ is the cycle composed of the edge $e_i$ and the path that connects its endpoints in $T$.}%The \textit{fundamental cycle basis} induced by $T$ is the set of cycles $\mathcal{F}$ constructed by joining an edge in $G$ outside $T$ to its unique path in $T$.
	\label{def:FCB}
\end{definition}
\noindent An example of a fundamental cycle basis is given in Fig.~\ref{fig:FCB} (a). To get the expansion of a cycle $C$ with respect to $\calF$ we focus on the edges $\{e_i\}$. Note, there exists a unique basis cycle $B_i\in\calF$ for each edge $e_i$. 
So, $B_i$ is in the expansion of $C$ with respect to $\calF$ if and only if $\{e_i\}$ is contained in $C$. Furthermore, the coefficient $b_i$ as defined in Eq.~\eqref{eq:cyc_exp_format_full} is the indicator that $e_i\in C$ and $(b_1,\hdots,b_\nu)$ is the incidence vector of $C$ restricted to edges in $E(G)\setminus E(T)$.
%So, the coefficient $b_i$ as defined in Eq.~\eqref{eq:cyc_exp_format_full} is the indicator that $e_i$ is contained in $C$. Furthermore, the vector $(b_1,\hdots,b_\nu)$ is the incidence vector of $C$ restricted to edges in $G\setminus T$. See Fig.~\ref{fig:FCB} (b). 
%{\color{red}-- this has a lot of names, e.g. cyclomatic number, but introducing more notation feels annoying --} 
Following this procedure, we find the rank of the cycle space is the number of edges in $E(G)\setminus E(T)$:
\begin{equation}
	\nu = m-n+(\text{\# Connected Components}),
\end{equation}
where a tree with $n$ nodes has $n-1$ edges, and $T$ has a spanning tree of each connected component of $G$. However, we have not yet verified that $\calF$ is a basis.


%To see that $\mathcal{F}$ is linearly independent observe that each basis cycle $B\in\mathcal{F}$ has an edge the other cycles in $\mathcal{F}$ do not -- namely the edge $B$ has in $G\setminus T$. To see that $\mathcal{F}$ spans the cycle space, we find the decomposition of an arbitrary cycle $C$.



It is easy to see that $\calF$ is linearly independent, since each cycle $B_i$ has an edge $e_i$ that the other cycles in $\calF$ do not. To see that $\calF$ spans $\CycSp$, we take an arbitrary cycle $C\in\CycSp$. Then, for each basis cycle $B_i\in\calF$ we add it to $C$ if the corresponding edge $e_i\in E(G)\setminus E(T)$ is in $C$:
\begin{equation}
	C' = C\oplus \bm{1}_{[e_1\in C]}B_1\oplus\hdots\oplus\bm{1}_{[e_\nu \in C]}B_\nu
\end{equation}
where $\bm{1}$ is the indicator function. By cancellation from the $\oplus$-operator, $C'$ contains no edges in $E(G)\setminus E(T)$. In other terms, $C'$ is contained in the spanning forest $T$. Trees lack loops, so $C'$ is necessarily the empty cycle $\emptyset$, and $C$ is contained in the span of $\calF$.

%Consider an edge $e$ in $C$ that lies outside the spanning tree, i.e. in $G\setminus T$. There is only one cycle $C_e\in\mathcal{F}$ that shares this edge with $C$, so $C_e$ is necessarily in the decomposition of $C$ with respect to $\mathcal{F}$. If for edge $e$ in $G\setminus T$ contained in $C$ we add the corresponding cycle $C_e\in\mathcal{F}$ to $C$ using the $\oplus$ operator, then by cancellation of the edges in $G\setminus T$ we achieve a cycle in $T$. The spanning tree $T$ has no loops, so it only contains the null cycle $\emptyset$. By this cancellation, we find that $C$ is contained in $\mathcal{F}$. 
%
%Furthermore, the decomposition of $C$ with respect to $\mathcal{F}$ is given by the edges it has in $G\setminus T$. Since $T$ is a spanning tree of a connected graph, it has $n-1$ edges. So, the number of edges in $G\setminus T$ is \begin{equation}\nu:=m-n+1.\end{equation}


%If we add all cycles that share an edge in $G\setminus T$ with $C$ to $C$ using the $\oplus$ operator, then all edges in $G\setminus T$ cancel and we obtain a cycle in $T$. A tree contains no loop so the only cycle in $T$ is the empty set $\emptyset$. Thus, $C$ is in the

% -- making the cycles in $\mathcal{F}$ linearly independent. To see that $\mathcal{F}$ spans the cycle space consider an arbitrary cycle $C$. Then, we add each cycle in $\mathcal{F}$ to $C$ with the $\oplus$ operator if it shares an edge in $G\setminus T$ with $C$. The resulting cycle has cancellation between all the edges in $G\setminus T$, so it is contained in $T$. A tree has no loops, so this cycle is necessarily empty $\emptyset$. Thus, we have shown that $C$ is contained in the span of $\mathcal{F}$. We have also computed the decomposition in $\mathcal{F}$.
%If we add each cycle in $\mathcal{F}$ that shares an edge in $G\setminus T$ with $C$ to $C$ using the $\oplus$ operator, then the edges in $G\setminus T$ cancel and we achieve a cycle in the tree $T$. A tree contains no loops so the only cycle in $T$ is the empty set $\emptyset$. 


\subsection{Minimum cycle bases and relevant cycles}

\label{sub:MCB+Vis_Fam}
%{\bf P1: Why need to care about cycle lengths.}
Fundamental cycle bases are useful for analyzing the cycle space as a vector space. However, these bases should not be used when sampling loop statistics. Specifically, fundamental cycle bases are not unique and depend on the spanning tree used to construct them. Furthermore, a given fundamental cycle basis may contain large cycles that are not intuitively relevant.

To elaborate on this issue consider the graph depicted in Figure~\ref{fig:FCB} with two hexagons that share an edge. The spanning tree in Figure~\ref{fig:FCB} (a) results in a fundamental cycle basis composed of these hexagons, as desired. However, the spanning tree in Figure~\ref{fig:FCB} (c) produces a fundamental cycle basis containing the larger cycle with 10 nodes surrounding the perimeter of this graph. We argue, that it is more natural to describe this larger cycle as the sum of smaller cycles.

%{\bf P2: Definition of MCB}
A common approach taken to address this concern is to take a cycle basis that contains only short cycles:
\begin{definition}
	A \emph{mimimum cycle basis (MCB)} $\calM$ is a cycle basis that minimizes the sum of cycle lengths
	\begin{equation*}
		{\sf Cost}(\calM) = \sum_{B\in\calM} |B|
	\end{equation*}
	where $|B|$ is the number of edges in cycle $B$.
\end{definition}
\noindent As one might hope, the two hexagons in Figure~\ref{fig:FCB} form a unique MCB.

This definition has an interesting history in many fields. Ref.~\cite{kavithaCycle2009} offers a survey of algorithms to compute MCBs. 
Ref.~\cite{bergerCounterexamples2004} reviews methods for counting cycles in chemical systems. 
%for a review of algorithms used to compute minimum cycle bases and Ref.~\cite{bergerCounterexamples2004} for a review of these definitions for chemical systems. 
We remark that, in chemical systems, a cycle is typically called a ring, and an MCB is referred to as the ``Smallest Set of Smallest Rings''. In Section~\ref{sub:MCB_comp}, we review an algorithm from Ref.~\cite{kavithaFaster2004} for computing MCBs. This algorithm runs in $O(mn^2)$ time. In Section~\ref{sub:MCB_COB}, we show how one can expand an arbitrary cycle in terms of an MCB output by this algorithm. 

\begin{figure}[t]
	\centering
	\includegraphics{BaralleneGraphic.pdf}
%	\begin{tikzpicture}[scale = 1.25,
%		reg_node/.style={circle, draw, fill=white,
%			inner sep=0pt, minimum size = 12pt},
%		reg_edge/.style={line width = 2pt},
%		loop_edge/.style={line width = 2pt,color=red},
%		loop_arrow/.style={-Latex,thick,black}]
%		{\footnotesize
%			
%%		\node[anchor=south west] () at (-3pt-6cm,-.65in-3pt) {\includegraphics[height=2in]{Figures/Barallene_Print.jpg}};
%		\node[] () at (-4.5,.4) {\includegraphics[height=5cm]{Figures/Barallene_Print.jpg}};
%		
%		\fill[black!10] (-6.5,1.9) rectangle (-6,2.4);
%		\node () at (-6.25,2.15) {\normalsize(a)};
%		
%		\fill[black!10] (-2.5,1.9) rectangle (-2,2.4);
%		\node () at (-2.25,2.15) {\normalsize(b)};
%		
%		\node[namednode] (a) at (0,0) {1};
%		\node[namednode] (b) at (1,0) {2};
%		\node[namednode] (c) at (-.5,0.883) {3};
%		\node[namednode] (d) at (-.5,-0.883) {4};
%		\node[namednode] (e) at (.75,.5) {8};
%		\node[namednode] (f) at (1.75,.5) {5};
%		\node[namednode] (g) at (0.25,1.383) {6};
%		\node[namednode] (h) at (0.25,-0.383) {7};
%		\draw[namednode] (a)--(b)--(f)--(e)--(g)--(c)--(a)--(d)--(h)--(e);
%		}
%		\node () at (0.25,0.5) {$G$};
%		
%		\def\dxo{1.25}
%		\def\dxt{1.116}%.25+.866}
%		\node () at (\dxo+1.5,\dxt+.5/3) { $C_1$};
%		\node[regnode] (1a) at (\dxo,\dxt) {};
%		\node[regnode] (1b) at (\dxo+1*2/3,\dxt) {};
%		\node[regnode] (1c) at (\dxo-.5*2/3,\dxt+0.883*2/3) {};
%%		\node[regnode] (1d) at (\dxo-.5/2,\dxt-0.883/2) {};
%		\node[regnode] (1e) at (\dxo+.75*2/3,\dxt+.5*2/3) {};
%		\node[regnode] (1f) at (\dxo+1.75*2/3,\dxt+.5*2/3) {};
%		\node[regnode] (1g) at (\dxo+0.25*2/3,\dxt+1.383*2/3) {};
%%		\node[regnode] (1h) at (\dxo+0.25/2,\dxt-0.383/2) {};
%		\draw[] (1a)--(1b)--(1f)--(1e)--(1g)--(1c)--(1a);
%		
%		\node[regnode] (2a) at (\dxo,-\dxt+.5) {};
%		\node[regnode] (2b) at (\dxo+1*2/3,-\dxt+.5) {};
%%		\node[regnode] (2c) at (\dxo-.5*2/3,\dxt+0.883*2/3) {};
%		\node[regnode] (2d) at (\dxo-.5*2/3,-\dxt-0.883*2/3+.5) {};
%		\node[regnode] (2e) at (\dxo+.75*2/3,-\dxt+.5*2/3+.5) {};
%		\node[regnode] (2f) at (\dxo+1.75*2/3,-\dxt+.5*2/3+.5) {};
%%		\node[regnode] (2g) at (\dxo+0.25*2/3,\dxt+1.383*2/3) {};
%		\node[regnode] (2h) at (\dxo+0.25*2/3,-0.383*2/3-\dxt+.5) {};
%		\draw (2a)--(2b)--(2f)--(2e)--(2h)--(2d)--(2a);
%		\node () at (\dxo+1.5,-\dxt+.5/3+.5) { $C_2$};
%		
%		
%		\node[regnode] (3a) at (-1.25,0) {};
%%		\node[regnode] (1b) at (\dxo+1*2/3,\dxt) {};
%		\node[regnode] (3c) at (-1.25-.5*2/3,0.883*2/3) {};
%		\node[regnode] (3d) at (-1.25-.5*2/3,-0.883*2/3) {};
%		\node[regnode] (3e) at (-1.25+.75*2/3,.5*2/3) {};
%%		\node[regnode] (1f) at (\dxo+1.75*2/3,\dxt+.5*2/3) {};
%		\node[regnode] (3g) at (-1.25+0.25*2/3,1.383*2/3) {};
%		\node[regnode] (3h) at (-1.25+0.25*2/3,-0.383*2/3) {};
%		\draw (3a)--(3c)--(3g)--(3e)--(3h)--(3d)--(3a);
%		\node () at (-1.75,0.5/3) { $C_3$};
%		
%		\draw[thick] (-6.5,-1.6) rectangle (3.5,2.4) 
%		(-2.5,-1.6)--(-2.5,2.4)
%		(-6,1.9) rectangle (-6.5,2.4)
%		(-2,1.9) rectangle (-2.5,2.4);% (-6.5,.25)--(3.5,.25);
%	\end{tikzpicture}
	\caption{
		A graph $G$ that does not have a unique MCB. Ignoring double bonds, $G$ is the carbon skeleton of barallene with formula $\rm C_8 H_8$. \textbf{(a)} 3D printed visualization of $G$. \textbf{(b)} Two-dimensional representation of $G$. The three hexagonal relevant cycles $C_1$, $C_2$, and $C_3$ are shown separately. Any pair of these cycles forms an MCB, where the third cycle is the sum of the other two, e.g., $\calM_1=\{C_1,C_2\}$ is an MCB and $C_3=C_1\oplus C_2$.
%		A graph $G$ that does not have a unique MCB. $G$ is the carbon skeleton of {bicyclo[2,2,2]octane}.  
%		(a) 3D printed visualization of the graph to show that its structure is best visualized as a pinwheel shape. (b) Classical depiction of the graph with labeled nodes. In (b), we separately depict the three hexagonal relevant cycles in $G$:
%		$C_1=(1,2,5,8,6,3)$, $C_2=(1,2,5,8,7,4)$, and $C_3=(1,3,6,8,7,4)$. 
%		Any pair of these cycles forms an MCB, where the third cycle is the sum of the other two, e.g., $\calM_1=\{C_1,C_2\}$ is an MCB and $C_3=C_1\oplus C_2$. 
	}
	\label{fig:MCB_Unique}
\end{figure}

Unfortunately, MCBs are not always unique -- see Figure~\ref{fig:MCB_Unique}. Therefore, if we evaluate summary statistics from cycles in a given MCB, then we obtain inconsistent results depending on which MCB is used. For instance, consider the number of times node 2 in Figure~\ref{fig:MCB_Unique} participates in a cycle. If we compute this statistic using the MCB $\calM_1=\{C_1,C_2\}$ then node 2 participates in two cycles $C_1$ and $C_2$, but if we take $\calM_2=\{C_2,C_3\}$ as an MCB then node 2 participates in only one cycle $C_2$. As mentioned in Ref.~\cite{kolodzikUnique2012}, this is an issue with pattern matching of chemical rings in SMARTS, {a programming library for identifying molecular patterns and properties}~\cite{SMARTS}. Additionally, the OEChem {library for chemistry and cheminformatics} avoids using MCBs in chemical ring perception citing lack of uniqueness, and states ``Smallest Set of Smallest Rings (SSSR) considered Harmful''~\cite{SSSR_Harmful}.% {\color{red} short clarification }

%{\color{red} change outline, because of this... We do not summarize all of these works here, and refer ... Instead, we discuss the notion of relevant cycles from (Vismara)...}
Because of this issue, several works in computational chemistry consider definitions that extend MCBs to larger sets of cycles. We do not review all of these definitions here, and refer the interested reader to Ref.~\cite{bergerCounterexamples2004}. Instead, we review the notion of relevant cycles from Ref.~\cite{vismaraUnion1997}, which serves as a foundation for the definitions discussed in this manuscript. The relevant cycles have two equivalent definitions:
%These cycles are defined as follows
%{\color{red} Move 2 definitions here}
\begin{definition}
	The set of relevant cycles $\RelCyc$ is the union of all MCBs.\label{def:rel}
\end{definition}
\begin{definition}
	A cycle $C\in\CycSp$ is relevant ($C\in\RelCyc$) if and only if it is not generated by smaller cycles, i.e., there does not exist a set of smaller cycles $\calX$, $|C'|<|C|$ for $C'\in\calX$, such that
	\begin{equation}
		C = \bigoplus_{C'\in\calX} C'
	\end{equation}
	where $\bigoplus$ is the series representation of the $\oplus$-operator.
	\label{def:rel_2}
\end{definition}
% move forward
%\noindent Alternatively, the relevant cycles are the cycles that cannot be written as the sum of a set of smaller cycles. We verify this equivalent definition at the end of this subsection. 
We include a proof of the equivalence of these definitions after introducing additional theory about the cycle space -- see Corollary~\ref{cor:rel_def_eq} in Sec.~\ref{sub:basis_ex}. For the graph depicted in Figure~\ref{fig:MCB_Unique}, the relevant cycles are the three hexagons $\RelCyc=\{C_1,C_2,C_3\}$. This set of relevant cycles has several useful properties~\cite{bergerCounterexamples2004}. % {(\color{red} This list uses the same headings as the review cited, should this be avoided?)}
\begin{itemize}
	\item \textit{Uniqueness}: The set $\RelCyc$ is unique, so statistics on the relevant cycles are well-defined.
	\item \textit{Completeness}: The relevant cycles can be used to generate the cycle space since $\RelCyc$ contains a cycle basis.
	\item \textit{Computability}: There exists a polynomial time algorithm for computing the set of relevant cycles.~\cite{vismaraUnion1997}.
	\item \textit{Chemical relevance}: 
	The set of relevant cycles contains all building blocks of chemical structures present in molecules. 
\end{itemize}


One downside is that the number of relevant cycles may grow exponentially for certain types of graphs.
For example, in the bracelet-like graph composed of $k$ diamonds chained together into a loop shown in Figure~\ref{fig:intro_graphs} (b). 
All MCBs in this graph are composed of the $k$ diamonds and one large loop of length $3k$. Each large loop either takes the top or bottom path of a given diamond, so the number of large relevant cycles grows exponentially as $2^k$, while the number of nodes grow linearly as $4k$.~\cite{vismaraUnion1997}
	
It is remarkable that the algorithm from Ref.~\cite{vismaraUnion1997} for computing the relevant cycles of a graph runs in polynomial time.
In broad steps, the relevant cycles are computed by (1) generating a candidate set of cycles (a superset of $\RelCyc$) and (2) filtering out the cycles that are irrelevant. The key idea that makes this procedure feasible in polynomial time is that the cycles are grouped into families of equal length. 

%See Fig.~\ref{fig:Vis_Fams} for their description.% as depicted in Fig.~\ref{fig:Vis_Fams}.

\begin{figure}[t]
	\centering
	\includegraphics{VFam.pdf}
	%	\includegraphics[width=.5\linewidth]{Vis_Cyc_Fam_Graphic.png}
	\caption{%Method used by Vismara for constructing candidate relevant cycles~\cite{vismaraUnion1997}. If we take a root node $r$ in the cycle and look at the opposite end of the cycle we find two nodes $p$ and $q$ equidistant to $r$ joined by either an edge or a common node $x$.
%		Depiction of the families of candidate relevant cycles introduced by Vismara~\cite{vismaraUnion1997}. Families of odd length cycles are indexed by a root node $r$ and edge $(p,q)$ at the opposite end of the cycle. Cycles in an odd family consist of a shortest path from $r$ to $p$ (red line), a shortest from $r$ to $q$ (blue line), and the edge $(p,q)$ (black line). 
%		The construction of families of even cycles is similar except $p$ and $q$ are separated by an intermediate node $x$.
%		A cycle family may have multiple cycles if there are multiple shortest paths from $r$ to $p$ and $q$. The size of a family is the number of shortest path pairs, i.e., the number of shortest paths from $r$ to $p$ multiplied by the number of shortest paths from $r$ to $q$.
%		Such a family may have many cycles since there are potentially many shortest paths from $r$ to $p$ and from $r$ to $q$. 
%		The construction of families of even cycles is similar except $p$ and $q$ are separated by an intermediate node $x$.% so that shortest paths from $r$ to $p$ and $q$ are equal in length which is less than half the length of the cycle.
%		-----
%		Shorter: (1) depiction of the descriptors used to organize cycles into families, (2) there may be additional cycles with the same descriptors
		Depiction of the descriptors used to organize the relevant cycles into V-families as in Ref.~\cite{vismaraUnion1997}. 
%		For an odd length cycle (left), the destcriptors are the largest index node $r$ and opposing edge $(p,q)$. For an even cycle (right), they are the largest index node $r$ and opposing triple $(p,x,q)$.
		For an odd length cycle (left), we fix its largest index node $r$ and the edge $(p,q)$ at the opposite end of the cycle. For an even length cycle (right), we fix its largest index node $r$ and the triple $(p,x,q)$ at the opposite end of the cycle. 
		In the left, we also draw the exchange of the path $P_1$ for $P_2$ by adding the cycle $C_e=P_1\oplus P_2$ to $C_1$ to obtain $C_2=C_1\oplus C_e$.
	}
	\label{fig:Vis_Fams}
\end{figure}

\subsubsection{Vismara's families of cycles}

%We review the construction of cycle families from Ref.~\cite{vismaraUnion1997}.

Here we review the construction of cycle families from Ref.~\cite{vismaraUnion1997}. 
The key idea to make the algorithm for computing the relevant cycles polynomial is to group the relevant cycles into a polynomial number of families.
%
%This definition results in a polynomial time algorithm that groups the relevant cycles into a polynomial number of families. 
At the same time, it has a theoretical drawback that it is attached to a given enumeration of the vertices.

Consider a simple cycle $C$. We select antipodal parts of $C$ to determine the family it belongs to.
%To determine the family $C_1$ belongs to, we fix antipodal objects in $C_1$. 
First, we fix the node $r$ with largest index in $C$. The largest index node is selected because it is unique up to vertex labels. Second, if $C$ is an odd length cycle, then we fix the edge $(p,q)$ at the opposite end of the cycle. If $C$ has even length, then we fix the triple $(p,x,q)$ at the opposite end of the cycle. In either case, the paths from $r$ to $p$ and $q$ in $C$ are of equal length. See Figure~\ref{fig:Vis_Fams}. We call these families V-families. %The cycle family $C$ belongs to is simply the set of cycles that share its descriptors $r$, $p$, $q$ (and $x$ if $|C|$ is even).% $r$ and $(p,q)$ or $(p,x,q)$.

\begin{definition}[Vismara cycle family]\phantom{.}
	%{\color{red}bulleted list}
	
	\begin{itemize}
		\item The odd V-family with descriptors $r,p,q$ is the set of cycles with largest index node $r$ consisting of shortest paths from $r$ to $p$ and $q$ joined by the edge $(p,q)$.
		
		\item The even V-family with descriptors $r,p,x,q$ is the set of cycles with largest index node $r$ consisting of shortest paths from $r$ to $p$ and $q$ joined by the triple $(p,x,q)$.
	\end{itemize}
	
%	Let $C$ be a simple cycle with an odd number of vertices, $r$ its node with largest index and $(p,q)$ the edge opposite to this node. The Visamra cycle family with descriptors $r,p,q$
	
	
%	The odd Vismara cycle family associated with a given node $r$ and edge $(p,q)$ is the set of cycles with largest index node $r$ composed of two shortest paths one from $r$ to $p$ and one from $r$ to $q$ joined into a loop by the edge $(p,q)$. The even cycle family associated with a given node $r$ an triple $(p,x,q)$ is the set of cycles with largest index node $r$ composed of shortest paths from $r$ to $p$ and $q$ joined by the triple $(p,x,q)$.
	
	\label{def:vis_fam}%\color{red} Use Vismara's phrasing for not passing through later nodes
\end{definition}

\noindent All cycles in a V-family are of equal length. We only consider shortest paths from $r$ to $p$ and $q$ because it is a necessary condition for a cycle to be relevant -- see Lemma~\ref{lem:shortest_paths}.%It is not immediately obvious why we restrict ourselves to shortest paths from $r$ to $p$ and $q$. We will return to this later in Lemma~\ref{lem:shortest_paths}.

\begin{figure}[t]
	\centering
	\includegraphics{VFamEx.pdf}
	\caption{A labeled graph composed of three diamonds chained into a loop. The relevant cycles of this graph are the three diamonds and eight large loops of length nine. For the given node labels, the large cycles form three V-families depicted separately with their descriptors.}
	\label{fig:three_diamonds}
\end{figure}


%There is a polynomial number of cycle families with respect to the number of nodes. %The construction of these families depends on node labels. 
%As an example, the graph depicted in Figure~\ref{fig:exp_rel_cyc} (b), the $2^3$ large cycles of length nine are split into three families
As an example, consider the graph depicted in Figure~\ref{fig:three_diamonds}. 
This graph has eight large relevant cycles of odd length nine, distinguished by the path they take through each diamond.
%This graph has $2^3=8$ larger cycles of length nine.
Four of these cycles take the path $(10,12,11)$ in the 
%{\color{red} \st{right diamond }} 
diamond $(9,11,12,10)$
making their largest index node $r=12$. 
%Four of these cycles have the largest index node $r=12$ in the graph, 
%each of which can be identified by the path they take in the upper and left diamonds. %taking either path in the left and upper diamonds.
All four cycles share the same opposing edge $(2,5)$ and belong to the same family. The remaining large relevant cycles take $(10,9,11)$ instead of $(10,12,11)$ in the diamond $(9,11,12,10)$, so they have largest index node $r=11$. The opposing edge for these remaining cycles is $(5,6)$ if they take the path $(5,6,8)$ in the diamond $(5,6,8,7)$ or $(5,7)$ if they take the path $(5,7,8)$. This results in two additional V-families, each of which has two cycles, one for each path in the diamond $(1,2,4,3)$.%the $2^3=8$ large cycles of length nine in Figure~\ref{fig:three_diamonds} (b) are split into three families:
%\begin{itemize}
%	\item $r=12$, $(p,q)=(2,5)$, four cycles,
%	\item $r=11$, $(p,q)=(5,6)$, two large cycles passing through node $6$ in the top diamond and node $9$ in the bottom right diamond,
%	\item $r=11$, $(p,q)=(5,7)$, two large cycles passing through node $7$ in the top diamond and node $9$ in the bottom right diamond.
%\end{itemize}
%These families can be hard to interpret, and their construction depends on node labels. Thus, other works consider alternative partitions of the relevant cycles that are easier to interpret and unique~\cite{gleissInterchangeability2000,kolodzikUnique2012}. We discuss these in Sec.~\ref{sub:Int_Back}. Nevertheless, the relevant cycles and Vismara's cycle families are important as the foundation for these definitions.
%It can be hard to interpret these output cycle families, so other works~\cite{gleissInterchangeability2000,kolodzikUnique2012} consider alternate groupings of the relevant cycles, which we discuss in Sec.~\ref{sub:Int_Back}. Nevertheless, Vismara's cycle families are useful as they are the computational foundation for these definitions.% are a useful starting point for these definitions.

%If we swap the labels of node $$ and $16$ in Figure~\ref{fig:exp_rel_cyc} (b), then the large cycles are split into 7 families.

%We give two examples to clarify how the relevant cycles are grouped into families. For the graph depicted in Figure~\ref{fig:MCB_Unique} (b), the hexagons $C_1$, $C_2$, $C_3$ form separate families with a shared node $r=8$ and distinct triples $(2,1,3)$, $(2,1,4)$, and $(3,1,4)$ respectively.

%Let us consider a simple cycle of odd number of edges. Let $r$ be a arbitrarily selected node and $(p,q)$ the opposite edge.
%Now we will review how the cycles are split into families that Vismara's algorithm runs in polynomial time.

Ultimately, splitting cycles into V-families as in Definition~\ref{def:vis_fam} is effective because shortest paths can be exchanged via cycle addition. For instance, let us consider two cycles $C_1$ and $C_2$ in a V-family that differ by taking disjoint paths $P_1$ and $P_2$ from $r$ to $p$ respectively. See, e.g., Figure~\ref{fig:Vis_Fams}, left. Additionally, consider the cycle $C_e$ obtained by merging these paths, given by%\footnote{We note, $C_e=P_1\oplus P_2$ is a cycle, even when $P_1$ and $P_2$ are not disjoint. This is verified by showing nodes in $C_e$ have even degree. First, consider a endpoint node $u$ that participates in one edge $e_1$ in $P_1$ and one edge $e_2$ in $P_2$.  If $e_1$ and $e_2$ are distinct, then the degree of $u$ is $\text{deg}(u)=2$. Otherwise, if $e_1=e_2$, then the edge cancels through the symmetric difference, and $\text{deg}(u)=0$. In both cases, $\text{deg}(u)$ is even as desired. Next, consider a node $v$ that is in the interior of $P_1$ or $P_2$. Let $A$ and $B$ be the respective edges in $P_1$ and $P_2$ that contain $v$. $A$ and $B$ contain an even number of edges: for the edges in $P_1$, $|A|=0$ if $v$ does not lie on $P_1$, and $|A|=2$ otherwise. From this, it can be shown that $v$ has even degree using the inclusion-exclusion principle $\text{deg}(v)=|A\oplus B| = |A|+|B|-2|A\cap B|$.}
%Ultimately, splitting cycles into families is effective because shortest paths can be exchanged via the addition of smaller cycles. We limit our discussion to paths from $r$ to $p$ in odd families as depicted in Fig.~\ref{fig:Vis_Fams}, left. The steps for all other cases are similar.
%-----
%Let us consider $C_1$ contain $r$ and $(p,q)$... Let $C_2$ be a cycle except taking...
%-----
%\st{In particular,} consider the cycles $C_1$ and $C_2$ in which take distinct paths $P_1$ and $P_2$ from $r$ to $p$ respectively. Here, $P_1$ and $P_2$ should be interpreted as sets of edges, like cycles. 
%Additionally, consider the cycle $C_e$ obtained by merging these paths. Formally, $C_e$ is given by the symmetric difference\footnote{By this definition, $C_e$ is a cycle, even when the paths $P_1$ and $P_2$ are not disjoint. This is verified by showing the nodes in $C_e$ have even degree. First consider a endpoint node $u$ that participates in one edge $e_1$ in $P_1$ and one edge $e_2$ in $P_2$.  If $e_1$ and $e_2$ are distinct, then the degree of $u$ is 2. Otherwise, if $e_1=e_2$, then the edge cancels through the symmetric difference, and the degree of $u$ is 0. In both cases, $\text{deg}(u)$ is even as desired. Next, consider a node $v$ that is in the interior of $P_1$ or $P_2$. Let $A$ and $B$ be the respective edges in $P_1$ and $P_2$ that contain $v$. $A$ and $B$ contain an even number of edges: for the edges in $P_1$, $|A|=0$ if $v$ does not lie on $P_1$, and $|A|=2$ otherwise. From this, it can be shown that $v$ has even degree using the inclusion-exclusion principle $\text{deg}(v)=|A\oplus B| = |A|+|B|-2|A\cap B|$.}
\begin{equation}
	C_e \coloneq P_1\oplus P_2.\label{eq:Ce_Path_Merge}
\end{equation}
%in case $P_1$ and $P_2$ overlap. 
%-----
%Discuss situation for unique shortest path
%-----
Then, we exchange $C_1$ for $C_2$ as follows
\begin{equation}
	C_1 = C_2 \oplus C_e\label{eq:Ce_Ex_Vis}
\end{equation}
where the path $P_2$ cancels between $C_2$ and $C_e$ and is replaced by $P_1$. %Note, $P_1$ is less than half the length of $C_1$, because the paths from $r$ to $p$ and $q$ are equal length, and $C_1$ has an additional edge $(p,q)$. 
%Thus, $C_e$ is smaller than $C_1$ and $C_2$
%\begin{equation}
%	|C_e|\le |P_1| + |P_2| = 2|P_1| < |C_1|.
%\end{equation}

%\textbf{New Par.?}
%The exchange of cycles via the addition of shorter ones can be used to show two useful properties about cycle families for computing $\RelCyc$.
In general, one cycle $C_1$ in a family can be exchanged for another cycle $C_2$ in the same family by adding at most two cycles. Specifically, one cycle is used to swap paths from $r$ to $p$, and the other is used to swap paths from $r$ to $q$. Note, these exchange cycles are shorter than the equal length cycles $C_1,C_2$, e.g., for $C_e$ given by Eq.~\eqref{eq:Ce_Path_Merge}
\begin{equation}
	|C_e|\le |P_1|+|P_2|=2|P_1|<|C_1|.\label{eq:Ce_Ex_Len}
\end{equation}
This exchange procedure can be used to show two useful properties for computing cycle families.

First, this can be used to show that the paths from $r$ to $p$ and $q$ in a relevant cycle are shortest paths, because otherwise they can be exchanged for shortest paths. The proof of this claim relies on the fact that the paths from $r$ to $p$ and $q$ are less than half the length of the cycle. 



\begin{lemma}[Ref.~\cite{vismaraUnion1997}, Lemma 2] %\cite[Lemma 1]{vismaraUnion1997}]
	Let $C$ be a relevant cycle and $P$ a path in $C$ that is at most half the length of $C$, $|P|\le|C|/2$. Then, $P$ is a shortest path.
	\label{lem:shortest_paths}
\end{lemma}

\begin{proof}
	Suppose towards a contradiction that $P$ is not a shortest path. Let $P'$ be a shortest path with the same endpoints, such that $|P'|<|P|$. Let $C_e$ be the cycle used to exchange paths $P$ and $P'$ given by
	\begin{equation}
		C_e = P\oplus P'.
	\end{equation}
	Let $C'$ be given by
	\begin{equation}
		C' \coloneq C\oplus C_e = (C\setminus P)\oplus P'.
	\end{equation}
	where the second equality uses associativity and that the path $P$ belongs to $C$. Then, our initial cycle is given by the sum
	\begin{equation}
		C = C'\oplus C_e.
	\end{equation}
	Both $C'$ and $C_e$ are strictly smaller than $C$
	\begin{align}
		|C_e| &\le |P|+|P'|< 2|P|\le |C|,\\
		|C'| &\le |C| - |P| + |P'| < |C|.
	\end{align}
	Therefore, $C$ is not relevant by Definition~\ref{def:rel_2}, because it is the sum of shorter cycles. This contradicts our initial assumptions, so $P$ is a shortest path.
\end{proof}

%First, this exchange procedure justifies why the path $P_1$ from $r$ to $p$ needs to be a shortest path. To see this, suppose instead that $P_1$ is not a shortest path and that $|P_2|<|P_1|$. Then, the cycle $C_2$ is shorter than $C_1$, and $C_1$ is the sum of shorter cycles by Eq.~\eqref{eq:Ce_Ex_Vis}. Thus, if $P_1$ is not a shortest path then $C_1$ is not relevant by Definition~\ref{def:rel_2}. 


Shortest paths from a given root node $r$ to the other vertices can be obtained by breadth-first search. 
Neglecting some detail -- how to avoid nodes with index greater than $r$ -- the algorithm for generating candidate cycle families is to perform breadth-first search on all vertices $r\in V$ and identify node pairs $p$ and $q$ that form structures as in Fig.~\ref{fig:Vis_Fams}. 
We discuss this algorithm and our modifications to it in Sec.~\ref{sub:Vis_Fam}.

%------
%Perhaps add a formal statement
%------

Another benefit of this exchange procedure is that we only need to test if one cycle in a V-family is relevant. In Ref.~\cite{vismaraUnion1997}, a single cycle \textit{prototype} is taken arbitrarily from a V-family. If this prototype is relevant, then all cycles in the V-family are relevant. 
This is computationally feasible to test because there are a polynomial number of cycle prototypes.
%This makes testing relevance of the V-families computationally feasible, since we only need to test if a polynomial number of cycle prototypes are relevant. %The proof of this relies on the cycles used to exchange shortest paths being shorter than the cycles in the family, see Eq.~\eqref{eq:Ce_Ex_Len}. 
In the following lemma, $C_1$ and $C_2$ are cycles in a family, and $\calX$ is a set of two cycles used to exchange the paths from $r$ to $p$ and $q$. The cycles $\calX$ used to exchange paths are shorter than $C_1$ and $C_2$. See Eq.~\eqref{eq:Ce_Ex_Len}.

\begin{lemma}
	Let $C_1,C_2\in\CycSp$ be equal length cycles that differ by smaller cycles, i.e., there exists a set of cycles $\calX$, $|C'|<|C_1|$ for $C'\in\calX$, such that
	\begin{equation}
		C_1 = C_2\oplus\bigoplus_{C'\in\calX} C'.\label{eq:C1_C2_sm_c}
	\end{equation}
	Then, $C_1$ is relevant if and only if $C_2$ is relevant. 
\end{lemma}

\begin{proof}
%	Let $C_1,C_2\in\CycSp$, $|C_1|=|C_2|$, differ by smaller cycles, and 
	Let $C_2$ not relevant. We want to show that $C_1$ is not relevant. By Definition~\ref{def:rel_2}, $C_2$ is the sum of smaller cycles $\calY$, $|C'|<|C_2|$ for $C'\in\calY$, such that
	\begin{equation}
		C_2 = \bigoplus_{C'\in\calY}C'.
	\end{equation}
	Then, by Eq.~\eqref{eq:C1_C2_sm_c}, $C_1$ is given by the pair of sums
	\begin{equation}
		C_1 = \bigoplus_{C'\in\calY} C' \oplus \bigoplus_{C''\in\calX}C''.
	\end{equation}
	Cycles in the intersection $C'\in\calX\cap\calY$ cancel in the above double sum, since $C'\oplus C'$. Thus, $C_1$ is given by the single sum
	\begin{equation}
		C_1 = \bigoplus_{C'\in\calX\symdiff\calY} C'
	\end{equation}
	where $\symdiff$ is the symmetric difference $\calX\symdiff\calY = (\calX\setminus\calY)\cup(\calY\setminus\calX)$ for more general sets. Finally, $C_1$ is not relevant by Definition~\ref{def:rel_2}, because it is the sum of smaller cycles $\calX\symdiff\calY\subseteq\calX\cup\calY$. 
	
	To see the reverse statement, we note that the relationship given by Eq.~\eqref{eq:C1_C2_sm_c} is symmetric. In particular, we can add cycles $C'\in\calY$ to both sides of Eq.~\eqref{eq:C1_C2_sm_c} to move them to the left-hand side
	\begin{equation}
		C_1 \oplus \bigoplus_{C'\in\calY} C' = C_2.
	\end{equation}
	By the same argument as above, if $C_1$ is not relevant, then $C_2$ is not relevant.
\end{proof}







\subsection{Lemmas about exchanging basis cycles}
\label{sub:basis_ex}

	Much of the theory in this text relies on the procedure of swapping a given cycle in a basis for another cycle. This notion of basis exchange can be used to define a matroid, but it is easier to understand these properties as a consequence of the cycle space being a vector space. In this section, we state this notion of basis exchange in Lemma~\ref{Lem:Bas_Ex}; construct conditions for a cycle to be relevant in Lemma~\ref{Lem:Relevant_MCB_Test}; and verify that Definitions~\ref{def:rel} and~\ref{def:rel_2} are relevant in Corollary~\ref{cor:rel_def_eq}.
%	that the definitions of relevance, Definitions~\ref{def:rel} and~\ref{def:rel_2}, are equivalent in Corollary~\ref{cor:rel_def_eq}. %Here, we will use Definition~\ref{def:rel} as our initial definition of relevance.

%{\color{red} Introduce somehow:\\
	
	
%	{\color{red} Make this into a definition.} 
	For these lemmas, it is useful to introduce further notation for expanding a cycle with respect to a basis. Recall from Eq.~\eqref{eq:cyc_exp_format_full}, a cycle $C$ can be expanded as the sum of basis cycles $B_i$ with binary coefficients $b_i$. Thus, $C$ is simply the sum of cycles $B_i$ with coefficients $b_i=1$. We define the following function to output these cycles. 
	\begin{definition}
		Let $\calB$ be a cycle basis. The \emph{expansion operator} $\expand{\calB}{\cdot}$ is the function on cycles $C\in\CycSp$ that outputs the unique subset $\expand{\calB}{C}\subseteq\calB$ that sums to $C$
		\begin{equation}
			C = \bigoplus_{\mathclap{B\in\expand{\calB}{C}}} B,\quad \forall C\in\CycSp.
		\end{equation}
		\label{def:ex_op}
	\end{definition}
	
%	We define the expansion operator $\expand{\calB}{C}\subseteq \calB$ to be the cycles in the basis $\calB$ used to expand $C$, i.e., those with coefficient $1$. %For instance, the cycle $C_3$ depicted in Figure~\ref{fig:cyc_space}(b), right, has expansion $\expand{\calB}{C_3}=\{B_1,B_4\}$. 
%	By definition, the expansion operator satisfies
%	Alternatively, we can remove the cycles in Eq.~\eqref{eq:cyc_exp_format_full} with coefficient $b_i=0$, and represent $C$ as the sum of cycles $B_i$ with coefficient $b_i=1$. 
%	We find it convenient to introduce the expansion operator $\expand{\calB}{C}\subseteq \calB$ as the set of cycles in the basis $\calB$ used to represent $C$. In particular, this function satisfies
	%We introduce the expansion operator $\expand{\calB}{C}\subseteq \calB$ to be these basis cycles, such that
%	\begin{equation}
%		C = \bigoplus_{\mathclap{B\in\expand{\calB}{C}}} B,\quad \forall C\in\CycSp.
%	\end{equation}
%%}
%	\add{Additionally, the expansion of the sum of cycles is equal to the symmetric of their expansions
%	\begin{equation}
%		\expand{\calB}{C_1\oplus C_2} = \expand{\calB}{C_1}\symdiff\expand{\calB}{C_2}.
%		\label{eq:expand_sum}
%	\end{equation}
%	This identity is obtained by individually expanding $C_1$ and $C_2$ with respect $\calB$
%	\begin{equation}
%		C_1\oplus C_2 = \bigoplus_{{C'\in\expand{\calB}{C_1}}}C' \oplus \bigoplus_{{C''\in\expand{\calB}{C_2}}}C'' = \bigoplus_{{C'\in\expand{\calB}{C_1}\symdiff\expand{\calB}{C_2}}} C'
%	\end{equation}
%	where the final inequality is a consequence of cancellations of the form $C\oplus C=\emptyset$ between the two sums.}
	
	In the following lemma we show a cycle $C$ can be exchanged for cycles in its expansion $\expand{\calB}{C}$. {This well-known lemma has many equivalent forms, and is often referred to as the matroid property of the cycle space, see e.g., Proposition 4 in Ref.~\cite{gleissInterchangeability2000}.}

\begin{lemma}[Basis exchange]
	Let $\calB$ be a cycle basis, $C\in\CycSp$ be a cycle, and $B\in\expand{\calB}{C}$ be an arbitrary cycle in the expansion of $C$ with respect to $\calB$. If we replace $B$ with $C$ in $\calB$ then we obtain a new cycle basis $\calB' {=}\, (\calB{\setminus} \{B\}){\cup}\{C\}$.\nolinebreak % adds an e
	\label{Lem:Bas_Ex}
\end{lemma}
\begin{proof}
	To start, we expand $C\in\CycSp$ with respect to the basis $\calB$ as follows
	\begin{equation}
		C = B_1{\oplus} B_2 {\oplus}\hdots {\oplus} B_k, \enspace B_i{\in}\calB\text{, }1{\le} i{\le} k.\label{eq:C_exp_reg_BasEx}
	\end{equation}
	Note, in the above equation we fix $\expand{\calB}{C}=\{B_1,\hdots,B_k\}$. We want to show that if we swap $B_i$ with $C$ in $\calB$, then we obtain a new cycle basis. Without loss of generality, we will show that $\calB'\coloneq (\calB\setminus\{B_1\})\cup\{C\}$ is a cycle basis.
	To show this, we add $C$ and $B_1$ to both sides of Eq.~\eqref{eq:C_exp_reg_BasEx} to obtain
	\begin{equation}
		B_1 = C\oplus B_2\oplus\hdots\oplus B_k.\label{eq:B1_exp_reg_BasEx}
	\end{equation}
	Equation~\eqref{eq:B1_exp_reg_BasEx} shows that $B_1$ is in the span of $\calB'$, and that our original set $\calB$ is contained in the span of $\calB'$. $\calB$ is a cycle basis, so $\calB'$ must span the full cycle space. $C$ is linearly independent of $\{B_2,\hdots,B_k\}$, otherwise $B_1$ would be in the span of $\{B_2,\hdots,B_k\}$ by Eq.~\eqref{eq:B1_exp_reg_BasEx}. Thus, $\calB'$ is a basis.
\end{proof}

%	We now apply this property to minimum cycle bases to better understand the relevant cycles. Here, we take Definition~\ref{def:rel} as the primary definition of relevance, and we will show Definition~\ref{def:rel_2} is equivalent in Corollary~\ref{cor:rel_def_eq}. 
	
	{In the following lemma, we introduce an alternate definition of relevance in terms of the expansion of a cycle with respect to a MCB. We will use this lemma as a backbone for our theoretical results in this manuscript and for testing the relevance of cycles in V-families.
%	Additionally, we use Lemma~\ref{Lem:Relevant_MCB_Test} to algorithmically test if a cycle is relevant. 
	% the relevance of a prototype cycle in a Vismara cycle family.
	}% This approach is somewhat different than the one taken in Ref.~\cite{vismaraUnion1997}, which uses Gaussian elimination to test relevance, and outputs a minimum cycle basis.

\begin{lemma}[Relevance criterion]
	Let $C\in\CycSp$ be a cycle, $\calM$ an MCB, and
	$$L = \max\{|C'|:C'\in\expand{\calM}{C}\}$$
	the length of the largest cycle in the expansion of $C$ with respect to $\calM$.
	Then, the following are true
	\begin{enumerate}
		\item $L\le |C|$, i.e., no cycles in $\expand{\calM}{C}$ are larger than $C$,
		\item $L=|C|$ if and only if $C$ is relevant $C\in\mathcal{C_R}$. Furthermore, if $C\in\mathcal{C_R}$, then for any $C'\in\expand{\calM}{C}$, $|C'|=|C|$, the set $\calM'=(\calM\setminus\{C'\})\cup\{C\}$ is an MCB.
	\end{enumerate}
	\label{Lem:Relevant_MCB_Test}
\end{lemma}
\begin{proof}
	%	We begin by proving statement 1.
	(1) Let $C'\in\expand{\calM}{C}$ be an arbitrary cycle in the expansion of $C$ with respect to $\calM$. By Lemma~\ref{Lem:Bas_Ex}, $\calB=(\calM\setminus \{C'\})\cup\{C\}$ is a cycle basis. The sum of cycle lengths in $\calB$ is
	\begin{equation}
		\sum_{B\in\calB} |B| = \sum_{B\in\calM} |B| - |C'| + |C|.
	\end{equation}
	$\calM$ is an MCB, so the sum of cycle lengths in $\calB$ is greater than or equal to the sum of cycle lengths in $\calM$, and by the above equation
	\begin{equation}
		\begin{split}
		\sum_{B\in\calM} |B| - |C'| + |C| \ge \sum_{B\in\calM} |B| \quad \\\longrightarrow \quad |C| \ge |C'|.
		\end{split}
	\end{equation}
	Thus, $C$ is not smaller than any of the cycles in its expansion with respect to $\calM$ as claimed.
	
	(2, $\Rightarrow$) 
		Let $\calM'$ be an MCB containing $C$. 
		Let $C'\in\expand{\calM}{C}$ be independent of $\calM\setminus\{C\}$. By item 1, $|C'|\le |C|$. 
		Then, $\calB=\calM'\setminus\{C\}\cup\{C'\}$ is another cycle basis.
		\begin{equation}
			\sum_{\mathclap{B\in\calM'}} |B| \le \sum_{\mathclap{B\in\calB}} |B| = \sum_{\mathclap{B\in\calM'}} |B| - |C| + |C'|.
		\end{equation}
		Hence $|C'|\ge|C|$. Thus $|C'|=|C|$.
%	Let $C'\in\expand{\calM}{C}$ be a cycle of length $|C'|=|C|$. We will use $C'$ to show $C$ is relevant. Consider $\calM'=(\calM\setminus\{C'\})\cup\{C\}$ a cycle basis by Lemma~\ref{Lem:Bas_Ex}. $\calM'$ is a minimum cycle basis since it has the same sum of cycle lengths as $\calM$
%	\begin{equation}
%		\sum_{\mathclap{B\in\calM'}} |B| = \sum_{\mathclap{B\in\calM}} |B| - {|C'| + |C|}=\sum_{\mathclap{B\in\calM}} |B|.
%	\end{equation}
%	$C$ is a relevant cycle since it is contained in a minimum cycle basis $C\in\calM'$.
%	
%	(2, $\Leftarrow$) Let $C\in\RelCyc$, we want to show the largest cycle in $\expand{\calM}{C}$ is of length $|C|$. Using the relevance of $C$, Definition~\ref{def:rel}, consider another minimum cycle basis $\calM'$ where $C\in\calM'$. We want to use the basis exchange property with this new basis, swapping $C$ with a cycle in $\expand{\calM}{C}$ in the basis $\calM'$.
%	
%	Next, we find there is a cycle in our initial basis, specifically in $\expand{\calM}{C}$, that can be exchanged for $C$ in $\calM'$ to obtain a new basis. To see that this is possible, note $C$ is independent of the other cycles in $\calM'$ and the cycles $\expand{\calM}{C}$ sum to $C$. Thus, the cycles $\expand{\calM}{C}$ cannot be in the span of $\calM'\setminus\{C\}$.
%	 
%	Let $C'\in\expand{\calM}{C}$ be independent of $\calM'\setminus\{C\}$ such that $\calB=(\calM'\setminus \{C\})\cup\{C'\}$ is another cycle basis. Since $\calM'$ is a minimum cycle basis, we obtain the following inequality
%	Consider expanding the cycles in $\expand{\calM}{C}$ with respect to our new basis $\calM'$. The cycles $\expand{\calM}{C}$ sum to $C$ and $C$ is not in the span of $\calM'\setminus\{C\}$ by linear independence. Thus, $\expand{\calM}{C}$ cannot be in the span of $\calM'\setminus\{C\}$. Furthermore, we can fix $C'\in\expand{\calM}{C}$ where $C$ is in its expansion with respect to the new basis, i.e., $C\in\expand{\calM'}{C'}$. Finally, by Lemma~\ref{Lem:Bas_Ex}, $\calB=(\calM'\setminus\{C\})\cup\{C'\}$ is a cycle basis as well. $\calM'$ is a minimum cycle basis, so it minimizes the sum of cycle lengths, and
%	\begin{equation}
%		\begin{split}
%		\sum_{\mathclap{B\in\calM'}} |B| \le \sum_{\mathclap{B\in\calB}} |B| = \sum_{\mathclap{B\in\calM'}} |B| - |C| + |C'|\\\longrightarrow |C'|\ge |C|.
%		\end{split}
%	\end{equation}
%	Statement 1 bounds the length of cycles in $\expand{\calM}{C}$ above by $|C|$, so $C'$ is a maximal length cycle in $\expand{\calM}{C}$ and $|C'|=|C|$.
\end{proof}

\begin{corollary}[{Ref~\cite{vismaraUnion1997}, Lemma 1}]
	A cycle $C$ is not relevant if and only if it is generated by smaller cycles. In other terms, Definition~\ref{def:rel_2} is equivalent to Definition~\ref{def:rel}.
	\label{cor:rel_def_eq}
\end{corollary}
\begin{proof}
	($\Rightarrow$) Let $C$ not be relevant. Let $\calM$ be a given MCB. $C$ is the sum of shorter cycles
%	($\Rightarrow$) First, let $C$ not be relevant. Then, by Statement 2 of Lemma~\ref{Lem:Relevant_MCB_Test}, $C$ is the sum of smaller cycles. In particular,
	\begin{equation}
		C = \bigoplus_{B\in\expand{\mathcal{M}}{C}} B
	\end{equation}
%	for any MCB $\calM$, where $|B|<|C|$ for $B\in\expandd{\calM}{C}$.
	where $|B|<|C|$ for $B\in\expandd{\calM}{C}$ by Lemma~\ref{Lem:Relevant_MCB_Test}.
	
	($\Leftarrow$) Let $C$ be the sum of smaller cycles $\calX$, i.e., $C=\bigoplus_{C'\in\calX}C'$ where $|C'|<|C|$ for $C'\in\calX$. Then, the expansion of $C$ with respect to an MCB $\calM$ satisfies
	\begin{equation}
		\expand{\calM}{C} \subseteq \bigcup_{\mathclap{C'\in\calX}} \expand{\calM}{C'}.
	\end{equation}
	By construction of $\calX$ and Statement 1 of Lemma~\ref{Lem:Relevant_MCB_Test}, the cycles in $\expand{\calM}{C}$ are smaller 
	\begin{equation}
		|B|\le |C'|<|C|,\quad C'\in\calX,B\in\expand{\calM}{C'}
	\end{equation}
	as claimed.
	%Second, let $C$ be the sum of smaller cycles $\calX$, $|C'|<|C|$ for $C'\in\calX$. Then, we can expand the cycles in $\calX$ with respect to $\calM$ to obtain the double sum
%	\begin{equation}
%		C = \bigoplus_{C'\in\calX} C' = \bigoplus_{C'\in\calX}\bigoplus_{B\in\expand{\calM}{C'}}B
%		=\bigoplus_{\expand{\calM}{C}} B.
%		\label{eq:C_double_exp}
%	\end{equation}
%	{Note, $\expand{\calM}{C}\subseteq \bigcup_{C'\in\calX}\expand{\calM}{C'}$.}
%	By construction of $\calX$ and Statement 1 of Lemma~\ref{Lem:Relevant_MCB_Test}, the cycles in the right-hand side of Eq.~\eqref{eq:C_double_exp} are smaller 
%	\begin{equation}
%		|B|\le |C'|<|C|,\quad C'\in\calX,B\in\expand{\calM}{C'}.
%	\end{equation}
%	Finally, the cycles in $\expand{\calM}{C}$ are smaller than $C$, so $C$ is not relevant by Statement 2 of Lemma~\ref{Lem:Relevant_MCB_Test}.
\end{proof}

\subsection{Partitions of the relevant cycles}
\label{sub:background_partitions}

%\begin{figure}
%	\centering
%	\begin{tikzpicture}[scale=1]
%		\node[regnode] (1R) at (1.5,-1) {};
%		\node[regnode] (2Ra) at (1.1,-.33) {};
%		\node[regnode] (3Ra) at (1.1,0.33) {};
%		\node[regnode] (2Rb) at (1.9,-.33) {};
%		\node[regnode] (3Rb) at (1.9,0.33) {};
%		\node[regnode] (4R) at (1.5,1) {};
%		\draw[thick] (1R)--(2Ra)--(3Ra)--(4R)--(3Rb)--(2Rb)--(1R);
%		\node[regnode] (1T) at (-1,1.5) {};
%		\node[regnode] (2Ta) at (-.33,1.1) {};
%		\node[regnode] (3Ta) at (0.33,1.1) {};
%		\node[regnode] (2Tb) at (-.33,1.9) {};
%		\node[regnode] (3Tb) at (0.33,1.9) {};
%		\node[regnode] (4T) at (1,1.5) {};
%		\draw[thick] (1T)--(2Ta)--(3Ta)--(4T)--(3Tb)--(2Tb)--(1T);
%		\node[regnode] (1L) at (-1.5,-1) {};
%		\node[regnode] (2La) at (-1.1,-.33) {};
%		\node[regnode] (3La) at (-1.1,0.33) {};
%		\node[regnode] (2Lb) at (-1.9,-.33) {};
%		\node[regnode] (3Lb) at (-1.9,0.33) {};
%		\node[regnode] (4L) at (-1.5,1) {};
%		\draw[thick] (1L)--(2La)--(3La)--(4L)--(3Lb)--(2Lb)--(1L);
%		\node[regnode] (1B) at (-1,-1.5) {};
%		\node[regnode] (2Ba) at (-.33,-1.1) {};
%		\node[regnode] (3Ba) at (0.33,-1.1) {};
%		\node[regnode] (2Bb) at (-.33,-1.9) {};
%		\node[regnode] (3Bb) at (0.33,-1.9) {};
%		\node[regnode] (4B) at (1,-1.5) {};
%		\draw[thick] (1B)--(2Ba)--(3Ba)--(4B)--(3Bb)--(2Bb)--(1B);
%		\draw[thick] (4B)--(1R)(4R)--(4T)(1T)--(4L)(1L)--(1B);
%	\end{tikzpicture}
%	\caption{(a) Nested rings}
%	\label{fig:Nested_loops}
%\end{figure}

%Review
%\begin{enumerate}
%	\item Vismara's cycle families. %Note I will obtain them starting from an edge, not a node. Then the opposite is either an edge or a node. The total number is bounded by $mn/3+m^2/4$. The first term is the maximum number of odd cycle prototypes, where there are $m$ edges and $n$ nodes to create the prototype, and there are at least $L\ge3$ ways to choose the initial edge. The first term is tight for the complete graph. The second term states there are at most $m^2$ ordered pairs of edges possible and there are $L\ge4$ initial edges that can be chosen from an even cycle. The second term is tight for the complete bipartite graph. Check terminology but I believe a prototype is used in lieu of a cycle family.
%	\item Gleiss et al.'s interchangeability definition. Emphasize we will modify this definition.
%	\item Kolodzik et al.'s unique ring families \cite{kolodzikUnique2012}. 
%	\item Other useful lemmas. (1) \underline{Matroid property}. (2) \ul{A cycle is relevant if and only if it cannot be decomposed into smaller cycles} $\leftrightarrow$ (3) show that this is equivalent to saying \ul{``a cycle $|C|$ is relevant if and only if its decomposition with respect to a given minimum cycle basis $\mathcal{B}$ has maximum length cycle $|C|$.''} Don't prove (2) just reference Gleiss et al.
%\end{enumerate}

%****


The V-families output by Vismara's algorithm~\cite{vismaraUnion1997} are often more useful than the relevant cycles themselves. % because (i) Vismara's cycle families are polynomial in number and (ii) two cycles in a family are exchangeable for each other, as in Lemma~\ref{Lem:Bas_Ex}.
Unfortunately, these families are not unique, as they rely on a particular enumeration of the vertices. 
Here we review two works that extend V-families to unique partitions of the relevant cycles~\cite{gleissInterchangeability2000,kolodzikUnique2012}.
%Two works have since considered other partitions of the relevant cycles that are unique~\cite{gleissInterchangeability2000,kolodzikUnique2012}.

%The algorithm for computing relevant cycles introduced by Vismara... 

%The work by Gleiss {et al.}~\cite{gleissInterchangeability2000} was the first to suggest a unique partition of the relevant cycles. In Ref.~\cite{gleissInterchangeability2000}, the relation between cycles is abstracted from Vismara's cycle families, where two cycles are grouped together if they can be exchanged between bases, as in Lemma~\ref{Lem:Bas_Ex}. The formal definition is given below.
Gleiss {et al.}~\cite{gleissInterchangeability2000} extend the V-families to include cycles that can be obtained by the basis exchange property. They achieve this with the following.

%Gleiss {et al.}~\cite{gleissInterchangeability2000} extend Vismara's cycle families the following relation to partition the relevant cycles.

%{\color{red} Switch $C,C'$ to $C_1,C_2$ to match the easier notation of Koldzik et al.}
%{\color{red} for definitions I am merging notation from both refs for consistency. Defs 6\& 7 have minor aesthetic variations from their original format.}

%\begin{definition}[Ref.~\cite{gleissInterchangeability2000}, Definition 6]
%	Two relevant cycles $C_1,C_2\in\RelCyc$, $|C_1|=|C_2|$, are \emph{interchangeable} if there exists a set of cycles $\calX$, $ C_1,C_2\notin\calX$ such that
%	\begin{enumerate}
%		\item $\calX\cup\{C_2\}$ linearly independent set of relevant cycles,
%		\item $C_1 = C_2 \oplus \bigoplus_{C'\in\calX}C'$,
%		\item $|C'|\le|C_1|$ for $C'\in\calX$.
%	\end{enumerate}
%%	
%%	linearly independent subset of relevant cycles, such that $C=C'\oplus\bigoplus_{C''\in\calX}C''$ and $|C''|<|C|$ for $C''\in\calX$.
%	\label{def:intchg}
%\end{definition}
\begin{definition}[Ref.~\cite{gleissInterchangeability2000}, Definition 6]
	Two relevant cycles $C_1,C_2\in\RelCyc$ of equal length $|C_1|=|C_2|$ are \emph{interchangeable} if there exist cycles $\calX\subset\RelCyc$, $C_1,C_2\notin\calX$, such that
	\begin{enumerate}
		\item $\calX\cup\{C_2\}$ is a linearly independent set of relevant cycles,
		\item $C_1=C_2\oplus \bigoplus_{C'\in\calX} C'$,
		\item $|C'|\le |C_1|$ for $C'\in\calX$.
	\end{enumerate}
	\label{def:intchg}
\end{definition}

In words, $C_1$ and $C_2$ are interchangeable if they are of equal length, and $C_1$ is a sum of $C_2$ and cycles of smaller or equal length.

By linear independence (condition 1), there exists a basis $\calB_2$ such that $\calX\cup\{C_2\}\subseteq\calB_2$. 
%By condition 2 and Lemma~\ref{Lem:Bas_Ex}, the exchange $\calB_1=(\calB_2\setminus\{C_2\})\cup\{C_1\}$ is also a basis.
The expansion of $C_1$ is $\expand{\calB_2}{C_1} = \calX\cup\{C_2\}$ by condition 2.
By Lemma~\ref{Lem:Bas_Ex}, the exchange $\calB_1=(\calB_2\setminus\{C_2\})\cup\{C_1\}$ is a basis.
%By condition 2 $\expand{\calB_2}{C_1} = \calX\cup\{C_2\}$, so the exchange $\calB_1=(\calB_2\setminus\{C_2\})\cup\{C_1\}$ is a basis by Lemma~\ref{Lem:Bas_Ex}.


%Interchangeable cycles are related by the basis exchange property; if $C_1,C_2$ are interchangeable and $\calB_2$, $\calX\cup\{C_2\}\subseteq\calB_2$, is a basis, then $\calB_1=(\calB_2\setminus\{C_2\})\cup\{C_1\}$ is a basis by Lemma~\ref{Lem:Bas_Ex}.
Interchangeability is an equivalence relation that partitions the relevant cycles into classes. 
Cycles in a V-family are interchangeable, so the number of interchangeability classes is polynomial in number.
For the chained diamonds in Figure~\ref{fig:three_diamonds}, all large relevant cycles belong to the same interchangeability class.
Interchangeability classes are tricky to compute because they can not be obtained using a single cycle prototype from each V-family~\cite{gleissInterchangeability2000}. With additional theory, interchangeability classes can be computed in polynomial time~\cite{bergerComputing2017}.
%Interchangeability classes cannot be computed using a single cycle prototype from Vismara's cycle families, but 
%Interchangeability classes are tricky to compute because they cannot be computed using a single cycle prototype from each Vismara cycle family~\cite{gleissInterchangeability2000}. Interchangeability classes can be computed in polynomial time~\cite{


%{\color{red} why this definition is useful.}
%This definition of interchangeability is useful because it is an equivalence relation. 
%Thus, it can be used to group the relevant cycles into unique equivalence classes. 
%Thus, it naturally partitions the relevant cycles into equivalence classes.
%For the example of chained diamonds, Fig.~\ref{fig:overlapping_diamonds}, all of the large relevant cycles are interchangeable. {\color{red} hard to compute}



%Theoretically, there is a flaw with Definition~\ref{def:intchg}: two interchangeable cycles are not necessarily exchangeable between minimum cycle bases, e.g., as in Statement 2 of Lemma~\ref{Lem:Relevant_MCB_Test}. 

\begin{figure*}
	\centering
	\includegraphics{OverlapDiamond.pdf}
	\caption{ (Adapted from Ref.~\cite{gleissInterchangeability2000}, Figure 4). A graph consisting of two overlapping diamonds. The graph has five relevant cycles $\RelCyc=\{C_L,C_L',C_0,C_R,C_R'\}$, depicted separately. There are four MCBs $\calM_{1,1}=\{C_L,C_0,C_R\}$, $\calM_{1,2}=\{C_L,C_0,C_R'\}$, $\calM_{2,1}=\{C_L',C_0,C_R\}$, $\calM_{2,2}=\{C_L',C_0,C_R'\}$, where each MCB has one left cycle in $\{C_L,C_L'\}$, one right cycle in $\{C_R,C_R'\}$, and center cycle $C_0$.
		%		\textbf{(Draw example minimum cycle bases)} Adapted from~\cite{gleissInterchangeability2000}. A graph with five relevant cycles: $C_0=(1,2,4,3)$, $C_1=(1,2,4,10,9,8)$, $C_1'=(1,3,4,10,9,8)$, $C_2=(1,2,4,7,6,5)$, and $C_2'=(1,3,4,7,6,5)$. This graph has four minimum cycle bases of the form $\{C_0,C_a,C_b\}$ where $C_a\in\{C_1,C_1'\}$ and $C_b\in\{C_2,C_2'\}$. The set of independent, relevant cycles $\{C_1',C_2',C_2\}$ is not mutually relevant. Similarly, $C_1$ and $C_2$ are interchangeable ($C_1=C_2\oplus C_2'\oplus C_1'$) but not directly-interchangeable.
	}
	\label{fig:overlapping_diamonds}
\end{figure*}

%We distinguish, if $C,C'$ are interchangeable, then it is not necessarily true there exists a minimum cycle basis $\calM$ where $\calM'=(\calM\setminus\{C\})\cup\{C'\}$ is also minimal. 
$C_1,C_2$ being interchangeable does not guarantee the existence of an MCB $\calM_2$ where $\calM_1=(\calM_2\setminus\{C_2\})\cup\{C_1\}$ is also minimal. To see this, consider the graph in Figure~\ref{fig:overlapping_diamonds}.
%First, we observe $C_0$ can be added to another relevant cycle to change the path it takes through $C_0$, e.g., $C_L' = C_0\oplus C_L$. Using this observation twice, we obtain the relation
First, we observe that $C_L = C_0\oplus C_L'$ and $C_0 = C_L \oplus C_L'$. Similarly, for the right cycles%$C_R$ we exchange
\begin{equation}
	\begin{split}
	C_R = C_0\oplus C_R' &=( {C_L \oplus C_L'}) \mathrel{\oplus} C_R' \\&= C_L \oplus (C_L'\oplus C_R')
	\end{split}
	\label{eq:ind_ex}
\end{equation}
where $|C_L|=|C_L'|=|C_R'|=|C_R|$. 
%{\color{red} add another sentence, don't need as much description, except why it works.}
Thus, $C_L$ and $C_R$ are interchangeable, as in Definition~\ref{def:intchg}, where $\calX=\{C_L',C_R'\}$.
%$C_L,C_R$ are interchangeable by Eq.~\eqref{eq:ind_ex}, where $\calX=\{C_L',C_R'\}$. 
However, each MCB has exactly one left and right cycle, so we cannot exchange $C_L$ for $C_R$ between MCBs.


%{\color{red} start with example on its own
{The interchangeability relation splits cycles into classes that may be too coarse for some applications.~\cite{kolodzikUnique2012}} %The example that }
%Kolodzik {et al.} \cite{kolodzikUnique2012} argue that the notion of interchangeability from Ref.~\cite{gleissInterchangeability2000} splits the relevant cycles into too large of groups. 
As an illustrative example, consider the cube graph, where the relevant cycles are the six face cycles $F_1,F_2,\hdots,F_6$. Any face is the sum of the other face cycles, e.g.,%The sum of face cycles 
\begin{equation}
	F_6 = F_1 \oplus  F_2\oplus F_3 \oplus F_4 \oplus F_5.\label{eq:cube_face_intchg}
\end{equation}
See Figure~\ref{fig:cube}. 
%Thus, all face cycles are interchangeable and belong to a single equivalence class.
Thus, all face cycles belong to the same interchangeability class. 
Kolodzik {et al.} \cite{kolodzikUnique2012} argue that this interchangeability class is too coarse and propose the face cycles should be placed in separate groups. They achieve this using a more restrictive pairwise relation that closely resembles the relation of cycles in Vismara's families.
%Thus, all face cycles are interchangeable and form an equivalence class. Instead, they propose the following pairwise relation.
%\begin{equation}
%	F_1\oplus F_2\oplus F_3 \oplus F_4 \oplus F_5 \oplus F_6 = \emptyset.
%	\label{eq:cube_face_sum}
%\end{equation}
%Then, Eq.~\eqref{eq:cube_face_sum} can be used to show each pair of face


\begin{figure}
	\centering
	\includegraphics{cubeGraph.pdf}
	\caption{Depiction of the cube graph with labeled face cycles $F_1,F_2,\hdots,F_6$. The partial sums of $ F_1\oplus F_2\oplus\hdots \oplus F_5$ are depicted in red to show how they sum to $F_6$.}
	\label{fig:cube}
\end{figure}


%Kolodzik {et al.} argue that the notion of interchangeability is too permissive, and that it partitions the relevant cycles into
%The example they give to expand upon this critique is the simple cube, Fig~\ref{fig:cube}(a). Specifically, if we treat the cube as a graph and sum its face cycles $F_1,F_2,\hdots,F_6$ together, then we cancel to the empty cycle $\emptyset$:
%\begin{equation}
%	F_1 \oplus F_2 \oplus \hdots \oplus F_6 = \emptyset.
%	\label{eq:cube_null}
%\end{equation}
%We show this sum visually in Fig.~\ref{fig:cube}(b), and remark that we will come back to this example later in this subsection. For now, we simply observe that this sum can be rearranged to show that the faces of the cube are interchangeable. For instance, if we move $F_1$ to the opposite side of Eq.~\eqref{eq:cube_null}---by adding $F_1$ to both sides and canceling $F_1\oplus F_1=\emptyset$---then we find that $F_1$ and $F_2$ are interchangeable:
%\begin{equation}
%	F_1 = F_2\oplus(F_3\oplus F_4\oplus F_5\oplus F_6).
%\end{equation}
%However, it seems more natural to state that the cube has six distinct faces, rather than it having six equivalent faces.

% claim that all six cycles should be placed in their own family, which they 
%To remedy this, Kolodzik {et al.} propose an alternative relation between cycles
\begin{definition}[Ref.~\cite{kolodzikUnique2012}, Definition 1]
	Let $C_1,C_2\in\RelCyc$, where $|C_1|=|C_2|$, then $C_1$ and $C_2$ are unique ring family (URF)-pair-related if
	\begin{enumerate}
%		\item $|C_1|=|C_2|$
		\item $C_1$ and $C_2$ share at least one common edge%$E(C_1)\cap E(C_2)\neq \emptyset$
		\item There exists a set of strictly {smaller} cycles $\calX$ such that $C_1 = C_2\oplus \bigoplus_{C'\in \calX}C'$.
	\end{enumerate}
	\label{def:URF_Relation}
\end{definition}


The key distinction here is the cycles $\calX$ used to exchange $C_1$ and $C_2$ are strictly smaller. The faces of the cube are of the same length, so Eq.~\eqref{eq:cube_face_intchg} does not satisfy the URF-pair-relation. 
For the graph composed of chained diamonds in Figure~\ref{fig:three_diamonds}, the large cycles are URF-pair-related because they are exchanged by adding shorter diamonds. More generally, cycles in a V-family satisfy the URF-pair-relation. 


The URF-pair-relation is not transitive, so it is not an equivalence relation. Instead, the URFs are given by the transitive closure of the URF-pair-relation. 
{The URFs can be computed in polynomial time using the V-families.~\cite{kolodzikUnique2012}}
%Definition~\ref{def:URF_Relation} is not transitive because the cycles must share edges (Condition 2, $E(C_1)\cap E(C_2)\neq \emptyset$). If Condition 2 is removed from Definition~\ref{def:URF_Relation}, then we obtain an equivalence relation. We call this relationship \emph{trivial-interchangeability} and discuss it in further detail in Section~\ref{sub:TICE_def}.
%
%Both the URFs and trivial-interchangeability classes can be computed efficiently using the relevant cycle prototypes from Vismara's cycle families. We outline the algorithm for computing the trivial-interchangeability classes in Section~\ref{sub:DICE_TICE_Comp}.
%
%Looking ahead, we find these trivial and direct-interchangeability relations serve two separate purposes. The trivial-interchangeability classes adapted from Ref.~\cite{kolodzikUnique2012} are a useful proxy for cycles when minimum cycle bases are not unique. The direct-interchangeability classes adapted from Ref.~\cite{gleissInterchangeability2000} are useful for identifying 3-dimensional structures. In particular, we interpret direct-interchangeability classes as clusters of polyhedra, as discussed in Section~\ref{sub:Poly}.

%{\color{red} Don't talk about things in Sec. 3. Instead mentiong we address shortcomings in next section.}

%Unfortunately, this relation is not transitive so it does not form an equivalence relation. Thus, they define ``unique ring families'' to be the transitive closure of this relationship. For instance, consider the inner and outer ring in Fig.~\ref{fig:Nested_loops}(b).

%For the cube, this definition places each face cycle into its own unique ring family. 
%The reason for this is that the cycles $\calX$ used to exchange $C_1$ and $C_2$ must be strictly smaller, while the faces of the cube are of equal length.
%The primary reason for this is Condition 2 in Definition~\ref{def:URF_Relation}, since all face cycles are the same length, and there are no smaller cycles that can be added to swap one for another.


%****
%
%
%We now return to the notion of interchangeability introduced by Gleiss {et al.}\cite{gleissInterchangeability2000}, and argue that its output has a secondary interpretation that is more useful. Functionally, interchangeability is an equivalence relation that shows two cycles can be swapped between bases. At a glance, this suggests that two interchangeable cycles should be interpreted as roughly the same cycle, e.g., we say there is one large cycle Fig.~\ref{fig:Nested_loops}(a). As mentioned by Kolodzik {et al.}~\cite{kolodzikUnique2012}, this intuition does not work for the cube, where all the faces form a single interchangeability class. The notion of unique ring families resolves this issue, where each face of the cube.
%
%To introduce our second interpretation of interchangeability, we focus on polyhedra in more detail. Each edge in a polyhedron participates in two faces cycles. If we add each of the faces in a polyhedron together, then all of the edges cancel and we obtain the empty cycle, as in Eq.~\eqref{eq:cube_null} for the cube. Further, we showed that these zero sums can be rearranged to show that face cycles are interchangeable (assuming they are of equal length requirements as in Def.~\ref{def:intchg}). What we argue is that interchangeability can be used to find zero sums of cycles, and that these zero sums should be interpreted as polyhedra. In this context, two interchangeable cycles are faces of the same polyhedron.
%
%We note, this is a three-dimensional generalization of biconnected components. Specifically, biconnected components are clusters of loops that overlap by edges.

%To introduce our second interpretation of interchangeability, it is important to understand how this definition works with polyhedra. Each edge in a polyhedron participates in two face cycles. Thus, if we sum over the faces of a polyhedron, then all edges cancel, and we obtain the empty cycle $\emptyset$. As with the cube, this sum can then be used to show that the faces of our given polyhedron are interchangeable (assuming the face cycles are of equal length). 

%Going further, consider two connected polyhedra. For concreteness, we recommend first focusing on the stacked cubes in Fig.~\ref{fig:cube}(c). For th
%We end the background with a restatement of our thesis: the notion of interchangeability introduced by Gleiss {et al.}~\cite{gleissInterchangeability2000} is complex, but it is useful for the description of three-dimensional materials. In particular, one common network structure type that results in non-trivial interchangeability classes is polyhedra. The reason for this is that each edge in a polyhedron participates in two face cycles. If we add these two face cycles together, then this edge cancels and is not present in their sum. Furthermore, if we add all face cycles together, then all edges cancel and the total sum cancels to the empty cycle $\emptyset$.

%With this background given, we re-state the outline of the following theory section. It appears that both of these definitions may be useful in tandem.
%
%The opinion of the author on how these definitions should be interpreted are as follows. The notion of interchangeability

\section{Theory: Concepts}
\label{sec:Theory}

\begin{figure}[t]
	\centering
	\includegraphics{PartitionList.pdf}
	\caption{
		Relative granularity of each partition of the relevant cycles. The more granular partitions are refinements of the coarser ones. The axis represents the number of sets in each partition. Notably, a {\tt pi} class is a collection of {\tt sli} classes, and a {\tt sli} class is a collection of V-families. 
		An MCB is not a partition of the relevant cycles, but it is useful for comparing the granularity of each partition.
		A {\tt pi} class has a fixed number of cycles in each MCB;
		a {\tt sli} class contains at most one cycle in any given MCB. 
	}
	\label{fig:partition_granularity}
\end{figure}

In this section, we introduce two partitions of the relevant cycles and related theory. The definitions of these partitions are inspired by those in Refs.~\cite{gleissInterchangeability2000} and~\cite{kolodzikUnique2012}. We summarize the granularity of these partitions in Figure~\ref{fig:partition_granularity}.

At a coarse level, we define the polyhedron-interchangeability ({\tt pi}) relation, 
a modification of the interchangeability relation from Ref.~\cite{gleissInterchangeability2000}.
Unlike the interchangeability relation, the {\tt pi} relation is defined using the basis exchange property between MCBs.
%These definitions are distinct because direct-interchangeability is defined in terms of exchanging cycles between minimum cycle bases, whereas interchangeability is defined algebraically.
This allows us to more precisely characterize the lack of uniqueness in MCBs.
%The modification we make is that direct-interchangeability (Definition~\ref{def:dir_intchg}) is defined by exchanging cycles between minimum cycle bases whereas interchangeability (Definition~\ref{def:intchg}) is defined algebraically.
%This relation is a modification of the interchangeability relation from Ref.~\cite{gleissInterchangeability2000} to be defined in terms of exchanging cycles between minimum cycle bases. 
%See Definitions~\ref{def:intchg} and~\ref{def:dir_intchg}. 
%Direct-interchangeability is an equivalence relation so it naturally partitions the relevant cycles into equivalence classes.
We define the {\tt pi} relation and elaborate on its advantages in Section~\ref{sub:DICE_Def}. 
Also, we introduce a physical interpretation of pairs of cycles that are {\tt pi-interchangeable} as being shared faces of an abstract polyhedron. We describe this in Section~\ref{sub:Poly}. Thus, we suggest the intuition to interpret {\tt pi} classes as polyhedron clusters.
%In this section, we also outline how this definition uncovers three-dimensional structures, which we call polyhedra.

%In Section~\ref{sub:DICE_Props}, we prove theoretical properties about the direct-interchangeability relation, e.g., direct-interchangeability is an equivalence relation, so it partitions the relevant cycles into equivalence classes.

%We find that directly-interchangeable pairs of cycles can also be related as being faces of an abstract polyhedron. We describe this in Section~\ref{sub:Poly}. Thus, we suggest the intuition to interpret direct-interchangeability classes as polyhedron clusters.
%While this intuition is useful, it should be taken with a grain of salt since the direct-interchangeability classes is a property of a graph that does not necessarily have a spatial component.
%Of course this intuition should be taken with a grain of salt, as the direct-interchangeability classes are combinatorial structures, and they have no intrinsic spatial structure.
%We also characterize pairs of directly-interchangeability cycles to be faces of an abstract polyhedron, making the direct-interchangeability classes polyhedron. See Section~\ref{sub:Poly}. 
%Thus, we also interpret the direct-interchangeability classes as polyhedron clusters.

At a granular level, we define the short loop-interchangeability {(\tt sli)} relation, which modifies the URFs from Ref.~\cite{kolodzikUnique2012}. In particular, the {\tt sli} relation is equivalent to the URF-pair-relation except that cycles do not need to share edges. Our definition is more amenable to theoretical results, e.g., {\tt sli} is an equivalence relation. In computation, the {\tt sli} classes are used in place of cycles since they are polynomial in number. The {\tt sli} relation is introduced in Section~\ref{sub:TICE_def}.

\subsection{Polyhedron-interchangeability: definition and basic concepts}
\label{sub:DICE_Def}

Here, we define and investigate the {\tt pi} relation. %polyhedron-interchangeability. 
\begin{definition}
	$C_1,C_2\in\RelCyc$ are \emph{polyhedron-interchangeable} if there exists an MCB $\calM_2$, $C_2\in\calM_2$, such that $\calM_1=(\calM_2\setminus\{C_2\})\cup\{C_1\}$ is also an MCB.
	\label{def:dir_intchg}
\end{definition}
\noindent We define the operator $C_1\pInt C_2$ to denote $C_1$ is polyhedron-interchangeable for $C_2$.

An equivalent definition was introduced in Ref.~\cite{gleissInterchangeability2000} as a ``stronger interchangeability relation'', but it was not analyzed in detail. In particular, they argue this relation should be avoided because it is not symmetric. However, we find this is incorrect and that {\tt pi} is an equivalence relation. Furthermore, the {\tt pi} relation has two important benefits:
\begin{itemize}
	\item \emph{(Interpretability)} Two cycles $C_1$ and $C_2$ such that $C_1{\pInt}\, C_2$ are related by the basis exchange property between MCBs. %A pair of cycles that are {\tt pi}-interchangeable are related by the basis exchange property between MCBs. 
	This is not the case for interchangeability.
%	A pair of {\tt pi}-interchangeable cycles are related by the basis exchange property between MCBs. This is not the case for interchangeability. 
	$C_L$ and $C_R$ in Figure~\ref{fig:overlapping_diamonds} are interchangeable by Definition~\ref{def:intchg}, but they are not related by the exchange of cycles in MCBs.
	\item \emph{(Computability)} The {\tt pi} classes are computed in polynomial time using the individual cycle prototypes from the V-families. See Theorem~\ref{Thm:DirInt_Comp} and the following discussion. Interchangeability classes are also computable in polynomial time, but they require additional theory~\cite{bergerComputing2017}.
\end{itemize}

%Our notion of direct-interchangeability over the relevant cycles $\RelCyc$ is as follows
%\begin{definition}
%	$C_1,C_2\in\RelCyc$ are \emph{directly-interchangeable} if there exists a minimum cycle basis $\calM_2$, $C_2\in\calM_2$, such that $\calM_1=(\calM_2\setminus\{C_2\})\cup\{C_1\}$ is a minimum cycle basis.
%	\label{def:dir_intchg}
%\end{definition}
%\noindent We quickly show Definition~\ref{def:dir_intchg} is symmetric. If $C_1,C_2$ are directly-interchangeable, then there exists a minimum cycle basis $\calM_2$, $C_2\in\calM_2$, such that $\calM_1=(\calM_2\setminus\{C_2\})\cup\{C_1\}$ is also minimal. Clearly, the reverse exchange $\calM_2=(\calM_1\setminus\{C_1\})\cup\{C_2\}$ results in a minimum cycle basis. 
%We prove transitivity later in Theorem~\ref{Thm:DirInt_Eq_Rel}.
%In Sections~\ref{sub:DICE_Props} and~\ref{sub:Poly} we state the theoretical properties and a physical interpretation of the direct-interchangeability relation respectively.

%{\color{red} transition: useful to have an algebraic relation...}
%{\color{red} Definition~\ref{def:dir_intchg} is easy to interpret, but it can be tricky to verify in practice. }

\subsubsection{Theoretical properties}
\label{sub:DICE_Props}


\begin{figure*}
	\centering
	\includegraphics[width=6.5in]{Dice_Alg.pdf}
	\caption{
		An example showing relevant cycles and the {\tt pi} classes.
		\textbf{(a)} Graph with 6 nodes composed of 3 overlapping tetrahedra. % corresponding to a minima of the Lennard-Jones LJ6 potential. 
		\textbf{(b)} Depiction of the relevant cycles $\RelCyc=\{C_1,\hdots,C_{10}\}$. 
		\textbf{(c)} Bipartite graph representation used to compute the {\tt pi} classes.
	}
	\label{fig:DICE_alg}
\end{figure*}


Here, we state our theoretical results for the {\tt pi} relation and their implications. 
%The proofs for these results are contained in Appendix~\ref{app:DICE_Proofs}.

%First, we state a result that allows us to define the direct-interchangeability classes.

\begin{theorem}
	{\tt pi} is an equivalence relation on $\mathcal{C}_{\mathcal{R}}$.
	\label{Thm:DirInt_Eq_Rel}
\end{theorem}

The proofs of Theorem~\ref{Thm:DirInt_Comp} and the other results in Section~\ref{sub:DICE_Props} are in Appendix~\ref{app:DICE_Proofs}.

As an equivalence relation, {\tt pi} partitions the relevant cycles into equivalence classes. 
Although {\tt pi} satisfies the mathematical definition of an equivalence relation, we do not consider a pair of {\tt pi} cycles to be equivalent. For instance, the faces of the cube are {\tt pi}, but they are visually distinct. We find the {\tt pi} classes are better interpreted as polyhedron clusters as discussed in Section~\ref{sub:Poly}.

Definition~\ref{def:dir_intchg} for {\tt pi} is valuable because it explicitly uses the basis exchange property. 
However, Definition~\ref{def:dir_intchg} is hard to test in practice.
Thus, it is also useful to have an equivalent, algebraic statement of this definition, which we state below. 
%First, recall from Definition~\ref{def:ex_op} the expansion operator $\expand{\calM}{C}$ outputs the unique subset of $\calM$ with sum $\bigoplus_{C'\in\expand{\calM}{C}}C' = C$.

%Definition~\ref{def:dir_intchg} is easy to interpret, but additional theory is needed to describe when such a basis $\calM_2$ exists. First, we generalize the notion of relevance to sets of cycles.

%{\color{red} Was mutual relevance}

%We now introduce an equivalent definition of direct-interchangeability in terms of mutually relevant sets.
%We use mutual relevance to create an algebrai
%In the following lemma, we use mutually relevant cycles to create necessary and sufficient conditions for direct-interchangeability. 
%{Lemma~\ref{Lem:dir_intchg_2} is the stronger interchangeability relation mentioned in Ref.~\cite{gleissInterchangeability2000}}.
%Both definitions of direct-interchangeability, Definition~\ref{def:dir_intchg} and Lemma~\ref{Lem:dir_intchg_2}, are useful in proofs. 
%As a prerequisite, recall the expansion operator $\expand{\calM}{C}$ (Definition~\ref{def:ex_op}) outputs the unique subset of $\calM$ with sum $\bigoplus_{C'\in\expand{\calM}{C}}C' = C$.

\begin{lemma}
	Let $C_1,C_2\in\RelCyc$ where $|C_1|=|C_2|$. $C_1\pInt C_2$ %are {\tt pi}-interchangeable 
	if and only if there exist cycles $\calX\subseteq\RelCyc$, $C_1,C_2\notin \calX$, such that
	\begin{enumerate}
		\item there exists an MCB $\calM_2$ such that $\calX\cup\{C_2\}\subseteq\calM_2$ and,
		\item $C_1$ is given by the sum
		\begin{equation}
			C_1 = C_2\oplus\bigoplus_{C'\in\calX}C',\label{eq:dir_int_C1_Ex}
		\end{equation}
		i.e., the expansion of $C_1$ in $\calM_2$ is given by $\expand{\calM_2}{C_1}=\calX\cup\{C_2\}$.
	\end{enumerate}
%	{\color{red} Keep this form. Still mention it's often easier to work with a set of cycles $\calX$ rather than a minimum cycle basis. This format is useful for polyhedra...}
	\label{Lem:dir_intchg_2}
\end{lemma}
\noindent Note, the cycles in $\calX$ are bounded in length, $|C'|<|C_1|$ for $C'\in\calX$, by Statement 1 of Lemma~\ref{Lem:Relevant_MCB_Test}. Moreover, $\calX$ satisfies the conditions for $C_1,C_2$ to be interchangeable as in Ref.~\cite{gleissInterchangeability2000} (Definition~\ref{def:intchg}) with the added constraint that $\calX\cup\{C_2\}$ belongs to a MCB. 
Thus, {\tt pi} implies interchangeability. The reverse is not true, e.g., $C_L$ and $C_R$ in Figure~\ref{fig:overlapping_diamonds} are interchangeable but not {\tt pi}-interchangeable.

In brief, Lemma~\ref{Lem:dir_intchg_2} states a cycle $C_1$ is {\tt pi}-related to the equal length cycles in its expansion with respect to a MCB $\calM_2$. One may wonder, do we need many MCBs to construct the {\tt pi} classes? Fortunately, it is sufficient to generate a single MCB, as stated below.

%Lemma~\ref{Lem:dir_intchg_2} clarifies that $C_1$ is directly-interchangeable for $C_2$ via $\calM_2$ if $C_2$ is an equal length cycle in the expansion of $C_1$ in $\calM_2$. This makes Statement 2 of Lemma~\ref{Lem:Relevant_MCB_Test} bidirectional. This second definition of direct-interchangeability is easier to verify computational, and it is useful for proofs.

%Before proving Lemma~\ref{Lem:dir_intchg_2}, we observe that this result can be interpreted in two ways. First, Lemma~\ref{Lem:dir_intchg_2} elaborates that $C_1,C_2$ are directly-interchangeable via the minimum cycle basis $\calM_2$ if and only if $C_2\in\expand{\calM_2}{C_1}$. This makes Statement 2 of Lemma~\ref{Lem:Relevant_MCB_Test} bidirectional. Second, direct-interchangeability can be verified using mutually relevant sets of cycles. In later proofs, it is easier to manipulate mutually relevant sets rather than explicitly defining minimum cycle bases. 


%Definition~\ref{def:intchg} with $\calX\cup\{C_2\}$ a subset of a minimum cycle basis is the stronger interchangeability relation mentioned in Ref.~\cite{gleissInterchangeability2000}.

%\begin{proof}
%	($\Leftarrow$) Let $\calX$ be a set of cycles such that $\calX\cup\{C_2\}$ is mutually relevant and $C_1$ is given by the Eq.~\eqref{eq:dir_int_C1_Ex}. By Definition~~\ref{def:mut_rel} of mutual relevance, there exists a minimum cycle basis $\calM_2$ such that $\calX\cup\{C_2\}\subseteq\calM_2$. The expansion of $C_1$ in $\calM_2$ is $\expand{\calM_2}{C_1}=\calX\cup\{C_2\}$ by Eq.~\eqref{eq:dir_int_C1_Ex}. Thus, we may exchange $C_2$ for $C_1$ to obtain a new basis $\calM_1=(\calM_2\setminus\{C_2\})\cup\{C_1\}$, which is minimal by $|C_1|=|C_2|$. See Statement 2 of Lemma~\ref{Lem:Relevant_MCB_Test}. In summary, we have shown that $C_1\DirEx C_2$ via $\calM_2$ by Definition~\ref{def:dir_intchg}.
%	
%	($\Rightarrow$) Let $C_1,C_2$ be directly-interchangeable. By Definition~\ref{def:dir_intchg}, let $\calM_2$ be a minimum cycle basis such that $\calM_1=(\calM_2\setminus\{C_2\})\cup\{C_1\}$ is also a minimum cycle basis. 
%	We observe, the cycles $\{C_1\}\cup\expand{\calM_2}{C_1}$ are linearly dependent because they sum to the empty set $\emptyset$ by 
%	\begin{equation*}
%		C_1=\enspace\bigoplus_{\mathclap{C'\in\expand{\calM_2}{C_1}}}\enspace C'.
%	\end{equation*}
%	Of course, the basis cycles $\calM_1$ are linearly independent. Thus, the cycles $\{C_1\}\cup\expand{\calM_2}{C_1}$ are not a subset of $\calM_1$, i.e., $\expand{\calM_2}{C_1}\not\subseteq\calM_2\setminus\{C_2\}$ by definition of $\calM_1$. Therefore, $C_2$ is contained in the expansion $C_2\in\expand{\calM_2}{C_1}$, since $\expand{\calM_2}{C_1}\subseteq\calM_2$. This is our desired result, as the mutually relevant cycles $\calX\cup\{C_2\}\coloneq\expand{\calM_2}{C_1}$ satisfy the given hypothesis.
%\end{proof}







%In this subsection, we prove two key properties about direct-interchangeability: 
%\begin{itemize}
%	\item Direct-interchangeability is an equivalence relation. See Theorem~\ref{Thm:DirInt_Eq_Rel}.
%	\item The resulting equivalence classes can be computed using the cycle prototypes from Vismara's families and any minimum cycle basis $\calM$. See Theorem~\ref{Thm:DirInt_Comp} and discussion after it.
%\end{itemize}
%
%%Before proving these theorems, we find it helpful to introduce a lemma that helps with bookkeeping. 
%Before proving these theorems, we find it useful to introduce a few preliminary results that help with bookkeeping mutually relevant sets -- see Lemma~\ref{Lem:SmallCycs2MCB} and Corollaries~\ref{Cor:dir_intchg_3} and~\ref{Cor:mut_rel_2}. We start with a simple generalization of Statement 1 of Lemma~\ref{Lem:Relevant_MCB_Test}.


%{\color{red} Was Lem: set of small cycles to a set of small cycles}


%Recall, to verify that $C_1$ and $C_2$ are directly-interchangeable as in Lemma~\ref{Lem:dir_intchg_2} we must generate a set of cycles $\calX$ such that $\{C_2\}\cup\calX$ is a subset of a minimum cycle basis and $C_1=C_2\oplus\bigoplus_{C\in\calX}C$.
%The utility of Lemma~\ref{Lem:SmallCycs2MCB} is that it allows us to focus our attention on the cycles in $\calX$ of equal length $|C|=|C_1|$. 
%The utility of Lemma~\ref{Lem:SmallCycs2MCB} is that it allows us to focus our attention on cycles of a single fixed length. 
%In the following corollary, we use Lemma~\ref{Lem:SmallCycs2MCB} to loosen the constraints on the smaller cycles in $\calX$ (those of length $|C|<|C_1|$), where they no longer need to be mutually relevant, relevant, or linearly independent.
%In the following corollary of Lemma~\ref{Lem:SmallCycs2MCB}, we remove the constraints on the shorter cycles in $\calX$ (those of length $|C|<|C_1|$), where they need not be mutually relevant, relevant, or even linearly independent.


%{\color{red} Was the smaller cycles }

%In practice, we do not use the smaller cycles $\calY$ in Corollary~\ref{Cor:dir_intchg_3}, so we will define them implicitly.
%We will state that $C_1$ is directly-interchangeable for $C_2$ via $\calX$ up to smaller cycles if $\{C_2\}\cup\calX$ is mutually relevant and
%\begin{equation}
%	C_1 = C_2\oplus\bigoplus_{C'\in\calX}C'\oplus\smCyc
%\end{equation}
%where $\smCyc$ is the sum of a set of smaller cycles $\calY$.

%Next, we use Lemma~\ref{Lem:SmallCycs2MCB} to provide conditions for when a set of cycles of equal length is mutually relevant. 
%Interestingly, if we replace $\smCyc$ with the empty set $\emptyset$ in Eq.~\eqref{eq:mut_rel_ind} in Corollary~\ref{Cor:mut_rel_2}, then we obtain a condition for linear dependence. Thus, Corollary~\ref{Cor:mut_rel_2} defines mutual relevance as an abstract notion of linear independence over the relevant cycles.
%We observe, a set of cycles $\calX$ is linearly dependent if there exists a subset $C_1,\hdots,C_k$ such that
%\begin{equation}
%	C_1\oplus\hdots\oplus C_k=\emptyset.
%\end{equation}
%If we replace the empty set with the sum of smaller cycles then we obtain Equation~\eqref{eq:mut_rel_ind} in Corollary~\ref{Cor:mut_rel_2}, making $\calX$ not mutually relevant. Thus, Corollary~\ref{Cor:mut_rel_2} defines mutual relevance as an abstract notion of linear independence
%We note, Eq.~\eqref{eq:mut_rel_ind} in Corollary~\ref{Cor:mut_rel_2} resembles a condition for linear dependence, i.e.,  cycles $\calX$ are linearly dependent if there exists a collection of cycles $C_1,\hdots,C_k\in\calX$ such that
%\begin{equation}
%	C_1\oplus\hdots\oplus C_k=\emptyset.
%\end{equation}
%Thus, Corollary~\ref{Cor:mut_rel_2} can be viewed as an abstracted notion

%Remark, no conditions are placed on the cycles $\calY$, i.e., independence or relevance. We find it easier to define the cycles $\calY$ implicitly to avoid extra bookkeeping. Specifically, we state $C_1$ is directly-interchangeable for $C_2$ via cycles $\calX$ satisfying the given hypothesis if
%\begin{equation}
%	C_1 = C_2 \oplus \bigoplus_{C'\in\calX}C'\oplus\smCyc
%\end{equation}
%where (smaller cycles) is the sum of an arbitrary set of smaller cycles $\calY$.







\begin{theorem}
	Let $\calM$ be a given MCB. Let $C_2\in\calM$ and $C_1\in\RelCyc{\setminus}\calM$. If $C_2\in\expandd{\calM}{C_1}$ and $|C_1|=|C_2|$ then $C_1\pInt C_2$.
	The {\tt pi} classes are given by the transitive closure of these relations.\nolinebreak % gives an extra line break...
	\label{Thm:DirInt_Comp}
\end{theorem}


%\newpage

Theorem~\ref{Thm:DirInt_Comp} reduces the task of computing the {\tt pi} classes to breadth-first search. As input, we compute an MCB $\calM$ and the relevant cycles $\RelCyc$. Then, we expand the relevant cycles in $\calM$. Using this, we create a bipartite graph, where nodes represent the relevant cycles split by if they are in $\calM$. Two cycles $C_1\in\calM$ and $C_2\in\RelCyc\setminus\calM$ are joined by an edge if $C_2\in\expand{\calM}{C_1}$. The {\tt pi} classes are the connected components of this bipartite graph. See Figure~\ref{fig:DICE_alg}~(c) for an instance of this bipartite graph.

%{\color{red} ***Lacks context now*** To reduce the computational complexity of this procedure, we can restrict ourselves to a single prototype cycle from each of Vismara's families. This approach is valid because cycles in Vismara's families differ by shorter cycles. Then, the nodes in the bipartite representation are the prototypes for Vismara's families, and a direct-interchangeability class is the union of Vismara's families in a connected component. We discuss this further in Section~\ref{sub:DICE_TICE_Comp}.}

%\begin{algorithm}
%	\caption{text}
%	
%	\KwIn{graph $G$}
%	\KwOut{direct and trivial interchangeability classes}
%	\color{magenta}
%	Compute Vismara's prototypes decomposed w.r.t. FCB (matrix $\bar{R}\in\{0,1\}^{\nu\times}$)\;
%	Compute a minimum cycle basis\;
%	Compute inverse matrix for MCB $A^{-1}$\;
%	Sort MCB and  $A^{-1}$ by decreasing length\;
%	$R\gets A^{-1}R$\Comment*[r]{Expand Vismara's prototypes in $\mathcal{B}$}
%	Sort $R$ in lexographic order\;
%	Remove rows of $R$ where the largest cycle is larger than the prototype\;
%	Extract trivial-interchangeability classes\;
%	Extract direct-interchangeability classes via breadth-first search\;
%\end{algorithm}


%We end with a pair of convenient properties for the direct-interchangeability relation. 
The {\tt pi} relation can be used to construct a path from one MCB to another by only swapping one pair of cycles at a time. 
%First, we can always swap a cycle from one minimum cycle basis for a cycle from another:
\begin{lemma}
	Let $\calM_1,\calM_2$ be distinct MCBs. There exists $C_1\in\calM_1\setminus\calM_2$ and $C_2\in\calM_2\setminus\calM_1$ such that $\calM_1'\,{=}\,(\calM_1\setminus\{C_1\})\cup\{C_2\}$ is an MCB.%\nolinebreak
	\label{Lem:DICE_Irreducible}
\end{lemma}
%\noindent Moreover, we find that any minimum cycle basis is reachable from any other minimum cycle basis. We can repeatedly swap cycles from $\calM_1$ for cycles in $\calM_2$ to relate these bases by the basis exchange property. 
In this procedure, each cycle in $\calM_1$ also belongs to $\calM_2$ or it is {\tt pi}-interchangeable for a cycle in $\calM_2\setminus\calM_1$. This results in a bijective mapping from $\calM_1$ and $\calM_2$ between {\tt pi}-interchangeable cycles. From this, we obtain the following result.
%Second, There is a fixed number of cycles from a direct-interchangeability class that belong to a given minimum cycle basis:
\begin{corollary}
	Let $\calM_1,\calM_2$ be MCBs and $\mathcal{PI}\subseteq\RelCyc$ be a {\tt pi} class. Then, $\calM_1$ and $\calM_2$ have the same number of cycles in $\mathcal{PI}$:
	\begin{equation}
		|\calM_1\cap\mathcal{PI}| = |\calM_2\cap\mathcal{PI}|.
	\end{equation}
\end{corollary}
\noindent This introduces a notion of rank for the {\tt pi} classes, given by the number of cycles in the {\tt pi} class that belong to a MCB $|\calM\cap\mathcal{PI}|$.



\subsubsection{{\tt pi} classes as polyhedron clusters}
\label{sub:Poly}



%Here, we describe how direct-interchangeability may be used to identify 3-dimensional structures. In particular, we describe how clusters of polyhedra belong to a single direct-interchangeability class. {\color{red} ...}

%Our notions of interchangeability and those in Refs.~\cite{gleissInterchangeability2000,kolodzikUnique2012} 
The partitions of the relevant cycles 
were introduced as a means of addressing the lack of uniqueness in MCBs. In the examples considered so far, lack of uniqueness can be caused by symmetries or by three-dimensional data. Here, we will discuss the latter. Specifically, we will discuss how {\tt pi} classes may be used to identify three-dimensional structures.

First, consider a graph composed of a convex polyhedron with face cycles $F_1,\hdots,F_k$ ordered by decreasing length $|F_1|\ge|F_2|\ge\hdots\ge|F_k|$.
For instance, our graph could be one of the platonic solids, e.g., the tetrahedron or cube. 
We observe, that each edge in the polyhedron participates in two face cycles. Thus, all edges cancel in the sum $F_1\oplus\hdots\oplus F_k$. In particular, the face cycles sum to the empty set
\begin{equation}
	F_1\oplus \hdots\oplus F_k=\emptyset.
	\label{eq:poly_null_sum}
\end{equation}
Next, we re-arrange Eq.~\eqref{eq:poly_null_sum} to express $F_1$ as the sum of the remaining face cycles
\begin{equation}
	F_1 = F_2\oplus (F_3\oplus\hdots \oplus F_{k}).\label{eq:dir_int_poly}
\end{equation}
See Figure~\ref{fig:cube} for the case of the cube. 
Eq.~\eqref{eq:dir_int_poly} resembles Eq.~\eqref{eq:dir_int_C1_Ex} in the definition of the {\tt pi} relation given by Lemma~\ref{Lem:dir_intchg_2}.
%Equation~\eqref{eq:dir_int_poly} takes a form that resembles our previous notions of interchangeability. 
Indeed, if we add two constraints (i) $|F_1|=|F_2|$ and (ii) $F_2,\hdots,F_{k}$ form an MCB, then $F_1$ and $F_2$ are {\tt pi}-interchangeable by Lemma~\ref{Lem:dir_intchg_2}.

\begin{figure}[t]
	\centering
	\includegraphics{PolyEx.pdf}
	\caption{\textbf{(a)} The square pyramid, a polyhedron whose faces do not form a {\tt pi} class. %does not result in a {\tt pi} relation. 
%		The triangular faces of this graph form a unique MCB.
		The unique MCB of this graph is the set of triangular faces. 
		\textbf{(b)} Carbon skeleton of adamantane, $\rm C_{10}H_{16}$, the smallest diamondoid, or a small molecule with the same structure as diamond. The faces of adamantane are a {\tt pi} class, %are {\tt pi}-interchangeable, 
		but they are not flat as in a classical polyhedron.}
	\label{fig:poly_counter}
\end{figure}


We remark, not all polyhedra can be used to obtain a pair of face cycles $F_1,F_2$ such that $F_1\pInt F_2$. % $F_1,F_2$ of {\tt pi}-interchangeable face cycles. 
For example, if $F_1$ is strictly longer than the other faces $|F_1|>|F_i|$, $i=2,\hdots,k$, then $F_1$ is not relevant as the sum of smaller cycles by Eq.~\eqref{eq:dir_int_poly}. This is the case for the square pyramid, where the square base is the largest cycle in Figure~\ref{fig:poly_counter} (a). Nevertheless, polyhedra are useful for constructing examples for the {\tt pi} relation.

The {\tt pi} classes can be used to identify three-dimensional structures. If $C_1,C_2$ are {\tt pi}-interchangeable, then we can use Lemma~\ref{Lem:dir_intchg_2} to obtain a set of cycles $\calP=\{C_1,C_2\}\cup\calX$ with null sum
\begin{equation}
	C_1\oplus C_2 \oplus\bigoplus_{C'\in\calX}C' = \emptyset.
\end{equation}
This matches the sum of face cycles for polyhedra, Eq.~\eqref{eq:poly_null_sum}. Thus, we propose to interpret these sets $\calP$ with null sums obtained by Lemma~\ref{Lem:dir_intchg_2} as polyhedra.

It is not necessarily the case that we can draw these null sums as polyhedra, e.g., adamantane in Figure~\ref{fig:poly_counter} (b) does not have flat faces. Instead, the sets $\calP$ obtained via {\tt pi} relations might be more appropriately defined as abstract polyhedra~\cite{mcmullenAbstract2002} with possibly curved faces and edges. 
%We will not discuss this distinction further and, for convenience, we will still call the 3-dimensional structures obtained via direct-interchangeability relations polyhedra.

\begin{figure}[t]
	\centering
	\includegraphics{TwoCubes.pdf}
	\caption{Graph composed of two cubes $\calP_1$ and $\calP_2$. Both $\calP_1$ and $\calP_2$ are sets of six face cycles. The symmetric difference $\calP_1\symdiff\calP_2$ (depicted right) gives the ten faces of the larger rectangular prism containing $\calP_1$ and $\calP_2$. The interior face cycle (shaded with gray stripes) belonging to both $\calP_1$ and $\calP_2$ is not present in $\calP_1\symdiff\calP_2$.}
	\label{fig:polyhedron_pair}
\end{figure}

We remark, a {\tt pi} class is often not a single polyhedron. For instance, if we glue two cubes together so that they share a face as shown in Figure~\ref{fig:polyhedron_pair}, then all square faces are {\tt pi}-interchangeable by transitivity. Thus, in the same way that a biconnected component is a cluster of loops, we argue that a {\tt pi} class should be interpreted as a cluster of polyhedra.

Furthermore, we may extend the polyhedra obtained by a {\tt pi} class to a vector space. As before, the addition operator of this vector space is the symmetric difference, and the scalars are the binary GF(2) field. This ``polyhedron space'' is simply the null space of the relevant cycles, i.e., if polyhedra $\calP_1$ and $\calP_2$ have null sums then so does their symmetric difference
\begin{equation}
	\bigoplus_{C'\in\calP_1\symdiff\calP_2} C' = \bigoplus_{C'\in\calP_1}C'\oplus\bigoplus_{C''\in\calP_2}C''.
\end{equation}
A polyhedron space is likely useful for graphs representing material structures, where there are three intrinsic spatial dimensions. It appears pertinent to define a notion of ``minimum polyhedron basis'' to serve as a summary statistic for this space. We lack the algorithms and data to justify such a construction. 
%Instead, we will provide a conjecture on how such a basis might be defined in our discussion, Section~\ref{sec:discussion}.

%\add{Why I am not fully defining minimum polyhedron basis here? It's nuanced, and it helps to have trivial-interchangeability to give a definition that has better characteristics. The biggest issue is if there is an exponential number of relevant cycles, then the number of basis polyhedra is the number of dependent relevant cycles $|\RelCyc|- \nu$ where $\nu$ is the rank of the cycle space. Should I clarify this here?}


%\textbf{Paragraph outline:}
%\begin{itemize}
%	\item Lemma 6 can be used to obtain a set of cycles $\calP = \{C_1,C_2\}\cup\calX$ that sum to 0. Resembling Eq.~\eqref{eq:poly_null_sum} for the polyhedra.
%	\item We interpret these null sums as polyhedra.
%	\item It is not necessarily the case that we can draw these null sums in 3-dimensional space as polyhedra, e.g., barallene Figure~\ref{fig:MCB_Unique} which doesn't have flat faces. More correctly, defined as an abstract polytope [add a cite]. don't discuss just mention the technicality.
%	\item In practice, we might be able to use this to obtain unit cells (adamantane example?) for metals or for crystals using only combinatorial data.
%	\item \textbf{Where to place:} More useful generalization of cycle space to 3d data. probably at end, but hint earlier?
%\end{itemize}

%\textbf{New paragraph:}
%\begin{itemize}
%	\item These polyhedra as null sums form a vector space. Addition operation symmetric difference $\calP_1\symdiff\calP_2$ like the cycle space.
%	\item In some cases it may be useful to construct a minimum polyhedron basis. This likely only works well when there is not a large variance of cycle lengths. 
%\end{itemize}



\subsection{Short loop-interchangeability: definition, equivalence classes, \& computation}

\label{sub:TICE_def}

%\add{Items I want this section to address:
%	\begin{itemize}
%		\item Basic: what is trivial-interchangeability? What is the theory?
%		\item More theory is annoying. Where is direct-interchangeability insufficient? 
%		\item Our notion of trivial-interchangeability modifies the unique ring families of Kolodzik {et al.}. What benefits does this definition have? (Mostly theoretical, I do not aim to argue that URFs are not practically useful)
%		\item For what materials might this be useful?
%		\begin{itemize}
%			\item Perhaps in polymer-like materials with diverse loop lengths are common. Then this definition collapses
%		\end{itemize}
%		\item What do trivial-interchangeability classes typically look like?
%	\end{itemize}}


%{\color{red} reminder of direct-interchangeability. Start with, remind generalized polyhedra.}

%{\color{red} get to definition}

%{\color{red} need to say trivial-interchangeability implies direct-interchangeability}

%{\color{red} (dissertation? Check if phd) mention introduced in Gleiss. We extend of Gleiss. Lemma 12 and Corollary 14 yes!}

%\textbf{
%	1. Motivation (DICE Context, nanotube), 
%	2. Definition, 
%	3. Citations, 
%	4. Comparisons 
%	5. Theory
%	}

% 1. DICE context, where polyhedra are ineffective

%Here, we introduce a second equivalence relation to partition the relevant cycles called short loop-interchangeability ({\tt sli}). 
Here, we introduce the second {\tt sli} relation used to partition the relevant cycles.
To motivate this, we discuss where the {\tt pi} classes are less effective as a descriptive tool.

In Section~\ref{sub:Poly}, the {\tt pi} relation was related to three-dimensional structures, i.e., polyhedra with non-flat faces.
%In Section~\ref{sub:Poly}, we related the direct-interchangeability %related lack of uniqueness 
%to three-dimensional structures (polyhedra). 
%Recall, the direct-interchangeability
%For instance, 
Consider the graph shown in Figure~\ref{fig:2ManyPoly}~(a). 
The polyhedra are sampled using an MCB $\calM$ by merging each cycle $C\in\RelCyc\setminus\calM$ with its expansion $\expand{\calM}{C}$ as in Figure~\ref{fig:2ManyPoly}~(b).
%We sample polyhedra by constructing a minimum cycle basis $\calM$ and merging each cycle $C\in\RelCyc\setminus\calM$ with its expansion $\expand{\calM}{C}$ as in Figure~\ref{fig:2ManyPoly}~(b).
These polyhedra describe three-dimensional structures like beads of a bracelet, as in Figure~\ref{fig:2ManyPoly}~(c).

Visually, the third polyhedron $\calP_3$ appears redundant, where its two beads are a combination of the other polyhedra. 
This is not resolved by treating the polyhedra as a vector space since $\calP_3\neq\calP_1\symdiff\calP_2$. 
%This redundancy is not solved by the symmetric difference, since $\calP_3\neq \calP_1\symdiff \calP_2$.
%However, the third polyhedra in the bottom of Figure~\ref{fig:2ManyPoly}~(c) is redundant, where its two beads can be seen as a combination of the other two polyhedra.
As the bracelet grows in length, the number of visually redundant polyhedra grows exponentially with the number of relevant cycles.
%this issue amplifies as the number of relevant cycles and polyhedra grows exponentially.
%Numerous cycles around the circumference of a nanotube (see Figure~\ref{fig:tubes}~(a)) form a large number of redundant polyhedra.
%However, these polyhedra are redundant, where the third polyhedra with both beads in the bottom of Figure~\ref{fig:2ManyPoly}~(c) is a combination of the other two polyhedra.

\begin{figure}[t]
	\centering
	\includegraphics{2ManyPoly2.pdf}
	\caption{
		\textbf{(a)} Bracelet graph composed of two diamonds chained together. 
		The planar faces form an MCB $\calM$. 
		%The minimum cycle basis $\calM$ composed of face cycles are labeled. 
		\textbf{(b)}~Three polyhedra $\calP_1\,{=}\,\{C_1{,}C_0{,}D_1\}, \calP_2\,{=}\,\{C_2{,}C_0{,}D_2\},$ and $\calP_3=\{C_3{,}C_0{,}D_1{,}D_2\}$ constructed using the cycles not in $\calM$ and their expansions in $\calM$.
		\textbf{(c)}~Three-dimensional representation of the polyhedra as the union of their face cycles.
	}
	\label{fig:2ManyPoly}
\end{figure}

To address this, 
we observe the cycles in the examples above vary in length. 
%we observe in the examples above that the cycles have varying lengths. 
We argue that larger cycles that differ by shorter ones should be treated as equivalent. 
For instance, we would group the circumference cycles around the carbon nanotube into an equivalence class without sampling polyhedra. 
%we attribute many hexagonal faces and a single circumference cycle to the carbon nanotube. 
%In this setting, lack of uniqueness is often caused by symmetries in the graph, rather than by three-dimensional structures. 
%We make this distinction using the following definition.
We formalize this equivalence with the {\tt sli} relation.

% 2.

\begin{definition}
	$C_1,C_2\in\mathcal{C}_{\mathcal{R}}$, $|C_1|=|C_2|$, are \emph{short loop-interchangeable} if
	\begin{equation}
		C_1=C_2\oplus\smCyc,
		\label{eq:triv_intchg}
	\end{equation}
	i.e., there exist cycles $\calY$, $|C| < |C_1|$ for $C\in\calY$, such that
	\begin{equation}
		C_1 = C_2\oplus\bigoplus_{C\in\calY}C.
	\end{equation}
	\label{def:triv_intchg}
\end{definition}
\noindent We use the operator $C_1\slInt C_2$ to denote $C_1$ is short loop-interchangeable for $C_2$.

The numerous cycles around the circumference of the carbon nanotube in Figure~\ref{fig:tubes}~(a) form a single {\tt sli} class.


%{\tt sli}-interchangeability is a stronger condition than {\tt pi}-interchangeability.
%The {\tt sli} relation is stronger than the {\tt pi} relation.
{\tt sli} is a stronger relation than {\tt pi}.
A pair of {\tt pi}-related cycles may differ by equal length cycles. Therefore, the faces of a polyhedron with uniform lengths (e.g., the cube) form a {\tt pi} class but not a {\tt sli} class. %are {\tt pi}-interchangeable, but not {\tt sli}-interchangeable.
We sample polyhedra from {\tt pi} classes that do not reduce to a {\tt sli} class. This is discussed in more detail in Section~\ref{sub:DICE_TICE_Comp}.

\begin{figure}[t]
	\centering
	\includegraphics{Tubes.pdf}
	\caption{
		\textbf{(a)} Small carbon nanotube. The large cycles around the circumference of the nanotube are both a {\tt sli} class and a URF. %{\tt sli}-interchangeable, and they are connected by the transitive closure of the URF-pair-relation.
		\textbf{(b)}~Hexagonal prism graph. The two hexagonal faces form a {\tt sli} class but not a URF. %are {\tt sli}-interchangeable, but they do not satisfy the URF-pair-relation.
	}
	\label{fig:tubes}
\end{figure}

%This definition of trivial-interchangeability 
Definition~\ref{def:triv_intchg} was first introduced in the dissertation by Gleiss~\cite{gleissShort2001} as a stronger interchangeability relation. Two of the results shown here, Lemma~\ref{Lem:TrInt_Eq_Rel} and Corollary~\ref{Cor:tr_intchg_rank}, were first shown there. 
We include them for completeness. 
Here, we add a method of computing the {\tt sli} classes, Lemma~\ref{Lem:TrEx_Alg}. 
Also, our two-fold partition using both {\tt pi} and {\tt sli} classes is novel.
The proofs of these results are in Appendix~\ref{app:TICE_Proofs}
%Also, our two-fold partition using both trivial-interchangeability and direct-interchangeability is novel.

Our approach of equivocating cycles that differ by shorter cycles is inspired by the URFs in Ref.~\cite{kolodzikUnique2012}. 
The {\tt sli} relation is weaker than the URF-pair-relation.
%{\tt sli} is a weaker relation than the URF-pair-relation (Definition~\ref{def:URF_Relation}). 
Specifically, cycles do not need to share edges to be {\tt sli}-interchangeable. 
For example, the two hexagonal cycles around the perimeter of the prism shown in Figure~\ref{fig:tubes}~(b) are {\tt sli}-interchangeable but not URF-pair-related. 
In MD data, we find these two definitions behave similarly. 
We use the {\tt sli} relation because it is simpler, making it easier to develop theory.
%, e.g.,
%trivial-interchangeability benefits from being an equivalence relation.

\begin{lemma}
	{\tt sli} is an equivalence relation on $\mathcal{C}_{\mathcal{R}}$.
	\label{Lem:TrInt_Eq_Rel}
\end{lemma}

By Lemma~\ref{Lem:TrInt_Eq_Rel}, the {\tt sli} relation partitions the relevant cycles into equivalence classes. 
To compute these classes, we redefine {\tt sli} using an MCB.
%To compute these equivalence classes, we redefine trivial-interchangeability using a minimum cycle basis.

\begin{lemma}
	\label{Lem:TrEx_Alg}
	Let $\calM$ be a given MCB. $C_1,C_2\in\RelCyc$, $|C_1|=|C_2|$, are {\tt sli}-interchangeable if and only if their expansions in $\calM$ are equal up to shorter cycles.
\end{lemma}

\noindent Cycles in Vismara's families differ by shorter cycles. Therefore, the relevant cycles can be broken into {\tt sli} classes in polynomial time using Vismara's cycle families.
%To compute the trivial-interchangeability classes, we redefine the trivial-interchangeability relation
%Because of this, the trivial-interchangeability classes naturally partition the relevant cycles. Typically, tube-like graphs (e.g., those in Figure~\ref{fig:tubes}) form a trivial-interchangeability class because the large cycles at each of end of the tube are separated by a series of short cycles. Not all trivial-interchangeability classes need to look like cylinders.

Note, if a cycle $C_1$ belongs to the MCB $\calM_1$, then its expansion in $\calM_1$ is $\expand{\calM_1}{C_1}=\{C_1\}$. 
If $C_2\slInt C_1$ then the expansion of $C_2$ in $\calM_1$ is $C_1$ and shorter cycles by Lemma~\ref{Lem:TrEx_Alg}. 
Thus, we may exchange $C_2$ for $C_1$ in $\calM_1$ to obtain a new MCB by Lemma~\ref{Lem:Relevant_MCB_Test}. 
%This is where we get the name ``trivial''-interchangeability, where we can freely exchange $C_2$ for $C_1$
We state this as a corollary.

\begin{corollary}
	Let $C_1,C_2$ be cycles such that $C_1{\slInt}\, C_2$. %be {\tt sli}-interchangeable 
	Let $\calM_1$ be an MCB where $C_1{\in}\,\calM_1$. Then, $\calM_2=(\calM_1\setminus\{C_1\})\cup\{C_2\}$ is an MCB.
	\label{Cor:TrInt_Exchanges}
\end{corollary}

Also, observe $C_2$ is not in the MCB $\calM_1$ since it is already in the span of $C_1$ and shorter cycles in $\calM_1$. 
Therefore, $C_1$ is the only cycle from its {\tt sli} class that belongs to $\calM_1$:

\begin{corollary}
	Let $\calM$ be an MCB and $\calS\subseteq\RelCyc$ a {\tt sli} class. At most one cycle in $\calS$ belongs to $\calM$, i.e., $|\calM\cap\calS|\le 1$. %the trivial-interchangeability class belongs to the minimum cycle basis, i.e., $|\calM\cap\calT|\le 1$.
	\label{Cor:tr_intchg_rank}
\end{corollary}

Note, the intersection may be empty, i.e., $|\calM\cap\calS|=0$. For instance, consider the cube where each face is its own {\tt sli} class, and only five faces belong to each MCB.








\section{Algorithms for sampling {\tt pi} and {\tt sli} classes}
\label{sec:algorithm}

In this section, we review and upgrade algorithms for computing the relevant cycle families, {\tt pi} classes, and {\tt sli} classes. In Section~\ref{sub:Vis_Fam}, we modify the relevant cycle families from Ref.~\cite{vismaraUnion1997} to be easier to compute in our setting. In Section~\ref{sub:MCB_comp}, we {show how a MCB can be extracted from these families using the algorithms from Ref.~\cite{kavithaFaster2004}.} In Section~\ref{sub:MCB_COB}, we show how to compute the expansion of cycles with respect to a MCB. In Section~\ref{sub:DICE_TICE_Comp}, we compute the {\tt pi} and {\tt sli} classes using the MCB representation of the cycle families.
%One motivation for the definitions of directly and trivially interchangeable cycles is that they are computationally easy to work with...
%
%{\color{red} Add a note on sparsity
%
%Similarly, we can consider a sparse representation of the $(b_1,\hdots,b_\nu)$-vector, where we instead store the list of indices $i$ such that $b_i=1$. For ReaxFF simulation data, these vectors are very sparse, i.e., one or two nonzero entries, so this is the representation we use in our codes.
%}

\subsection{Modified Vismara's cycle families}
\label{sub:Vis_Fam}

%{\color{red}$\bullet$ V family for Vismara cycle family}

\subsubsection{Definition}

{Here, we describe a method for computing candidate families of relevant cycles. 
We modify the construction of the V-families as introduced by Vismara~\cite{vismaraUnion1997} by rooting at an edge belonging to a fundamental cycle basis rather than a node. 
Compare Figures~\ref{fig:Vis_Fams} and~\ref{fig:mod_vis_fams}.
We call the modified cycle families V'-families.}
%In the original Vismara cycle families the descriptors use node labels. We use edge labels according to some fundamental cycle basis  -- compare Figures~\ref{fig:Vis_Fams} and~\ref{fig:mod_vis_fams}.
%\add{We call the original Vismara cycle families V-families and the modified families V'-families.}
%We call these modified cycle families V'-families.
%\add{We use MV-family as an abbreviation for modified Vismara cycle family.}
%First, recall that we can expand a cycle with respect to a fundamental cycle basis using the edges outside of a spanning tree $T$ -- see Definition~\ref{def:FCB}. Using this observation, we redefine cycle families as follows.
\begin{definition}[V'-family]
	Let $G$ be a graph and $T$ a given spanning tree of $G$. Let $e_j\in E(G)\setminus E(T)$ be an enumeration of the edges outside of the spanning tree $T$.
%	Let $T$ be a given spanning tree of $G$ and let the edges $e_i\in E(G)\setminus E(T)$ be enumerated arbitrarily. 
	\begin{itemize}
		\item The odd V'-family with descriptors $e_j = (u_j,v_j)$ and node $p$, where $d(u_j,p)=d(v_j,p)$, is the set of cycles with largest-indexed edge $e_j$ in $E(G)\setminus E(T)$ composed of $e_j$ and shortest path pairs from $u_j$ and $v_j$ to $p$. See Figure~\ref{fig:mod_vis_fams}, left.
		\item The even V'-family with descriptors $e_j=(u_j,v_j)$ and edge $(p,q)$, where $d(u_j,p)=d(v_j,q)$, is the set of cycles with largest-indexed edge $e_j$ composed of $e_j$, $(p,q)$, and shortest path pairs from $u_j$ to $p$ and from $v_j$ to $q$. See Figure~\ref{fig:mod_vis_fams}, right.
	\end{itemize}
	\label{def:mod_vis_fams}
\end{definition}

\begin{figure}[t]
	\centering
	\includegraphics{VpDescendant.pdf}
	\caption{(Left) A cycle of length 5 whose descriptors are the largest-indexed edge $e_j\in E(G)\setminus E(T)$ for some spanning tree $T$ and the node $p$ antipodal to $e_j$. 
	(Right) A cycle of length 6 whose descriptors are its largest-indexed edge $e_j=(u,v)$ and antipodal edge $(p,q)$. 
	Nodes are colored red or blue if they are closer to $u$ or $v$ respectively. 
	The left node $p$ is purple because it is equidistant to $u$ and $v$.}
	\label{fig:mod_vis_fams}
\end{figure}
%\noindent\add{We will use MV-families as an abbreviation for the modified Vismara cycle families.}


%More succinctly, Definition~\ref{def:mod_vis_fams} modifies Vismara's cycle families by rooting the cycles at an edge $e_j$ instead of a node $r$ -- compare Figures~\ref{fig:Vis_Fams} and~\ref{fig:mod_vis_fams}. 
This has the following benefits:
\begin{itemize}
	\item \textbf{Node labeling.} The shortest paths for the V'-families are computed using breadth-first search starting at the root edge $e_j=(u_j,v_j)\in E(G)\setminus E(T)$. By rooting at an edge, we may recursively label whether a node is closer to $u_j$ or $v_j$. %A node's label is used to determine its role in a family. 
	See the blue and red nodes in Figure~\ref{fig:mod_vis_fams}.
	
	\item \textbf{Computation time.} 
	Repeatedly traversing the graph to construct V'-families is the primary bottleneck of our algorithm.
	To compute the V'-families, breadth-first search is run once per root edge in $E(G)\setminus E(T)$, the number of which is equal to the dimension of the cycle space $\nu$. 
	For the original Vismara families, breadth-first search is run once per node. 
	In our simulations of hydrocarbon pyrolysis, the dimension of the cycle space ranges from $0.01 n \le \nu \le 0.2 n$. 
	This reduces the number of times we traverse the graph by a scaling factor that depends on initial conditions.
%	In our simulations of hydrocarbon pyrolysis the dimension of the cycle space is much smaller than the number of nodes, typically $\nu \le 0.1 n$. Running breadth-first search is the computational bottleneck of our algorithm, so this reduces computation time.
	 
	For diamond, with mean degree $\langle k\rangle=4$, the dimension of the cycle space is $\nu = m - n + 1 \sim n$ since $n\langle k\rangle = 2m$. Thus, even for reasonably dense materials, the number of times we traverse the graph does not increase. This result does not extend to arbitrarily dense graphs where the assumption $m=O(n)$ fails.
	
	\item \textbf{Simple descriptors.} 
	By modifying the original V-families, the descriptors of even cycles change from a node and connected triple to a pair of edges. Connected triples are harder to work with, e.g.,
%	Edges are more convenient to work with than connected triples, e.g., 
	the number of edges $m$ is given whereas the number of connected triples as a function of the degree sequence $\{k_i\}$ is given by $\sum_i \binom{k_i}{2}$.
%	The original descriptors of an even Vismara family are a node and a connected triple. The modified descriptors are both edges, which we argue are more convenient. 
%	For instance, it is easier to count edges than connected triples. As a result, the number of even modified Vismara families is bounded by $\nu m$, where $\nu$ is the number of feasible root edges and $m$ is the number of feasible antipodal edges. By a similar argument, the number of odd modified Vismara families is bounded by $\nu n$. In total, the number of modified Vismara families is bounded by $\nu(m+n)$.
\end{itemize}

All relevant cycles belong to a V'-family, but not all V'-families contain relevant cycles.
The method for filtering out irrelevant V'-families is discussed later in Section~\ref{sub:DICE_TICE_Comp}. 

\subsubsection{An algorithm for computing V'-families}

%{\color{red} Update to match for loop on edges in Algorithm}


%\begin{algorithm}
%	\KwIn{$\bullet\ G$ = an undirected, unweighted graph,\newline
%			$\bullet\ \calF$ = fundamental cycle basis of $G$ with cycles associated with edges $e_k\in E(G)\setminus E(T),$ 
%			\newline\phantom{$\bullet\ \calF$ = }
%			$k=1,\hdots,\nu$, where $T$ is a spanning tree of $G$}
%	\KwOut
%\end{algorithm}

\begin{algorithm*}
	\footnotesize
	\KwIn{$\bullet\ G$ = an undirected, unweighted graph,\newline
		$\bullet\ \calF$ = fundamental cycle basis of $G$ with cycles associated with edges ${e_j}\in E(G)\setminus E(T),$ 
		\newline\phantom{$\bullet\ \calF$ = }
		$j=1,\hdots,\nu$, where $T$ is a spanning tree of $G$}%,\newline
%		$\bullet\ e_{j} = \tt(u,v)$ = a fixed root edge associated with $\mathcal{F}$\;}
	\KwOut{$\bullet$ V'-families represented by their descriptors,\newline
		$\bullet\ D_j=$ a directed, acyclic network for computing shortest paths to $e_{j}$ for $j=1,\hdots,\nu$,\newline
		$\bullet$ node properties {$\tt ancestor_j$}, {$\tt distance_j$}, {$\tt num\_paths_j$} for $j=1,\hdots,\nu$\;}
	
	\For{${e_j}=({\tt u_j,v_j})$ {\normalfont in} $E(G)\setminus E(T)$}{
	Initialize root nodes $\tt ancestor_j[u_j]\gets u_j, distance_j[u_j]\gets 0, num\_paths_j[u_j]\gets 1,$ and $\tt ancestor_j[v_j]\gets v_j, distance_j[v_j]\gets 0, num\_paths_j[v_j]\gets 1$\;
	For each remaining node $\tt x$, mark $\tt observed[x]\gets False$ and $\tt valid[x]\gets True$\;
%	Mark remaining nodes as {\color{red}unobserved}\;
	Initialize (FIFO) $\tt queue\gets [u_j,v_j]$\;
	\For{{\tt p} {\normalfont in} $\tt queue$}{
%		Pop $\tt p$ from the beginning of $\tt queue$\;
		\# {\tt p} is a descriptor for an odd V'-family if it has two ancestors $\tt u_j$ and $\tt v_j$\;
%		{\color{red} \# Check if {\tt p} is a descriptor of an odd cycle family ** odd family because it has two ancestors **}\;
		\uIf{$\tt ancestor_j[p]=\{u_j,v_j\}$ {\normalfont\textbf{and}} $\tt valid[p]$}{
			Count shortest paths from $\tt p$ to $\tt u_j$: $\tt n_1=\sum_{x, p\to x\in D_j}num\_paths_j[x]\cdot \textbf{1}[ancestor_j[x]=u_j]$\;
			Count shortest paths from $\tt p$ to $\tt v_j$: $\tt n_2=\sum_{x, p\to x\in D_j}num\_paths_j[x]\cdot \textbf{1}[ancestor_j[x]=v_j]$\;
			\uIf{$\tt n_1 n_2>0$}{
				Store the odd V'-family with descriptors $\tt e_{j}$ and $\tt p$ representing $\tt n_1 n_2$ cycles of length $\tt 2distance_j[p]{+}1$\;
			}
		}
		\For{$\tt q$ {\normalfont neighboring} $\tt p$ {\normalfont in} $G$}{
			\# {\tt p} is {\tt q}'s parent and {\tt q} is unobserved, initialize {\tt q}'s node properties\;
			\uIf{{\normalfont\textbf{not}} $\tt observed[q]$ }{
				%Mark $\tt q$ as {\color{red} observed} and add {\tt q} to {\tt queue}\;
				Mark $\tt observed[q]\gets True$ and add $\tt q$ to $\tt queue$\;
				$\tt ancestor_j[q]\gets ancestor_j[p]$\;
				$\tt distance_j[q]\gets distance_j[p]+1$\;
				\uIf{${\tt(p,q)} = {\tt e_k}\in E(G)\setminus E(T)$  {\normalfont where} $k>j$}{
					$\tt num\_paths_j[q]\gets 0$\;
				}\uElse{
					$\tt num\_paths_j[q]\gets num\_paths_j[p]$\;
					Add edge $\tt q\to p$ to $D_j$\quad\# child to parent\;
				}
			}
			\# {\tt p} is {\tt q}'s parent and {\tt q} has been observed, update {\tt q}'s node properties\;
			\uElseIf{$\tt distance_j[q]=distance_j[p]+1$}{
%				\uIf{$\tt ancestor[v]\neq ancestor[u]$}{
%					$\tt ancestor[v]\gets\{u_{j_0},v_{j_0}\}$\;
%				}
				$\tt ancestor_j[q]\gets ancestor_j[q]\cup ancestor_j[p]$\;
				\uIf{$\tt ancestor_j[p] = \{u_j,v_j\}$}{
					Mark {\tt q} invalid\;
				}
				\uIf{$\tt num\_paths_j[p]>0$ {\normalfont\textbf{and} \textbf{not}} ${\tt(p,q)} = e_k\in E(G)\setminus E(T)$  {\normalfont where} $k>j$}{
					$\tt num\_paths_j[q] \gets num\_paths_j[q]+num\_paths_j[p]$\;
					Add edge $\tt q\to p$ to $D_j$\;
				}
			}
			\# ({\tt p,q}) is a descriptor of an even V'-family if {\tt p} and {\tt q} have different ancestors and are equidistant to $\tt e_j$, $\tt ancestor_j[p]=u_j$ to break symmetry\;
%			\# {\tt p} and {\tt q} are neighbors equidistant to $\tt e_j$, check if they form an even cycle family\;
			\uElseIf{$\tt ancestor_j[p]=u_j$ {\normalfont\bf and} $\tt ancestor_j[q]=v_j$}{
				\uIf{{\normalfont \textbf{not}} ${\tt(p,q)} = e_k\in E(G)\setminus E(T)$  {\normalfont where} $k>j$}{
					\uIf{$\tt num\_paths_j[p]\cdot num\_paths_j[q]>0$}{
						Include the even C'-family with descriptors $\tt e_{j}$ and $\tt (p,q)$ representing $\tt num\_paths_j[p]\cdot num\_paths_j[q]$ cycles of length $\tt distance_j[p]+distance_j[q]+2$
					}
				}
			}
		}
		}
	}
	\caption{Compute V'-families }%{\color{red} mention Python notation}}
	\label{alg:vis_fam}
\end{algorithm*}

\begin{figure}[t]
	\centering
	\includegraphics{ModVisFamInfo.pdf}
	\caption{\textbf{(a)} The directed acyclic graph with edges pointing along shortest paths towards the root edge $e_j=(u_j,v_j)$. 
	There are two V'-families: one with descriptors $e_j$ and $(p_1,q_1)$ representing 2 cycles of length 6, and one with descriptors $e_j$ and $p_2$ representing 4 cycles of length 9.	
%	There is a modified Vismara family with descriptors $e_i$ and $(p,q)$ holding 2 cycles of length 6, one for each path from $q$ to $v_i$. There is another family with descriptors $e_i$ and $x$ holding 4 cycles of length 9, one for each pair of paths from $x$ to $u_i$ and $x$ to $v_i$. 
%	(b) A modified Vismara cycle family with descriptors $e_j$ and $p_2$ where $p_2$ follows a node $p_1$ with both ancestors $\{u,v\}$. 
%	The paths from $p_2$ to $u$ and $v$ in the exterior cycle $C$ are exchanged for paths that pass through $p_1$ by adding $C_1$ and $C_2$ to $C$ respectively. 
%	Moreover, $C$ is the sum of shorter cycles $C_0,C_1,C_2$ where $C_0$ is a cycle connecting $p_1$ and $e_i$. 
	\textbf{(b)} An odd V'-family with descriptors $e_j$ and $p_2$ where $p_2$ is a descendant of $p_1$ equidistant to $u_j$ and $v_j$. The paths from $p_2$ to $u_j$ and $v_j$ in the outer cycle $C$ are exchanged for paths through $p_1$ by adding $C_1$ and $C_2$ to $C$ respectively. This shows $C$ is not relevant because it is the sum of shorter cycles $C_0,C_1,C_2$.
%	Furthermore, $C$ is not relevant.
%		The blue path from $x_2$ to $u_i$ in the exterior cycle $C$ can be exchanged for a path through $x_1$ using the cycle $C_1$. Similarly, $C_2$ can be used to exchange the path from $x_2$ to $v_i$ for a path that passes through $x_1$. By construction, $C$ is the sum of shorter cycles $C_0,C_1,C_2$ where $C_0$ is a cycle with the descriptors $e_i$ and $x_1$. Thus, $C$ is not relevant.
	}
	\label{fig:ModVisFamInfo}
\end{figure}


%Now, we describe the algorithm for computing the modified Visamara families of a graph $G$. This method requires first computing a spanning tree $T$ and enumerating the edges $e_j\in E(G)\setminus E(T)$. We compute the modified Vismara families with largest indexed edge $e_j$ using Algorithm~\ref{alg:vis_fam}. To compute all of the modified Vismara families this algorithm is iterated over the edges $E(G)\setminus E(T)$. The modified Vismara families are stored using their descriptors. Also, Algorithm~\ref{alg:vis_fam} outputs a directed acyclic graph. The directed edges of this graph point along shortest paths towards the root edge $e_j$ -- see Figure~\ref{fig:mod_vis_fams} (a). Thus, the directed acyclic graph can be traced backwards to join the root edge $e_j$ and its antipodal descriptor $x$ or $(p,q)$ into a cycle.

%We now describe how the modified Vismara cycle families are computed. 
%We now describe how the V'-families are computed. 
%\remove{For brevity, we will call \emph{modified Vismara cycle families} just \emph{cycle families} from here on.}
We assume a spanning tree $T$ is given and the edges $E(G)\setminus E(T)$ are enumerated. 
Algorithm~\ref{alg:vis_fam} computes the V'-families by looping over each edge $e_j\in E(G)\setminus E(T)$.
%Algorithm~\ref{alg:vis_fam} is used to construct the cycle families with a fixed largest-indexed edge $e_j$. 
%This algorithm is looped over all edges outside of $T$ to generate the complete set of cycle families.
%This algorithm is looped over all edges outside of the spanning tree $T$ to generate the full set of modified Vismara cycle families.
%We emphasize this procedure outputs a superset of the relevant cycles, and not all cycle families contain relevant cycles.
%The method for filtering out irrelevant cycles is discussed in Section~\ref{sub:DICE_TICE_Comp}.
%We emphasize these cycle families do not necessarily contain relevant cycles. 
%Instead, we obtain a superset of the relevant cycles

%At a high level, Algorithm~\ref{alg:vis_fam} uses breadth-first search rooted at $e_j=(u,v)$ to obtain shortest paths from all nodes to the edge $e_j$. Here, the distance of a node $x$ to $e_j$ is
%\begin{equation}
%	d(x,e_j) = \min(d(x,u),d(x,v)),
%\end{equation}
%and a shortest path from $x$ to $e_j$ is a minimum length path from $x$ to either $u$ or $v$.
%As a data structure, we save these paths using a directed acyclic graph whose edges point along shortest paths towards $e_j$ -- see Figure~\ref{fig:ModVisFamInfo} (a). 


%The syntax for comments in Algorithm~\ref{alg:vis_fam}. In Figure~\ref{fig:VisFamEx} (d,e) we show two example calls to Algorithm~\ref{alg:vis_fam}.

Algorithm~\ref{alg:vis_fam} outputs V'-families represented by their descriptors from Definition~\ref{def:mod_vis_fams}.
%The cycle families output by Algorithm~\ref{alg:vis_fam} are represented by the descriptors from Definition~\ref{def:mod_vis_fams}. 
The largest-indexed edge $e_j=(u_j,v_j)$ is given, and the $p$ and $(p,q)$ descriptors are computed. 
Nodes $p$ for odd V'-families are those that are equidistant to $u_j$ and $v_j$, obtained in lines 7-11 of Algorithm~\ref{alg:vis_fam}. 
Edges $(p,q)$ for even V'-families are such that $p$ is closer to $u_j$ than $v_j$, $q$ is closer to $v_j$ than $u_j$, and $d(p,u_j)=d(q,v_j)$, obtained in lines 32-35.

\begin{figure*}
	\centering
	\includegraphics[width=6.5in]{VisFamAlgFull.pdf}
	\caption{\textbf{(a)} A graph composed of a six-sided loop joined to a triangular prism with two paths $(3,7,12)$ and $(3,8,12)$ connecting $3$ and $12$. The enumerated edges $e_0,\hdots,e_5$ are used to construct a fundamental cycle basis. 
		\textbf{(b)} V'-families sorted by length %Modified Vismara cycle families -- sorted by length -- 
		output by Algorithm~\ref{alg:vis_fam} iterated over $e_0,\hdots,e_5$. 
%	The edge descriptors $e_j=(u,v)$ and $(p,q)$ are shown in black and the node descriptor $p$ for odd families is shown in purple.
%		The descriptors $e_j=(u,v)$ and $(p,q)$ for even families are shown as black edges, and the descriptor $p$ for odd families is shown as a purple node. 
	The only V'-family with multiple cycles is $\{C_7,C_7'\}$.
	The secondary cycle $C_7'$ is plotted using dashed edges. 
%	The representative cycle $C_7$ is plotted using solid edges, and the secondary cycle $C_7'$ uses dashed. 
%	All cycle families except $\{C_7,C_7'\}$ are composed of a single cycle. The representative cycle $C_7$ from $\{C_7,C_7'\}$ is plotted using solid edges and the secondary $C_7'$ is plotted using dashed edges.
	\textbf{(c)} Matrix of representative cycles from each V'-family as fundamental cycle basis vectors.
	\textbf{(d)} Depiction of Algorithm~\ref{alg:vis_fam} rooted at $e_4$. Each frame represents a step through the central for loop over nodes $\tt p$ (circled in black). Steps ${\tt p}=5,6,4$ are omitted.
	Node colors denote ancestory, node labels denote distance to $e_4$, and directed edges point along valid shortest paths to $e_4$.
%	Node colors represent the ancestor property. Node labels denote the distance to $e_4$. The directed acyclic graph is shown using directed edges. 
%	Steps $p=5,6,4$ are omitted because they do not discover new nodes. 
	\textbf{(e)} Depiction of Algorithm~\ref{alg:vis_fam} rooted at $e_5$, where node labels denote the number of valid shortest paths to $e_5$.
	}
	\label{fig:VisFamEx}
\end{figure*}

%The largest-indexed edge $e_j=(u,v)$ in the output families is supplied as input to Algorithm~\ref{alg:vis_fam}.
%To find an odd cycle family, we identify complementary nodes $p$ that are equidistant to $u$ and $v$ as in line 6 of Algorithm~\ref{alg:vis_fam}.
%For an odd cycle family, the complementary descriptor to the root edge $e_j=(u,v)$ is a node $x$ equidistant to $u$ and $v$, as in line 6. 
%For an even cycle family, we identify edges $(p,q)$ such that $p$ is closer to $u$ than $v$, $q$ is closer to $v$ than $u$, and $d(p,u)=d(q,v)$.
%For an even cycle family, the complementary descriptor is an edge $(p,q)$ where 
%$p$ is closer to $u$ than $v$, $q$ is closer to $u$ than $v$, and $d(p,u)=d(q,v)$. See line 28.
%$p$ and $q$ are closer to $u$ and $v$, respectively, and $d(p,u)=d(q,v)$, as in line 28.

The cycles in a V'-family are constructed by joining the pair of descriptors by shortest paths as in Figure~\ref{fig:mod_vis_fams}. Algorithm~\ref{alg:vis_fam} computes shortest paths from all nodes to $e_j = (u_j,v_j)$ using breadth-first search. Here, we define the distance to $e_j$ as
\begin{equation}
	d(x,e_j) = \min(d(x,u_j),d(x,v_j)),
\end{equation}
%The actual cycles in a family are constructed by finding a pair of shortest paths to the nodes $u$ and $v$, and joining these paths to their descriptors. Algorithm~\ref{alg:vis_fam}, computes shortest paths from every node to the root edge using breadth-first search. Here, it is useful to define the distance from a node $x$ to the edge $e_j=(u,v)$ as
%\begin{equation}
%	d(x,e_j) = \min(d(x,u),d(x,v)),
%\end{equation}
and a shortest path from $x$ to $e_j$ as a minimum length path from $x$ to either $u$ or $v$.
As a data structure, we store a directed acyclic graph whose edges point along shortest paths to $e_j$ -- see Figure~\ref{fig:ModVisFamInfo} (a).
%As a data structure, we compute a directed acyclic graph whose edges point along shortest paths towards $e_j$ -- see Figure~\ref{fig:ModVisFamInfo} (a). 
We call a neighbor $q$ of node $p$ a child of $p$ if it is further from the root edge $e_j$. 
%We will refer to a neighbor $q$ of node $p$ as a child of $p$ if $d(q,e_j)>d(p,e_j)$. 
Edges in the directed acyclic graph point backwards from child to parent, which are assigned in lines 22 and 30.

We compute three additional node properties that further characterize the V'-families:
%These are listed below. Each of these node properties are computed recursively.

%The cycle families are obtained
%The modified Vismara cycle families are obtained by finding complementary descriptors to the root edge $e_j$. 
%An odd modified Vismara families is obtained via a node $p$ that are equidistant to $u$ and $v$. The even Vismara cycle families are obtained from a pair of neighboring nodes $p$ and $q$ which are closer to $u$ and $v$, respectively, where $d(p,u)=d(q,v)$.

%At a base level, Algorithm~\ref{alg:vis_fam} uses breadth-first search to compute shortest paths for the modified Vismara families. The complexity of Algorithm~\ref{alg:vis_fam} lies in the number of variables that need to be tracked: 

\begin{itemize}
%	\item Node properties
%	\begin{itemize}
		\item \textbf{Ancestor.} 
		A node's ancestor is the node $u_j$ or $v_j$ in the root edge $e_j$ closest to it. 
%		The ancestor of a node $q$ is the node in the root edge $e_j=(u,v)$ that is closer to $q$. 
%		The ancestor of a child node $q$ is inherited from its parent $p$, as in line 16. 
		If a node is equidistant to $u_j$ and $v_j$, then both $u_j$ and $v_j$ are its ancestors.
%		It is possible that the node $u$ will be equidistant to both nodes $u_i$ and $v_i$. If $u$ is equidistant to $u_i$ and $v_i$, then its ancestors are $\{u_i,v_i\}$. 
		The ancestors of a child node $q$ is the union of its parent's ancestors, computed in lines 16 and 25.
%		Node ancestors are used to characterize the role of nodes in a V'-family.
		Descendants of a node with both ancestors produce irrelevant cycles that can be decomposed into smaller cycles. See, Figure~\ref{fig:ModVisFamInfo} (b). We mark these descendants invalid in lines 26-27.
%		To achieve this, the ancestors of $q$ are given by the union of its parents ancestors, as in line 25.
%		The descendants of a node. 
%		The descendants of a node with both ancestors output cycles that are not relevant because they can be decomposed into smaller cycles. See Figure~\ref{fig:ModVisFamInfo} (b). We mark these descendants invalid in lines 26-27. 
		
%		If a node follows a node with both ancestors $\{u,v\}$ then it cannot be used to construct a relevant cycle, e.g., the cycle $C$ in Figure~\ref{fig:ModVisFamInfo} (b) is not relevant because it is the sum of shorter cycles. We mark all nodes following a node with both ancestors invalid, as in lines 23-24. Cycle families with invalid nodes are rejected, as in lines 6 and 28.
		\item \textbf{Distance.} %By the nature of breadth-first search, 
		A node $q$ will first be observed by a neighbor $p$ lying on a shortest path to $e_j$. Thus, $q$ is one edge further from $e_j$ than its parent node $p$, as in line 15. This distance is used to compute the lengths of cycles in a family, as in lines 11 and 35. %The distance from $u$ to the root edge $e_j=(u_j,v_j)$ is given by
%		\begin{equation}
%			d(u,e_j) = \min(d(u,u_j),d(u,v_j)).
%		\end{equation}
%		The distance of a child node $v$ is inherited from its parent $u$ as in line 15. This is used to compute the length of cycles in a family as in lines 10 and 31.
		\item \textbf{Number of paths.} %In addition to the directed acyclic graph, we store the number of shortest paths from a node $x$ to the root edge $e_j$. 
%		This value is computed recursively by summing over the number of shortest paths from each of the parents of $x$ to $e_j$ -- see lines 19 and 26.
		We count shortest paths from each node $q$ to $e_j$ as the sum of paths through its parents. See lines 19 and 26.
		This is used to compute the number of cycles in a V'-family, which is equal to the number of shortest path pairs to $u_j$ and $v_j$. 
		The size of a cycle family is computed in lines 8-11 and 35.
%		We count shortest paths from each node $q$ to % the root edge 
%		$e_j$ recursively by the sum of shortest paths from $p$ to $e_j$ for each parent $p$ of $q$.
%		See lines 19 and 26.
%		We also store the number of shortest paths from a node $u$ to the root edge. 
%		This is given by the sum of shortest paths through its parent nodes as in lines 19 and 26. 
%		This is used to count the number of cycles in a family. The number of cycles in an even family is given by the number of shortest paths from $p$ to $u$ times the number of shortest paths from $q$ to $v$, as in line 31. 
%		We use this value to compute the number of cycles in a family. For instance, the number of cycles in a even cycle family is given by the number of shortest paths from $p$ to $u$ times the number of shortest paths from $q$ to $v$, as in line 31. 
%		Similarly, 
%		The number of cycles in an odd family is the number of shortest paths from $p$ to $u$ times the number %of shortest paths 
%		from $p$ to $v$, as in lines 7-10.
%		The number of cycles in a family is given by the number of path pairs to $u$ and to $v$ (i.e., red and blue paths in Figure~\ref{fig:mod_vis_fams}) as in lines 10 and 31.
%		\item \textbf{Directed acyclic graph.} We assign directed edges from child nodes to parent nodes in lines 20 and 27. The directed graph is used to construct shortest paths to $u$ and $v$. These paths are used to construct cycles as in Figure~\ref{fig:mod_vis_fams}.
%		\item \textbf{Largest-indexed edge $\bm{e_j}$.} A cycle cannot include an edge $e_k\in E(G)\setminus E(T)$ with index $k > j$. We reject paths including these edges in lines 16-17, 23-24, and 29. It is possible that a family will not have a valid pair of paths to $u$ and to $v$. In this case, the empty cycle family is rejected as in lines 9 and 30.
%		\item \textbf{Identifying modified Vismara families.} An odd modified Vismara family is obtained via a node $p$ that has both $u$ and $v$ as ancestors as in line 6. An even modified Vismara family is obtained via a pair of nodes $p,q$ with distinct ancestors. Specifically, we choose the ancestors of $p$ and $q$ to be $u$ and $v$, respectively, to avoid double counting. See line 28.
%	\end{itemize}
\end{itemize}

As a final note, in each loop of Algorithm~\ref{alg:vis_fam} we compute cycles with largest-indexed edge $e_j$. We reject any path that includes an edge $e_k\in E(G)\setminus E(T)$ with index $k>j$ as in lines 16-17, 23-24, and 29. We still traverse edges $e_k$, $k>j$, in Algorithm~\ref{alg:vis_fam} to accurately compute the distance and ancestor node properties.

As a post-processing step, we sample a representative cycle from each V'-family. We order these cycles, and the associated V'-families, by length, i.e., $|C_i|\le |C_j|$ for $i\le j$. See Figure~\ref{fig:VisFamEx} (b). 
% These will be used in later sections. 
%First, we sample $C$ from a cycle family as a set of edges by obtaining a pair of paths from the directed acyclic graph discussed above. 
%{\color{red} Starting here more clear}
%It is convenient to order the cycles $\{C_i\}$ and their cycle families by increasing length, i.e., $|C_i|\le |C_j|$ for $i\le j$.
%by constructing a pair of paths from the directed acyclic graph.
We convert each cycle $C_i$ into a fundamental cycle basis vector using the non-tree edges $e_j\in E(G)\setminus E(T)$ that lie in $C_i$. 
Then, we construct a matrix whose rows are representative cycles from each V'-family as fundamental cycle basis vectors:
%These fundamental cycle basis vectors are stored as the rows of a matrix:
%{\color{red} matrix of representative cycles from each family as $\calF$-vectors... this is more clear}
%As a post-processing step, we sample a representative cycle from each cycle family. These representative cycles will be used in later sections. These cycles are sampled as a set of edges $C$ by obtaining a pair of paths to $u$ and $v$ using the directed acyclic graph and joining them to the descriptors 
%In the following section, we will need to sample a representative cycle from each cycle family. For each family we extract an arbitrary pair of paths to the root edge $e_j$ and use them to construct a cycle as a set of edges $C$. It is convenient to order these cycles $\{C_i\}$ and their respective families by length, i.e., $|C_i|\le|C_k|$ for $i<k$. Then each cycle $C_i$ is represented as a fundamental cycle basis vector by identify the edges $e_j\in E(G)\setminus E(T)$ that lie inside $C_i$. We store these fundamental cycle basis vectors as the rows of a matrix
\begin{equation}
	\bm{R_0}(i,:) = (\bm{1}[e_1{\in} C_i], \bm{1}[e_2{\in} C_i], \hdots, \bm{1}[e_\nu{\in} C_i]),%, \quad\text{for $e_j\in E(G)\setminus E(T)$, $C_i$ representing family $i$.}%\text{\color{red} for\ $e_j$\ in ...} \text{\color{red} for set of cycles}.
	\label{eq:R0_def}
\end{equation}
where $C_i$ represents family $i$. 
See Figure~\ref{fig:VisFamEx} (c).

\subsubsection{The computational cost}

The complexity of algorithms in this section depend on the size of $\bm{R_0}$. The number of columns in $\bm{R_0}$ is equal to the dimension of the cycle space $\nu$. The number of rows in $\bm{R_0}$ is equal to the number of V'-families. $\bm{R_0}$ has at least $\nu$ linearly independent rows that can be used to construct an MCB.
The number of V'-families is bounded by
%The number of rows in $\bm{R_0}$ is equal to the modified Vismara cycle families and the number of columns in $\bm{R_0}$ is equal to the dimension of the cycle space.
%The cycles in $\bm{R_0}$ can be used to construct a minimum cycle basis, so $\bm{R_0}$ has at least $\nu$ linearly independent rows. The number of rows in $\bm{R_0}$ is bounded above by
\begin{equation*}
	\begin{split}
		\nu \,&{\le}\, \textbf{\# of V'-families} \\&{\le}\,\big|E(G){\setminus}E(T)\big|\,(\#\text{V'-families per edge})\,{\le}\,\nu^{\mathrlap{2}}.
	\end{split}
\end{equation*}
We justify that Algorithm~\ref{alg:vis_fam} outputs at most $\nu$ cycles in each of its $\nu$ iterations below. This bound is tight, where the complete bipartite graph has $O(\nu^2)=O(n^4)$ cycle families. 

To bound the number of V'-families per iteration of Algorithm~\ref{alg:vis_fam}, consider a breadth-first spanning tree $T_j$ rooted at $e_j$. %with the same root edge $e_j$ as in Algorithm~\ref{alg:vis_fam}. 
Note, $T_j$ is not the spanning tree $T$ used to construct a fundamental cycle basis.
%This tree is different from the given spanning tree $T$ used to describe a fundamental cycle basis. 
$T_j$ is a subgraph of the directed acyclic graph $D_j$ used to construct paths to $e_j$, as in Figure~\ref{fig:ModVisFamInfo}, where each node has {out-degree} 1. 
The edge descriptor $(p,q)$ for an even V'-family lies outside $T_j$. 
The node descriptor $p$ for an odd V'-family has disjoint shortest paths to $u_j$ and $v_j$, so it has out-degree greater than or equal to 2 in $D_j$. 
Hence, one of these outgoing edges lies outside of $T_j$.
%Thus, $p$ has out degree greater than or equal to 2 in the directed acyclic graph, and one of these outgoing edges lies outside of $T_j$.
%Furthermore, the number of valid descriptors for constructing cycle families is bounded above by the number of edges outside of the tree $T_j$, i.e., by $\nu$.
In both cases, we associate an edge outside of the tree $T_j$ with the descriptor $p$ or $(p,q)$, and so the number of V'-families with root edge $e_j$ is bounded by the number of non-tree edges $\nu$.

In summary, $\bm{R_0}$ is a $\{0,1\}^{O(\nu^2)\times\nu}$ matrix. 
Assuming dense matrix operations, the computational cost of the algorithms in the following sections is $O(\nu^4)$, i.e., the cost of na\"ively multiplying $\bm{R_0}$ by a $\nu\times\nu$ matrix.
However, we find $\bm{R_0}$ is not too large and typically very sparse. 
The number of rows in this matrix is typically closer to the dimension of the cycle space. 
In a loop cluster sampled from a carbon-rich simulation of hydrocarbon pyrolysis (initial composition $\rm C_{10}H_{16}$), the dimension of the cycle space is $\nu=1304$ and the number of cycle families is $1604$.
This observation agrees with other chemical systems where the number of relevant cycles is approximately equal to $\nu$, see Table 8 in Ref.~\cite{mayEfficient2014}.
%For instance, in a representative sample of the giant bicomponent from a carbon-rich simulation of hydrocarbon pyrolysis (initial composition $\rm C_{10}H_{16}$), the dimension of the cycle space is $\nu=1304$ and the number of cycle families is $1604$. 
%This agrees with other chemical systems where the number of relevant cycles is approximately equal to $\nu$~\cite{mayEfficient2014}. 
%$\bm{R_0}$ is sparse because the relevant cycles are short, bounding the number of edges $e_j\in E(G)\setminus E(T)$. The $1604\times 1304$ $\bm{R_0}$ matrix described
%Also, relevant cycles are typically short, bounding 
Also, relevant cycles are typically short, bounding the number of fundamental cycle basis edges $e_j\in E(G)\setminus E(T)$ they have. 
As a result, their fundamental cycle basis vectors and the matrix $\bm{R_0}$ are very sparse. The $1604\times 1304$ $\bm{R_0}$ matrix described above has $7280$ non-zero entries. 
In the case where $\bm{R_0}$ is dense, breadth-first search may be used to compute matrix vector products $\bm{R_0}x$ in $O(\nu^2)$ time, as discussed in Section 5.2 in Ref.~\cite{mehlhornMinimum2010}, which can be used to reduce the complexity of matrix multiplication to $O(\nu^3)$. 


Using our Python codes on a Ryzen 9 9900X CPU, the computation of the matrix $\bm{R_0}$ for the system mentioned above takes approximately 5 seconds.
%the matrix $\bm{R_0}$ can be computed in around 5 seconds for the system described above. 
For comparison, a sparse matrix-matrix product $\bm{R_0 W}$ is computed in less than 0.5 seconds. Indeed, Algorithm~\ref{alg:vis_fam} is the main bottleneck of our approach. The complexity of a single iteration of Algorithm~\ref{alg:vis_fam} is equal to the complexity of breadth-first search, i.e., $O(m)$. The total cost of Algorithm~\ref{alg:vis_fam} is $O(\nu m)$. 
For the $\rm C_{10}H_{16}$ system, $\nu m = 0.32 n^2$.
For diamond, as a proxy for dense material systems, $\nu m\approx 2n^2$. Therefore, it is reasonable to write the complexity as $O(n^2)$ for material networks. This procedure can be further accelerated by running Algorithm~\ref{alg:vis_fam} in parallel, because no communication between iterations over $e_j$ in Algorithm~\ref{alg:vis_fam} is required.
%{\color{red} Say this in a less scary way.}
%In our present setting, we find that $\bm{R_0}$ is extremely sparse. The total time complexity of the computations in the following sections is $O(\nu^4)$... but in practice Algorithm~\ref{alg:vis_fam} is the most expensive step we consider. Each iteration of breadth-first search has complexity $O(m)$ as in Algorithm~\ref{alg:vis_fam}, which is repeated $\nu$ times. Thus, in the setting of hydrocarbon pyrolysis, the complexity our procedure is dominated by a $O(\nu m)=O(n^2)$ algorithm. %The complete bipartite graph, as in Figure..., has 
%{\color{red} Mention our system. Mention parallelization.}

%The node descriptor $p$ for an odd cycle family has out-degree
%The node descriptor $p$ for an odd cycle family has out-degree greater than or equal to $2$,  shortest paths to $u$ and to $v$. Thus, we can also associate to $p$ at least one 
% Since Algorithm~\ref{alg:vis_fam} is an iteration of breadth-first search, we could also consider an associated spanning tree $T_j$. This tree is different than the given tree $T$ used to construct our fundamental cycle basis. The breadth0first search tree $T_j$ is a subset of the directed-acyclic graph used to construct, as in Figure~\ref{fig:ModVisFamInfo} (a), such that all nodes have out degree 1. The edge descriptor $(p,q)$ for an even cycle family necessarily lies outside of $T_j$. Similarly,


%-------
%
%The first step of our procedure is to compute a super-set of the relevant cycles. We will later describe methods for extracting the relevant cycles from this super-set as well as an algorithm for splitting the relevant cycles into DICE and TICE classes. To do so we use methods introduced in Ref.~\cite{vismaraUnion1997}.
%
%We outline the approach used to compute relevant cycles as given in Ref.~\cite{vismaraUnion1997}. The key idea of this approach is to partition the relevant cycles into \textit{families} that share the same start and endpoint. Here, a family of cycles share a common edge associated with a given fundamental cycle basis.
%
%Describe algorithm shortly, remove a few if statements to make it read easily
%
%
%The choice of making cycle families start from an edge has a couple of minor benefits. For one, in the setting of hydrocarbon pyrolysis (where we take sample graphs) this reduces the amount of times we need to run breadth-first search. The number of times we run breadth-first search is equal to the number of edges associated with a fundamental cycle basis or the number of loops $\nu$. In our samples, this is roughly one order of magnitude smaller than the number of nodes $n$, which is the number of times needed to run breadth-first search . Second, it is easier to bound the number of relevant cycle families. Odd cycles families are defined by one edge and one node and even cycle families by two edges resulting in the bound
%\begin{equation}
%	\text{\# families} \le \underbrace{nm}_\text{odd families} + \underbrace{\binom{m}{2}}_\text{even families} = O(m^2).
%\end{equation}
%Last, this emphasizes the method of computing families of relevant cycles in Ref.~\cite{vismaraUnion1997} is flexible
%
%
%
%Discuss representation w.r.t. FCB
%
%Ultimately a size $\nu^2\times \nu$ matrix computed in $O(\nu^3)$ time.
%
%
%{\color{red}
%
%
%
%Describe parameters one-by-one
%
%}

\subsection{Minimum cycle basis}
\label{sub:MCB_comp}

%{\color{red} use $\oplus$ instead of mod 2}

Our tests for relevance and for the {\tt pi} and {\tt sli} relations rely on expanding cycles in terms of an MCB. See Lemma~\ref{Lem:Relevant_MCB_Test}, Theorem~\ref{Thm:DirInt_Comp}, and Lemma~\ref{Lem:TrEx_Alg} respectively. 
Here, we construct an MCB.
%Thus, it is critical that we are able to expand cycles in terms of a minimum cycle basis. First, we outline a method for computing minimum cycle bases.%to extract a minimum cycle basis from the cycle families computed in the previous subsection.

There is a large body of literature dedicated to the computation of MCBs. See Refs.~\cite{kavithaCycle2009,mehlhornMinimum2010} for a review. %In practice, we find that computing cycle families as outlined in the previous section is more computationally expensive. 
We opt for the approach in Ref.~\cite{kavithaFaster2004}, which we review below. This approach is computationally efficient, easy to implement in our present setting, and can be used to expand cycles in terms of the MCB.

This approach relies on an orthogonality argument using the inner product of binary vectors in the GF(2) field:
\begin{equation}
	\langle A,B\rangle_\mathcal{F} := \bigoplus_i a_i b_i.
\end{equation} 
We emphasize that this inner product is defined in terms of a fundamental cycle basis $\calF$. In this context, a cycle $C$ should be interpreted as a binary vector in $\{0,1\}^\nu$ where entry $i$ indicates whether $C$ contains the $i$th edge outside of the spanning tree used to construct $\calF$.
%, i.e. the $i$'th cycle of $\calF$ is in the expansion of $C$.
%Here, we assume that cycles are expanded in terms of a fundamental cycle basis $\mathcal{F}$, i.e. as binary vectors in $\{0,1\}^{m-n+1}$ where entry $i$ indicates whether edge $i$ outside of a given spanning tree is in our cycle. 
See Figure~\ref{fig:FCB}~(b) for a reminder of this decomposition.

%{\color{red} how do the Vismara cycle families relate to MCB... These representative cycles have a MCB, and redundant }


This construction of an MCB is iterative. Suppose cycles $B_1,\hdots, B_{j-1}$ belonging to an MCB are given. 
%Let $\calM = \{B_1,\hdots,B_{j-1}, \tilde{B}_j,\hdots,\tilde{B}_\nu\}$ be such a minimum cycle basis where $\tilde{B}_j,\hdots, \tilde{B}_\nu$ are unknown. 
To obtain the next basis cycle, we utilize a \emph{witness vector} defined as a non-zero vector $S_j\in\{0,1\}^\nu$ satisfying the orthogonality condition
\begin{equation}
	\langle B_i, S_j\rangle_\calF = 0,\quad i<j.
	\label{eq:Wit_Orth}
\end{equation}
{We construct the vectors $S_j$ later in this section.}
The next cycle $B_j$ in the MCB is given by the optimization problem
\begin{equation}
	\min_{B\in\CycSp} |B|,\text{ such that }\langle{B,S_j}\rangle_\calF =1
	\label{eq:len_const_opt}
\end{equation}
where $\CycSp$ is the set of all cycles. 
We obtain this cycle from the first row of the matrix $\bm{R_0}$ from Eq.~\eqref{eq:R0_def} with inner product 1.
%We perform this optimization problem over the set of representative cycles from each family, stored as the matrix $\bm{R_0}$ from Eq.~\eqref{eq:R0_def}.

%{\color{red} Move this as a lemma in appendix
%
%To verify $B_1,\hdots,B_j$ belong to a minimum cycle basis consider an arbitrary minimum cycle basis $\tilde{\calM} = \{B_1,\hdots,B_{j-1},\tilde{B}_j,\hdots,\tilde{B}_\nu\}$.
%In the following identity we expand $B_j$ in $\tilde{\calM}$ and take inner product with $S_j$ %{\color{red} Define $S_j$'s construction first}
%\begin{equation}
%	1 = \langle B_j , S_j \rangle_\calF = %\left\langle\bigoplus_{{B\in\expand{\tilde{\calM}}{B_j}}} B,S_j \right\rangle_\calF =
%	\sum_{B\in\expand{\tilde{\calM}}{B_j}}\langle B, S_j \rangle_\calF,
%	\label{eq:Bj_Exp}
%\end{equation}
%where the inner product moves into the sum by linearity.
%Matching both sides of Eq.~\eqref{eq:Bj_Exp}, there exists a cycle $\tilde{B}\in\tilde{\calM}$ in the expansion of $B_j$ in $\tilde{\calM}$ such that $\langle\tilde{B},S_j\rangle=1$.
%By Eq.~\eqref{eq:Wit_Orth}, $\tilde{B}$ is not one of the cycles $B_1,\hdots,B_{j-1}$. 
%The relation $|B_j|\le |\tilde{B}|$ is given by Eq.~\eqref{eq:len_const_opt}.
%The reverse relation $|\tilde{B}|\le|B_j|$ is given by statement 1 of Lemma~\ref{Lem:Relevant_MCB_Test} since $\tilde{B}$ is in the expansion of $B_j$ in $\tilde{M}$. 
%Thus $|B_j|=|\tilde{B}|$, and we can exchange $B_j$ for $\tilde{B}$ to obtain a new minimum cycle basis $\calM=(\tilde{\calM}\setminus\{\tilde{B}\})\cup\{B_j\}$ containing $B_1,\hdots,B_j$ by statement 2 of Lemma~\ref{Lem:Relevant_MCB_Test}.
%%By Eq.~\eqref{eq:len_const_opt}, we have the relation
%}

The last ingredient of this algorithm is to construct the witness vectors $S_j$. 
%These are constructed procedurally. 
%{\color{red} remove $e_j$ and just say standard unit vector... $e_j$ is an edge also not needed}
These vectors are computed dynamically, where the initial vectors $S_j,j=1,\hdots,\nu$, are the unit vector with value $1$ at coordinate $j$.
%These vectors are initialized as $S_j^{(1)}$, where $e_j$ is the unit vector with value $1$ at coordinate $j$. 
The witness vectors are updated each time a new basis vector $B_j$ is constructed. 
%At step $j$, the vectors $S_k=S_k^{(j)}$ for $k\le j$ are assumed to satisfy the orthogonality condition in Eq.~\eqref{eq:Wit_Orth}, and $B_j$ is drawn according to Eq.~\eqref{eq:len_const_opt}.
The future witness vectors are assumed to be orthogonal to the previous cycles, i.e., $\langle B_j,S_k\rangle=0$ for $j<k$.
At step $j$ if $\langle B_j,S_k\rangle=1$ for $k>j$ then we update $S_k$ as% by the rule
%The future witness vectors, $S_k^{(j)}$ for $k>j$, are updated as follows
\begin{align}
	S_k \gets S_k\oplus S_j.
\end{align}
%\begin{equation}
%	S_k^{(j+1)} \gets \begin{cases}
%		S_k^{(j)}\oplus S_j, & \langle B_j, S_k^{(j)}\rangle=1\\
%		S_k{(j)}, & \langle B_j, S_k^{(j)}\rangle=0.
%	\end{cases}
%\end{equation}
This ensures $\langle B_j,S_k\rangle=0$ for $k>j$.
%If $S_k^{(j)}$ and $B_j$ are not orthogonal, then after the update they are orthogonal by
%\begin{equation}
%	\begin{split}
%		\langle B_j,S_k^{(j+1)}\rangle &= \langle B_j,S_k^{(j)}\rangle+\langle B_j,S_j\rangle \\&= 1+1 = 0,\quad \text{(mod 2)}.
%	\end{split} 
%\end{equation}
%If a future witness vector is not orthogonal to the drawn cycle $\langle B_j, S_k^{(j)}\rangle=1$, $j<k$, then it is updated as follows
%\begin{equation}
%	S_k^{(j+1)} = S_k^{(j)}\oplus S_j.
%	\label{eq:Wit_update}
%\end{equation}
%The updated witness $S_k^{(j+1)}$ is orthogonal to $B_j$ by
%%The drawn cycle $B_j$ and the updated witness vector $S_k^{(j+1)}$ are orthogonal by
%\begin{equation}
%	\langle B_j,S_k^{(j+1)}\rangle = \langle B_j,S_k^{(j)}\rangle+\langle B_j,S_j\rangle = 1+1 = 0\quad \text{(mod 2)}.
%\end{equation}
%where $\langle B_j,S_j\rangle=1$ by Eq.~\eqref{eq:len_const_opt}. 
%Also, $S_k^{(j+1)}$ remains orthogonal to previous cycles, $B_\ell$ for $\ell<j$, by Eq.~\eqref{eq:Wit_Orth}
%\begin{equation}
%	\langle B_\ell,S_k^{(j+1)}\rangle = \langle B_\ell,S_k^{(j)}\rangle + \langle B_\ell,S_j\rangle = 0+0.
%\end{equation}
%If the future witness $S_k^{(j)}$ is already orthogonal
%When we draw cycle $B_j$ it is assumed the witness vectors $S_k = S_k^{(j)}$, $k\le j$, already satisfy Eqs.~\eqref{eq:Wit_Orth} and~\eqref{eq:len_const_opt}. After drawing $B_j$

\begin{algorithm}[t]
	\KwIn{$\bullet\ G = $ simple, connected graph \newline$\phantom{\bullet\ G=}$ with $\nu$ independent cycles\newline
%		$\bullet\ \bm{R_0} = $ matrix of representative cycles from each family as $\calF$-vectors
		$\bullet$ Matrix of cycles $\bm{R_0}$ from V'-\newline $\phantom{\bullet}$ families as fundamental cycle\newline $\phantom{\bullet}$ basis vectors
		}
	\KwOut{$\bullet\ $MCB $\calM = \{B_1,\hdots,B_\nu\}$\newline 
		$\bullet$ Witness vectors $\{S_1,\hdots,S_\nu\}$}
	Initialize $S_i$ to the standard unit vector in $\{0,1\}^\nu$ that is $1$ at coordinate $i$ and $0$ elsewhere\;
%	Initialize witness vectors as elementary vectors $S_i\gets e_i$ for $i=1,\hdots,\nu$\;
	\For{$j=1,\hdots,\nu$}{
		\# Basis cycle $B_j$ as first cycle $C_k$ from $\bm{R_0}$ with $\langle C_k,S_j\rangle_\calF=1$\;
		$B_j\gets$ row $k$ of $\bm{R_0}$ where $k$ is the first non-zero entry of $\bm{R_0} S_j$\;
		%(minimum cycle subject to $\langle C_j,S_j\rangle=1$)\;
		\For{$k=j+1,\hdots,\nu$}{
			\# Ensure $\langle C_j,S_k\rangle_\calF=0$ for $j<k$\;
			\uIf{$\langle C_j,S_k\rangle_\calF=1$}{
				$S_k\gets S_k\oplus S_j$\;
			}
		}
	}
	\caption{Compute MCB}
	\label{alg:MCB}
\end{algorithm}


To implement the approach above, we utilize the V'-families constructed in Section~\ref{sub:Vis_Fam}. This is shown in Algorithm~\ref{alg:MCB}. The cycles $B_j$ are drawn from the matrix $\bm{R_0}$ whose row vectors are the expansions of representative cycles from each family expanded with respect to our fundamental cycle basis $\calF$. If $C_k$ is the cycle represented by the $k$th row of $\bm{R_0}$, then the inner product $\langle C_k, S_j\rangle_\calF$ is the $k$th entry in the matrix-vector product $\bm{R_0} S_j$. Also, the cycles in $\bm{R_0}$ are sorted by length. Thus, the optimization problem in Eq.~\eqref{eq:len_const_opt} is solved by the row vector in $\bm{R_0}$ corresponding to the first non-zero entry in $\bm{R_0} S_j$, as in line 4 of Algorithm~\ref{alg:MCB}.



The main operations in Algorithm~\ref{alg:MCB} are matrix-vector products $\bm{R_0}S_j$, inner products $\langle C_j,S_k\rangle$, and vector additions $S_k\oplus S_j$.
If the entries are dense, the time complexity of Algorithm~\ref{alg:MCB} is $O(\nu^4)$. However, the actual cost of these operations is much less because $\bm{R_0}$ is very sparse. Additionally, the matrix-vector products $\bm{R_0} S_j$ only need to be computed up to the first non-zero entry.

\begin{lemma}
	The cycles $B_1,\hdots,B_\nu$ computed by Algorithm~\ref{alg:MCB} form an MCB.
\end{lemma}

\begin{proof}
	The correctness of Algorithm~\ref{alg:MCB} is proved via induction. Suppose $B_1,\hdots,B_{j-1}$ for $j\ge 1$ belong to an MCB $\calM_{j-1}$. The following identity is obtained by expanding $B_j$ in $\calM_{j-1}$ and taking the inner product with $S_j$
	\begin{equation}
		1 = \langle B_j,S_j\rangle_\calF = \enspace\bigoplus_{\mathclap{\calB\in\expandd{\calM_{j-1}}{B_j}}}\enspace \langle B,S_j\rangle_\calF.
	\end{equation}
	The cycles $B_1,\hdots,B_{j-1}$ are orthogonal to $S_j$. Hence, there exists $B\in\calM_{j-1}\setminus\{B_1,\hdots,B_{j-1}\}$ in the expansion of $B_j$ in $\calM_{j-1}$ such that $\langle B,S_j\rangle_\calF=1$. $|B|\le |B_j|$ because $B$ is in the expansion of $B_j$ in $\calM_{j-1}$. $|B_j|\le|B|$ because $B_j$ is a minimum length cycle such that $\langle B_j,S_j\rangle_\calF=1$. Therefore, $\calM_j=(\calM_{j-1}\setminus\{B\})\cup\{B_j\}$ is an MCB by Lemma~\ref{Lem:Relevant_MCB_Test}. By induction, $\calM_\nu=\{B_1,\hdots,B_\nu\}$ is an MCB.
\end{proof}

%The matrix-vector products $\bm{R_0} S_j$ can be efficiently computed by performing breadth-first search on each root edge $e_j$
% {\color{red} If the entries are dense, }Combined these suggest $O(\nu^4)$ time complexity where $\bm{R_0}$ is a $\{0,1\}^{O(\nu^2)\times \nu}$ matrix. Because of sparsity, the actual run time of Algorithm~\ref{alg:MCB} is much faster, and it is much faster than the computation of cycle families in Algorithm~\ref{alg:vis_fam}. Additionally, $\bm{R_0}\cdot S_j$ only needs to be computed to the first non-zero entry.



%The cycles $B_j$ are drawn from the representative cycles from each family discussed in Section~\ref{sub:Vis_Fam}. 
%Specifically, we use the matrix

%We use the cycles extracted in the previous section to construct a fundamental cycle basis.
%In particular, we use the matrix $\bm{R_0}$ whose rows are representative cycles from each cycle family stored as vectors using a fundamental cycle basis $\calF$. 
%Also recall these cycles are sorted in increasing length.
%
%Cycles in a family differ by smaller cycles, so they can be exchanged freely between minimum cycle bases.
%...
%The cycles in a cycle family can be exchanged freely between minimum cycle bases, since th

%To select a cycle $C_i$ for a minimum cycle basis consider a non-zero vector $S\in\{0,1\}^\nu$ -- which can also be interpreted as a cycle -- 
%and the optimization problem
%\begin{equation}
%	\min_C |C|, \text{ such that }\langle C,S\rangle_\mathcal{F}=1.
%\end{equation}
%Additionally, consider a minimal length cycle $C$ that satisfies the constraint $\langle C,S\rangle=1$. Let the expansion of $C$ with respect to some arbitrary minimum cycle basis $\calB$ is
%Let $C^\star$ be a optimizer to~\eqref{eq:len_const_opt}, and take its expansion with respect to a given minimum cycle basis $\calM$
%\begin{equation}
%	C^\star = B_1\oplus B_2\oplus\hdots \oplus B_k,\quad B_1,B_2,\hdots,B_k\in\calM.
%\end{equation}
%By linearity of the inner product 
%\begin{equation}
%	1 = \langle C^\star,S\rangle = \langle B_1,S\rangle+\langle B_2,S\rangle +\hdots +\langle B_k,S\rangle\quad\text{(mod 2)}
%\end{equation}
%and there is a cycle $C_j$ in the minimum cycle basis such that $\langle C_j,S\rangle = 1$. As an optimizer of~\eqref{eq:len_const_opt}, we achieve the length bound $|C_j|\ge |C^\star|$. Since $\calB$ is a minimum cycle basis, the length of $C$ is bounded by $|C|\ge |C_j|$ by statement 1 of Lemma~\ref{Lem:Relevant_MCB_Test}. So, $C_j$ is also a minimal length cycle subject to $\langle C_j,S\rangle = 1$, where $|C|=|C_j|$. Then, by statement 2 of Lemma~\ref{Lem:Relevant_MCB_Test}, we may exchange $C$ and $C_j$ to obtain a new minimum cycle basis $\calB'=(\calB\setminus\{C_j\})\cup \{C\}$.

%To recap, any minimum cycle basis has some minimizer of~\eqref{eq:len_const_opt}, and any minimizer of~\eqref{eq:len_const_opt} is in a minimum cycle basis. 
%
%
%------\\
%put somewhere
%\begin{equation}
%	\langle C_j,S_k\rangle=0,\quad\text{for }j<k.
%\end{equation}
%\begin{equation}
%	\langle C_j, S_j\rangle = 1.\label{eq:Sel_Const}
%\end{equation}
%
%
%Make a comment: it is critical that we output the witness vectors alongside the minimum cycle basis in Alg.~num. which will be used later in this section.

%------


%With the a candidate relevant cycle families it is pertinent to compute a minimum cycle basis. Specifically, our tests for relevance, trivial-interchangeability, and direct-interchangeability all rely on the decomposition of cycles with respect to an arbitrary minimum cycle basis. See Lemma~\ref{Lem:Relevant_MCB_Test}, Lemma~\ref{Lem:TrEx_Defs} (3), and Theorem~\ref{Thm:DirInt_Comp} respectively. There is an extensive body of research dedicated to the problem of computing minimum cycle bases, see e.g.~\cite{kavithaCycle2009}. Ultimately, we find that it is much more expensive

%--------
%
%Review witness algorithm from Kavitha (2004)?
%
%get roughly $O(\nu^3\log(\nu))$? complexity. Mention this is not the lowest complexity (Kavitha review and Michael and Mehlhorn). Give a remark how in practice this is typically better (sparsity makes the draw faster, $S_i$ sparse and don't need to sort the full lists in the dot products).

\subsection{Change of basis}
\label{sub:MCB_COB}



Next, we describe our procedure to compute the expansion of a cycle $C$ in terms of an MCB $\calM=\{B_1,\hdots,B_\nu\}$ computed by Algorithm~\ref{alg:MCB}.
We assume $C$ is given in vector form corresponding to its expansion with respect to the fundamental cycle basis $\calF$ used to construct $\calM$. 
Let the expansion of $C$ with respect to $\calM$ be given by
%We assume $C$ the minimum cycle $\calM$ is output by Algorithm~\ref{alg:MCB}, and that $C$ is represented as a fundamental cycle basis vector.
%Next, we want to expand cycles in terms of a given minimum cycle basis $\calB=\{C_1,\hdots,C_\nu\}$ given by Algorithm~\ref{alg:MCB}. %We will show that this can be done using the witness vectors $S_i$ used to compute the minimum cycle basis. 
%Let $C$ be an arbitrary cycle, with expansion in $\calB$ given by
\begin{equation}
	C = b_1 B_1\oplus\hdots\oplus b_\nu B_\nu,\quad b_i\in\{0,1\},\label{eq:C_exp_MCB_0}
\end{equation}
where the coefficients $b_i$ are unknown.
We will compute these coefficients using the witness vectors $S_i$.% used to construct $\calM$ in Algorithm~\ref{alg:MCB}.
%Ultimately, we want to compute the coefficients $b_i$. We do so recursively in reverse using the witness vectors $S_i$ used to compute $\calB$.

\begin{algorithm}[t]
	\KwIn{$\bullet$ Matrix of cycles $\bm{R_0}$ from V'-\newline $\phantom{\bullet}$ families as fundamental cycle\newline $\phantom{\bullet}$ basis vectors\newline
		$\bullet$ MCB $\calM=\{B_1{,}\hspace{1pt}.\hspace{1pt}.\hspace{1pt}.\hspace{1pt},B_\nu\}$ %as fundamental cycle basis vectors
		\newline
		$\bullet$ Witness vectors $S_1,\hdots,S_\nu$}
	\KwOut{$\bullet$ Matrix of cycles $\bm{R_1}$ from\newline $\phantom{\bullet}$ V'-families as MCB vectors}
	Initialize %modified witness vectors 
	$\tilde{S}_i\gets S_i$ for $i=1,\hdots,\nu$\;
	\For{$k = \nu,\nu-1,\hdots,2$}{
		\# $\tilde{S}_k$ is known, update $\tilde{S_j}$ for $j<k$\;
		\For{$j = 1,\hdots,k-1$}{
			\uIf{$\langle B_k, S_j\rangle_\calF=1$}{
				$\tilde{S_j}\gets \tilde{S}_j\oplus\tilde{S}_k$\;
			}
		}
	}
	$\bm{W} \gets [\tilde{S}_1,\hdots,\tilde{S}_\nu]$\;
	$\bm{R_1} \gets \bm{R_0} \bm{W}$\;
	\caption{Compute MCB representation}
	\label{alg:COB}
\end{algorithm}

We start with the coefficient $b_\nu$. By construction, $S_\nu$ is orthogonal to every cycle in $\calM$ except $B_\nu$. Thus, %if we take the inner product of Eq.~\eqref{eq:C_exp_MCB_0} with $S_\nu$ and simplify using Eqs.~\eqref{eq:Wit_Orth} and~\eqref{eq:len_const_opt} we find that
\begin{equation}
	\langle C,S_\nu\rangle_\calF = b_\nu.
	\label{eq:bnu}
\end{equation}
%This is the desired formula for $b_\nu$.

For the coefficient $b_{\nu-1}$ we use the witness vector $S_{\nu-1}$ which is orthogonal to $B_1,\hdots,B_{\nu-2}$ by Eq.~\eqref{eq:Wit_Orth}. We take the inner product of Eq.~\eqref{eq:C_exp_MCB_0} with $S_{\nu-1}$ and substitute $\langle C,S_\nu\rangle_\calF$ for $b_\nu$ using Eq.~\eqref{eq:bnu}
\begin{equation}
	\langle C,S_{\nu-1}\rangle_\calF = b_{\nu-1} \oplus \langle C,S_\nu\rangle_\calF \langle B_\nu,S_{\nu-1}\rangle_\calF.
\end{equation}
Then, we rearrange and solve for $b_{\nu-1}$
\begin{equation}
	b_{\nu-1} = \langle C,S_{\nu-1}\rangle_\calF \oplus \langle C,S_\nu\rangle_\calF \langle B_\nu,S_{\nu-1}\rangle_\calF.\label{eq:bnum1}
\end{equation}
We move the inner product $\langle C_\nu,S_{\nu-1}\rangle_\calF$ into the right term of $\langle C,S_\nu\rangle_\calF$ by linearity. After doing so -- and combining similar terms -- we find that Eq.~\eqref{eq:bnum1} simplifies into a single inner product
%More succinctly, $b_{\nu-1}$ is obtained via the inner product
\begin{align}
	b_{\nu-1} &= \langle C, \tilde{S}_{\nu-1}\rangle_\calF,\\ \tilde{S}_{\nu-1} &= S_{\nu-1}\oplus\langle B_\nu,S_{\nu-1}\rangle_\calF S_\nu.
\end{align}
We call $\tilde{S}_{\nu-1}$ a \emph{modified witness vector}.

In general, the coefficients $b_j$ can be computed via the inner product of $C$ with a modified witness $\tilde{S}_j$. This can be verified by induction where Eq.~\eqref{eq:bnu} is the base case. If we assume that $b_k = \langle C,\tilde{S}_k\rangle$ for $k>j$ then the inner product of Eq.~\eqref{eq:C_exp_MCB_0} with $S_j$ is
\begin{equation}
	\langle C,S_j\rangle_\calF = b_j\oplus\bigoplus_{{k=j+1}}^\nu \underbrace{\langle C,\tilde{S}_k\rangle_\calF}_{b_k} \langle B_k,S_j\rangle_\calF .
\end{equation}
Using the linearity of the inner product, we can combine terms to obtain $b_j$ and the associated modified witness $\tilde{S}_j$
%If we were to continue this procedure then we would find that all of the coefficients $b_j$ are given by an inner product with $C$. Specifically, the formula for the coefficients $b_j$ is
\begin{align}
	b_j &= \langle C,\tilde{S}_j\rangle_\calF,\\ \tilde{S}_j&= S_j\oplus\bigoplus_{{k=j+1}}^\nu\langle B_k,S_j\rangle_\calF \tilde{S}_k.\label{eq:mod_wit}
\end{align}
%This equation can be verified by induction by taking the inner product of Eq.~\eqref{eq:C_exp_MCB_0} with $S_j$ and re-arranging terms.



Once the modified witness vectors $\tilde{S}_j$ are computed using Eq.~\eqref{eq:mod_wit}, it is useful to collect them as the columns of a matrix
\begin{equation}
	\bm{W} = [\tilde{S}_1, \tilde{S}_2,\hdots,\tilde{S}_\nu].
\end{equation}
Then if we left-multiply $\bm{W}$ by our cycle $C$ (as a vector expanded in terms of a fundamental cycle basis) we obtain its expansion in terms of $\calB$
\begin{equation}
	C \cdot\bm{W} = [\langle C{,}\hspace{1pt}\tilde{S}_1\rangle_\calF{,}\hspace{1pt}\langle C{,}\hspace{1pt}\tilde{S}_2\rangle_\calF{,}\hspace{1pt}.\hspace{1pt}.\hspace{1pt}.\hspace{1pt}{,}\hspace{1pt}\langle C{,}\hspace{1pt}\tilde{S}_\nu\rangle_\calF].
\end{equation}
Recall, in Section~\ref{sub:Vis_Fam} we computed the matrix $\bm{R}_0$ whose rows are representative cycles from the V'-families as fundamental cycle basis vectors. If we right-multiply this matrix by $\bm{W}$,
\begin{equation}
	\bm{R}_1 = \bm{R}_0 \bm{W},
	\label{eq:R1_comp}
\end{equation}
then we obtain the expansion of these cycles in terms of our MCB $\calM$.


Algorithm~\ref{alg:COB} computes the matrix $\bm{R_1}$. Lines 2-8 compute the modified witness vectors according to Eq.~\eqref{eq:mod_wit}. The inner products and vector sums in these lines are $O(\nu^3)$. In line 10, $\bm{R_1}$ is computed by Eq.~\eqref{eq:R1_comp}. 
%we compute the matrix $\bm{R_1}$ using matrix multiplication as in Eq.~\eqref{eq:R1_comp}. 
Since $\bm{R_0}$ is a $O(\nu^2)\times \nu$ matrix, a simple implementation of this matrix multiplication is $O(\nu^4)$. In practice, Algorithm~\ref{alg:COB} is very efficient due to sparsity.

\subsection{{\tt pi} and {\tt sli} classes}
\label{sub:DICE_TICE_Comp}

%\textbf{Header}
Here we present a method for computing the {\tt pi} and {\tt sli} classes of an arbitrary unweighted, undirected graph. 
%Our pseudocode is given in Algorithm~\ref{alg:DICE_TICE_Comp}. 
In Figure~\ref{fig:DICE_TICE_alg}~(a-d), we show the major steps for this procedure for the graph shown in Figure~\ref{fig:VisFamEx}~(a). 
%To illustrate the steps of this algorithm, we consider the graph shown in Figure~\ref{fig:DICE_TICE_ex} (a). 
%This graph is interesting because it is relatively simple, has heterogeneous cycle lengths, and because the direct- and trivial-interchangeability classes achieve different levels of granularity.
We describe these steps below.


\begin{figure*}[t]
	\centering
	\includegraphics{Matrices.pdf}
	\caption{\textbf{(a)} Input matrix of representative cycles from each V'-family as MCB vectors. The blocks of this matrix separate cycles of differing lengths. \textbf{(b)} 
		%	Reduced matrix without $C_8$ which is not irrelevant cycle as the sum of shorter cycles.
		Reduced matrix without the irrelevant cycle $C_8$ since it is the sum of shorter cycles. 
		The orange arrows depict the procedure for traversing rows and columns to compute the {\tt pi} classes. \textbf{(c)} The same matrix where the rows and columns are sorted by {\tt pi} classes. The block matrices for the {\tt pi} classes are separated by dashed lines. \textbf{(d)} Output matrix $\bm{R}$ with $C_4$ and $C_6$ merged into a single row representing a {\tt sli} class, since $C_4=C_6\oplus\smCyc$. \textbf{(e)} The graph (left) and its relevant cycles separated into {\tt sli} classes (right). The representative and secondary cycles in the {\tt sli} classes are drawn using solid and dashed edges. The gray box highlights the large {\tt pi} class $\{C_3,C_4,C_6,C_7,C_7'\}$ with multiple {\tt sli} classes.
	}
	\label{fig:DICE_TICE_alg}
\end{figure*}

\textbf{Inputs.} The inputs for our algorithm are (i) the V'-families output by Algorithm~\ref{alg:vis_fam} and (ii) the matrix $\bm{R_1}$ of representative cycles from each V'-family as MCB vectors from Algorithm~\ref{alg:COB}. 
We sort the rows and columns of $\bm{R_1}$ by increasing cycle length. %which has time complexity $O(\nu^2\log\nu)$ where $\bm{R_1}$ is of size $O(\nu^2)\times\nu$.
We utilize the block structure of this matrix, where the row and column cycles have fixed length. See Figure~\ref{fig:DICE_TICE_alg} (a). 
This block matrix is lower triangular because the MCB representation of a cycle only includes equal or shorter length cycles -- see Lemma~\ref{Lem:Relevant_MCB_Test}.

\textbf{Remove irrelevant cycles.} By Lemma~\ref{Lem:Relevant_MCB_Test}, Statement 2, a cycle is relevant if and only if its MCB representation includes an equal length cycle. In Figure~\ref{fig:DICE_TICE_alg}~(a), the eight-sided cycle $C_8$ is not relevant as the sum of two six-sided cycles $C_4,C_5$ and a four-sided cycle $C_2$. Algorithmically, we loop through each row to find the largest cycle in each MCB vector to check if the row is relevant and remove it otherwise. 
%This has time complexity $O(\text{nnz}(\bm{R_1}))$ where $\text{nnz}(\bm{R_1})$ is the number of nonzero entries in $\bm{R_1}$. 
%We denote this trimmed matrix of \emph{relevant} cycles as $\bm{R_2}$

\textbf{{\tt pi} classes.} 
From Lemma~\ref{Lem:dir_intchg_2}, a relevant cycle is {\tt pi}-interchangeable for the equal length cycles in its MCB representation. These relations are identified by the non-zero entries in the diagonal blocks of $\bm{R_2}$, e.g., $C_7$ is {\tt pi}-interchangeable for $C_3$ and $C_4$ in Figure~\ref{fig:DICE_TICE_alg}~(b).
%, e.g., $C_7$ is directly-interchangeable for $C_3$ and $C_4$ in Figure~\ref{fig:DICE_TICE_alg} (b).
From Theorem~\ref{Thm:DirInt_Comp}, the {\tt pi} classes are computed using the transitive closure of these relations.
This is done by traversing the rows and columns in a diagonal block of $\bm{R_2}$ using a queue -- see the orange arrows in Figure~\ref{fig:DICE_TICE_alg}~(b). %This has time-complexity $O(\text{nnz}(\bm{R_2}))$.

%It is sufficient to use the representative cycles from each family since the cycles in a family differ by shorter cycles only.

We permute the rows and columns of $\bm{R_2}$ to group the {\tt pi} classes together.
We also permute rows within a {\tt pi} class by
\begin{enumerate}
	\item number of equal length cycles in the MCB vector, then by
	\item lexicographic ordering of the MCB vectors read \emph{right-to-left}.
\end{enumerate}
In Figure~\ref{fig:DICE_TICE_alg}~(c), the cycle $C_7$ follows $C_6$ because it has one more equal length cycle in its MCB representation. The cycle $C_6$ follows $C_4$ since $[0,1,0,1,0,0]$ follows $[0,1,0,0,0,0]$ in lexocographic ordering. We denote this sorted matrix by $\bm{R_3}$. 
%For instance, $C_6$ follows $C_4$ in $[0,1,0,1,0,0]$ follows $[0,1,0,0,0,0]$ in lexocographic ordering. Figure~\ref{fig:DICE_TICE_alg}~(c) because $[0,1,0,1,0,0]$ follows $[0,1,0,0,0,0]$ in lexocographic ordering. Also, $C_7$ follows $C_6$ because it has two equal length cycles in its minimum cycle basis representation instead of 1. We denote this sorted matrix by $\bm{R_3}$. Using dense operations, this has complexity $O(\nu^3\log\nu)$ where $O(\nu)$ is the time complexity to compare two rows of length $\nu$. In practice, this step is extremely fast.

\textbf{{\tt sli} classes.} 
Two cycles belong to the same {\tt sli} class if their MCB vectors are equal up to shorter cycles by Lemma~\ref{Lem:TrEx_Alg}. By using reverse lexicographic ordering, cycles in a {\tt sli} class will appear as consecutive rows of $\bm{R_3}$. See the cycles $C_4$ and $C_6$ in Figure~\ref{fig:DICE_TICE_alg}~(c). We merge cycles in a {\tt sli} class into a single row, and we keep track of the relevant V'-families they represent. We denote this final matrix as $\bm{R}$.
%By Lemma~\ref{Lem:TrEx_Alg}, two cycles are trivially-interchangeable if there minimum cycle basis vectors are equal up to shorter cycles. For instance,

\textbf{Time-Complexity.} Most steps of this procedure can be done using a single pass through the working matrix $\bm{R_i}$. This has time-complexity $O(\text{nnz}(\bm{R_i}))=O(\text{nnz}(\bm{R_1}))$ where $\text{nnz}(\bm{X})$ is the number of nonzero entries in $\bm{X}$. 
The largest exception to this is sorting the rows of each {\tt pi} class using reverse lexicographic ordering. 
Using dense operations, this has $O(\nu^3\log\nu)$ time-complexity, where $\bm{R_2}$ has $O(\nu^2)$ rows and comparing two rows takes $O(\nu)$ time. Using sparse operations, this has $O(\text{nnz}(\bm{R_1})\log\nu)$ time-complexity, where the time to compare rows reduces to the average number of nonzero entries per row.
For our MD simulations, this procedure is an order of magnitude faster than computing the relevant cycle families and an MCB.
%In MD data, the number of non-zero entries in the matrix 
%The matrix $\bm{R_2}$ has $O(\nu^2)$ rows, and using dense operations comparing two rows of $\bm{R_2}$ takes $O(\nu)$ time. This suggests $O(\nu^3\log\nu)$ time-complexity. In our simulations, the number of rows in $\bm{R_1}$ is very close to $\nu$, and there is a limited number of entries per row. For a simulation of $\rm C_{10}H_{16}$, this procedure takes significantly less time than the time needed to compute the relevant cycles families and 

\textbf{Output.} The output relevant cycle matrix $\bm{R}$ can be interpreted as follows. Each row of this matrix is a representative cycle for a {\tt sli} class as an MCB vector. The choice of representative cycle only changes these MCB vectors by shorter cycles. Moreover, the entries in the diagonal blocks of $\bm{R}$ do not depend on the choice of representative cycles.

The {\tt pi} classes form diagonal blocks in $\bm{R}$. The rows in these blocks are sorted so that the cycles belonging to the MCB appear first. Because of this, the format of each block is the identity matrix followed by additional rows $[I,X']'$. The polyhedra are obtained via these additional rows $X$. The number of polyhedra obtained is equal to the dimension of the (left) null space of $\bm{R}$, given by the number of rows minus the number of columns in $\bm{R}$. 

For instance, from the row corresponding to $C_7$ in Figure~\ref{fig:DICE_TICE_alg}~(c), we obtain the polyhedron
\begin{equation}
	C_7 = C_3\oplus C_4\oplus\smCyc.
\end{equation}
Here, $\{C_3,C_4,C_7\}$ are the large faces of the triangular prism. The smaller cycles used to construct the polyhedron depend on which representative cycles are sampled from each family. For the cycles given, the full polyhedron is $\{C_0,C_1,C_3,C_4,C_7\}$.
{Codes for computing the decomposition of this graph are given in the notebook {\tt cycle\_decomposition.ipynb} on Github~\cite{ruthCyclic2025}.}

If a {\tt pi} class has multiple polyhedra, then the polyhedra sampled depend on the MCB used to represent the columns of $\bm{R}$. For instance, for the graph shown in Figure~\ref{fig:polyhedron_pair}, this procedure may identify either two cubes or a cube and a large rectangular prism. 
Intuitively, the most desirable polyhedra are the smallest ones, i.e., the cubes in Figure~\ref{fig:polyhedron_pair}.
We leave the challenge of computing smallest polyhedra for future work.



\section{Theory and algorithms for sampling minimum cycle bases}
\label{sec:sampling}

\subsection{Uniform random sampling of minimum cycle bases}
\label{sub:rand_MCB}

%Recall, the primary issue with sampling cycles from a minimum cycle basis is that they are not unique. Furthermore, sample statistics from a minimum cycle basis are not unique, e.g., the number of edges each pair of cycles in the basis share. Of course, this is an issue for reproducibility. To address this issue, we introduce a method for sampling the minimum cycle bases of a graph uniformly at random. Thus, while our measured minimum cycle basis is not unique, it is statistically well-defined as a random variable.

\begin{figure}[t]
	\centering
	\includegraphics[]{TwistaneGraphic.pdf}
	\caption{\textbf{(a)}~3D print of the carbon skeleton of twistane, a molecule with formula $\rm C_{10}H_{16}$. \textbf{(b)}~The carbon skeleton of twistane as a graph and its smallest polyhedra $\calP_1=\{C_2,C_3,C_3'\}$ and $\calP_2=\{C_1,C_1',C_2\}$. %The set $\calP_3=\calP_1\symdiff\calP_2 = \{C_1,C_1',C_3,C_3'\}$ is the third (non-empty) polyhedron. 
	\textbf{(c)}~Transition from the MCB $\calM=\{C_1,C_2,C_3\}$ to $\calM'=\{C_1',C_2,C_3\}$ by swapping $C_1'$ for $C_1$ using the polyhedron $\calP_1$. 
%	Transition from the minimum cycle basis $\calM=\{C_1,C_2,C_3\}$ by swapping $C_1'$ for $C_1$ using the polyhedron $\calP_1$ to the new minimum cycle basis $\calM'=\{C_1',C_1,C_3\}$.
	}
	\label{fig:randMCB}
\end{figure}

Here, we introduce a method for sampling MCBs uniformly at random using a Markov Chain Monte Carlo approach.
We assume the reader is familiar with the basic theory of Markov chains~\cite{durrettEssentials2012}. 
%The random step for our procedure uses the generalized polyhedra obtained from the direct-interchangeability relation, defined in Section~\ref{sub:Poly}. 
%This is done efficiently using the trivial-interchangeability classes.
%First, we outline a simplified procedure without the trivial-interchangeability classes. %this approach without trivial-interchangeability classes first. 

\subsubsection{The basic algorithm}
A given MCB $\calM$ is randomly modified as follows.
\begin{enumerate}
	\item Randomly select a cycle $C'\in\RelCyc\setminus\calM$ to exchange with a cycle in $\calM$.
	\item Construct a polyhedron using $C'$ and $\calM$, i.e., $\calP = \{C'\}\cup\expand{\calM}{C'}$.
	\item Randomly select a cycle $C{\in}\calP$, $|C|{=}|C'|$, and return %the modified minimum cycle basis
	$\calM'=(\calM\setminus\{C\})\cup\{C'\}$. 
	If $C=C'$, return the same basis $\calM=\calM'$.
\end{enumerate} 
See Figure~\ref{fig:randMCB}~(c) for an example of this procedure in a graph with multiple polyhedra. These steps are repeated towards equilibrium.

To validate this approach, we make the following observation.
\begin{prop}
	The above Markov chain is ergodic.
\end{prop}
\begin{proof}
	First, we need the above procedure to be aperiodic. This is true because there is nonzero holding probability, i.e., the probability we return to the same state $P(\calM'=\calM)$. These holding probabilities cause the MCB to leave its current state at a random (geometric) time, which breaks any otherwise periodic behavior in the system.
	
	Second, the system must be irreducible, i.e., any MCB must be reachable by any other. Consider two distinct MCBs $\calM_1$ and $\calM_2$. By Lemma~\ref{Lem:DICE_Irreducible}, we can exchange a cycle in $\calM_1$ for another in $\calM_2$ to obtain a new MCB $\calM_1{\to}\calM_1'$ that shares an additional cycle with $\calM_2$. We may repeat this procedure to obtain $\calM_2$ from $\calM_1$.
\end{proof}


%Recall, a Markov chain is ergodic if it is aperiodic and irreducible. This procedure is aperiodic because it has non-zero holding probabilities, i.e., the probability we draw the same cycle $C=C'$ in step 3 above.
%This procedure describes an {ergodic} Markov chain because (i) it has non-zero holding probabilities and (ii) it is irreducible. 
%The positive holding probabilities are given by the probability that the cycle $C=C'$ is selected from $\calP$.
%To show irreducibility, consider distinct minimum cycle bases $\calM_1$ and $\calM_2$. Also, consider a cycle $C_2\in\calM_2\setminus\calM_1$ representing the cycle $C'$ in the procedure above.
%We make the following observations.
%The cycles of length $|C_2|$ in $\calM_2$ are mutually relevant by Definition~\ref{def:mut_rel}, so they cannot be used to sum to shorter cycles by Corollary~\ref{Cor:mut_rel_2}. In contrast, the faces of the polyhedron $\calP=\{C_2\}\cup\expand{\calM_1}{C_2}$ have null sum, so the cycles of length $|C_2|$ sum to shorter cycles:
%First, the cycles of length $|C_2|$ in $\calM_2$ are mutually relevant by Definition~\ref{def:mut_rel}.
%Next, the cycles of length $|C_2|$ in $\calP=\{C_2\}\cup\expand{\calM_1}{C_2}$
%To see that the Markov chain is irreducible, consider distinct minimum cycle bases $\calM_1$ and $\calM_2$ and a cycle $C_2\in\calM_2\setminus\calM_1$.
%We make the following observations about mutually relevant sets.
%By Definition~\ref{def:mut_rel}, the cycles of length $|C_2|$ belonging to $\calM_2$ are mutually relevant.
%The polyhedron $\calP=\{C_2\}\cup\expand{\calM_1}{C_2}$ is constructed to have null sum. Moreover, the cycles of length $|C_2|$ in $\calP$ are not mutually relevant by Corollary~\ref{Cor:mut_rel_2} because they sum to shorter cycles:
%\begin{equation}
%	\bigoplus_{C\in\calP}C=\emptyset \quad\to\quad \bigoplus_{\mathclap{\substack{C\in\calP\\|C|=|C_2|}}}C = \bigoplus_{\mathclap{\substack{C'\in\calP\\|C'|<|C_2|}}}C'.
%\end{equation}
%Cycles in M2 are mutually relevant... Cycles in P are not. 
%Therefore, the cycles of length $|C_2|$ in $\calP$ are not a subset of $\calM_1$, and there exists a cycle $C_1\in\calP$ of length $|C_2|$ such that $C_1\in\calM_1\setminus \calM_2$. 
%Therefore, there exists a cycle $C_1\in\calP$ of length $|C_2|$ belonging to $C_1\in\calM_1\setminus \calM_2$. 
%We may exchange $C_1$ for $C_2$ to obtain a minimum cycle basis $\calM_1'\,{=}\,(\calM_1\setminus\{C_1\})\cup\{C_2\}$ that shares one more cycle with $\calM_2$ than $\calM_1$. By inductive reasoning, we can reach $\calM_2$ from $\calM_1$ with these exchanges, and our Markov chain is irreducible. 


%{\color{red} Make this a proposition, and make it cleaner so it has aperiodic and irreducible. And make the aperiodic part more clear...}

By ergodicity, this Markov chain approaches a stationary distribution over the set of MCBs. This stationary distribution is the uniform distribution because the Markov chain has symmetric transition probabilities
\begin{equation}
	P(\calM{\to}\calM'){=}P(\calM'{\to}\calM){=}\frac{1}{|\RelCyc{\setminus}\calM|}\frac{1}{|\calP|}.
\end{equation}
The first term $1/|\RelCyc{\setminus}\calM|=1/(|\RelCyc|{-}\nu)$ is the constant probability with which we choose a given cycle $C'\in\RelCyc{\setminus}\calM$ in step 1 of the procedure above. 
The second term is the probability we select the cycle $C\in\calP$ to exchange for $C'$. 
%we exchange $C'$ for the cycle $C\in\calP$. 
The same polyhedron $\calP$ is used in both directions $\calM\to\calM'$ and $\calM'\to\calM$.

\begin{algorithm}[t]
	\KwIn{$\calM=$ input MCB\newline
		$\bm{R}=$ representative cycles for\newline $\phantom{\bm{R}=}$ {\tt sli} classes as vectors in $\calM$} %minimum cycle basis\newline $\phantom{\bullet}$  ($\calM$) vectors}
	\KwOut{$\calM_r=$ random MCB}
	\For{$i=1,\hdots,N_{\rm steps}$}{
		Randomly sample $C'\notin\calM$ from $\bm{R}$\;
		Using the row for $C'$ in $\bm{R}$ construct the polyhedron $\calP=\{C'\}\cup\expand{\calM}{C'}$\;
		Sample $C\in\calP, |C|=|C'|,$ with probability proportional to $1/|\calS(C)|$\;
		\uIf{$C\neq C'$}{
			$\calM\gets(\calM\setminus\{C\})\cup\{C'\}$\;
			\# Update $\bm{R}$ to new basis\;
			\For{Row cycle $C_i=\bm{R}[i,:]$}{
				\uIf{$C\in\expand{\calM}{C_i}$}{
%					\# Update basis representation of $C$ from $\calM$ to $\calM'$\;
					Set $\bm{R}[i,:]$ to $\expandd{\calM'}{C_i}=$  
					$ (\expandd{\calM}{C_i}\setminus\{C\})\symdiff\expandd{\calM'}{C}$\;
				}
			}
		}
	}
	
	$\calM_r\gets [\ ]$ \# empty list\;
	\For{$C\in\calM$}{
		Randomly sample a cycle from the {\tt sli} class $\calS(C)$ and add it to $\calM_r$\;%with probability proportional to its size\;
	}
	\caption{Random MCB}
	\label{alg:rand_MCB}
\end{algorithm}

%We simplify this procedure by working with a set of representative cycles from each trivial-interchangeability class. We store these cycles as a matrix $\bm{R}$ of minimum cycle basis vectors, computed in Section~\ref{sub:DICE_TICE_Comp}. These cycles are used to randomly select a valid set of trivial-interchangeability classes from which we draw our random minimum cycle basis. Recall, from Lemma~\ref{Lem:TrInt_Eq_Rel} 

%{\color{red}
%Modified exchange algorithm

\subsubsection{The efficient algorithm}

We simplify this procedure by drawing cycles from the matrix $\bm{R}$ computed in Section~\ref{sub:DICE_TICE_Comp}. %-- see Algorithm~\ref{alg:rand_MCB} for the pseudocode. 
Each cycle in $\bm{R}$ represents a {\tt sli} class. 
The MCB we will sample from $\bm{R}$ is used to select a set of {\tt sli} classes to draw cycles from. 
%represents the set of minimum cycle bases that share its trivial-interchangeability classes. 
%The Markov chain approach using polyhedra is used to randomize the trivial-interchangeability classes from which we draw a minimum cycle basis.
It is valuable to sample cycles from $\bm{R}$ since there is only a polynomial number of {\tt sli} classes. Cycles can be efficiently sampled from a {\tt sli} class directly, which we discuss at the end of this subsection. 

%Note, a minimum cycle basis drawn from this procedure represents

%The cycles we draw from $\bm{R}$ represent trivial-interchangeability classes. 
By Corollary~\ref{Cor:TrInt_Exchanges}, we can exchange pairs of cycles within a {\tt sli} class arbitrarily between MCBs. 
Thus, the total number of MCBs we can obtain using valid {\tt sli} classes $\calS_1,\hdots,\calS_\nu$ is $|\calS_1|\cdots|\calS_\nu|$. 
We want the stationary distribution to be proportional to the number of MCBs that $\calM$ drawn from $\bm{R}$ represents,
%We want the stationary distribution that MCB $\calM$ is drawn from $\bm{R}$ to be proportional to the number of MCBs it represents, i.e.,
\begin{equation}
	\pi_\calM \propto \prod_{C\in\calM} |\calS(C)|
	\label{eq:rand_stat_dist}
\end{equation}
where $\calS(C)$ is the {\tt sli} class containing $C$.
%}

To achieve the stationary distribution given by Eq.~\eqref{eq:rand_stat_dist}, we bias the transition probabilities to satisfy the detailed balance condition
\begin{equation}
	P(\calM\to\calM')\pi_\calM = P(\calM'\to\calM)\pi_{\calM'}.\label{eq:rand_det_bal}
\end{equation}
We re-arrange Eq.~\eqref{eq:rand_det_bal} to obtain the desired ratio between the forward and backward transition probabilities
\begin{equation}
	\frac{P(\calM\to\calM')}{P(\calM'\to\calM)} = \frac{\pi_{\calM'}}{\pi_\calM} = \frac{|\calS(C')|}{|\calS(C)|}
\end{equation}
where $\calM'=(\calM\setminus\{C\})\cup\{C'\}$. This is achieved by removing the cycle $C$ from our polyhedron with probability proportional to $1/|\calS(C)|$, as in Algorithm~\ref{alg:rand_MCB}, line 4. Then, the transition probabilities are given by
\begin{equation}
	\begin{split}
	&P(\calM\to\calM') \\&\quad= \frac{1}{\text{dim}(\text{coker}(\bm{R}))}\frac{1/|\calS(C)|}{\sum_{C'\in\calP}1/|\calS(C')|}
	\end{split}\label{eq:rand_trans_prob}
\end{equation}
where $\text{dim}(\text{coker}(\bm{R}))$ is the dimension of the left null space of $\bm{R}$, also the number of row cycles in $\bm{R}$ that do not belong to $\calM$.
Only the numerator $1/|\calS(C)|$ changes in Eq.~\eqref{eq:rand_trans_prob} between the forward ($\calM{\to}\calM'$) and reverse ($\calM'{\to}\calM$) exchange. The same polyhedron $\calP$ is used in both directions.

%{\color{red} Update the mimimum cycle basis vectors}

After performing the exchange $\calM\to\calM'$ we update the rows of $\bm{R}$ to be vectors in $\calM'$. We only need to update the representation of a cycle $C_i=\bm{R}[i,:]$ if its expansion $\expandd{\calM}{C_i}$ includes the cycle $C$ swapped for $C'$. If $C\in\expandd{\calM}{C_i}$, we replace $C$ with its expansion in $\calM'$
\begin{equation}
	\expandd{\calM'}{C_i} = (\expandd{\calM}{C_i}\setminus\{C\})\symdiff\expandd{\calM'}{C}.
\end{equation}


\begin{figure}[t]
	\centering
	\includegraphics{branch_rec.pdf}
	\caption{
		Procedure used to randomly sample a cycle from an even cycle family. 
		A cycle is sampled by joining two random paths $p\to u$ and $q\to v$ to the descriptors $(p,q)$ and $(u,v)$. 
		Node labels represent the number of valid shortest paths from a node to $(u,v)$. 
		Edge labels represent the probability a node is selected to construct a shortest path.
}
	\label{fig:rand_fam_backtrack}
\end{figure}

Last, we describe how cycles are sampled from each {\tt sli} class. A {\tt sli} class is the union of disjoint V'-families output by Algorithm~\ref{alg:vis_fam}. 
We choose a V'-family with probability proportional to the number of cycles it has. 
A cycle is sampled from the V'-family by backtracking from either the node $p$ or edge $(p,q)$ to the root edge $(u,v)$. 
At each step, the probability we step through a given node is proportional to the number of paths it has -- see Figure~\ref{fig:rand_fam_backtrack}. The final MCB is obtained by sampling a random cycle from each {\tt sli} class. The whole procedure is summarized in Algorithm~\ref{alg:rand_MCB}.

%As an example, consider the molecule twistane depicted in Figure~\ref{fig:twistane}.
%Starting with the minimum cycle basis $\calM=\{C_1,C_2,C_3\}$ we choose a random cycle outside of $\calM$, say $C_1'$. Then, we randomly select a cycle from the polyhedron $\calP_1 = \{C_1,C_1',C_2\}$, say $C_2$. Then, we obtain a new minimum cycle basis $\calM'=\{C_1,C_1',C_3\}$. This procedure is iterated until equilibrium.

%Using this procedure we can obtain any minimum cycle basis. Starting with the basis $\calM_1$ we can approach the new basis $\calM_2$ by exchanging a cycle $C_2\in\calM_2\setminus\calM_1$ for .

%In our codes, we use the trivial-interchangeability classes in place of cycles since the number of relevant cycles may grow exponentially. 

%As an example, consider the molecule twistane depicted in Figure~\ref{fig:twistane} composed of two polyhedra $\calP_1=\{C_1,C_1',C_2\}$ and $\calP_2=\{C_2,C_3,C_3'\}$. 
%Starting with an arbitrary minimum cycle basis such as $\calM=\{C_1,C_2,C_3\}$, we draw a cycle outside

%We randomly sample minimum cycle bases we use a Markov chain Monte Carlo (MCMC) method relying on the abstract polyhedra introduced in Section~\ref{Sec:Theory}. To introduce this procedure consider the graph depicted in Figure~\ref{fig:twistane} and the minimum cycle bases $\calM = \{C_0,C_1,C_2\}$. To randomize this basis, we repeatedly draw cycles...

%To generalize this procedure to arbitrary graphs we replace cycles with trivial-interchangeability classes.


\subsection{Pairwise intersections in a minimum cycle basis}
\label{sub:pair_intersect}

\begin{figure}[t]
	\centering
	\includegraphics[trim={0 .6cm 0 .6cm},clip]{DualFormat.pdf}
	\caption{Ring cluster and its dual graph from a $\rm C_8H_{18}$ simulation at 4000K and 40.5GPa. Blue nodes and edges denote carbon atoms and bonds that belong to the cluster. Orange nodes represent carbon rings in the dual graph. Orange edges represent a pair of rings that are joined by a path. The weights of the orange nodes and edges are equal to the number of edges in the corresponding ring or path.}
	\label{fig:dual_graph}
\end{figure}


The overarching goal of this manuscript is to develop mathematical theory that is capable of describing ring clusters in hydrocarbon pyrolysis. So far, we have discussed how to obtain the relevant cycles and how to partition these cycles into a polynomial number of {\tt pi} and {\tt sli} classes. These serve as the building blocks of our ring clusters. In this section, we present new theory for how these cycles intersect. 
This consists of two main results, the proofs of which are contained in Appendix~\ref{app:intersect_proofs}
%In this subsection, we introduce two new results: 

\begin{enumerate}
	\item We characterize the general format for how pairs of cycles in an MCB may intersect. See Theorem~\ref{Thm:MCB_Intersect}. 
	\item We introduce a method for modifying a MCB $\calM\to\calM'$ such that pairs of cycles in the new MCB $C_i,C_j\in\calM'$ intersect over a single path $P_{i,j}=C_i\cap C_j$. See Theorem~\ref{Thm:MCB_Pair_Swap}. This makes it possible to represent the loop cluster as a graph of cycles, also referred to as a \emph{dual graph}. The nodes of the dual graph are the cycles $\calM'=\{C_1,\hdots,C_\nu\}$, and the edges are the non-empty paths of intersections $\{P_{i,j}\}$. See Figure~\ref{fig:dual_graph}. %\\{\color{red} This proof is not complete! I am stuck on an inductive argument. Figure is not made yet.}
\end{enumerate}

%First, let $C\in\RelCyc$ be a relevant cycle. Recall, that a path belonging to $C$, $P\subseteq C$, of length $|P|\le|C|/2$ is a shortest path~\cite{vismaraUnion1997}. See Lemma~\ref{lem:shortest_paths}. 
%To start, recall that for a relevant cycle $C\in\RelCyc$, a path $P\subseteq C$ of length $|P|\le|C|/2$ is a shortest path~\cite{vismaraUnion1997}. 
%We reformulate this statement in terms of graph distances, where the distance $d(u,v)$ of nodes $u,v$ in a graph $G$ is the length of a shortest path connecting them.


\begin{figure}
	\centering
	\includegraphics{Intersect_Thm.pdf}
	\caption{%Depiction of the format in which pairs of cycles in a minimum cycle basis may overlap by Theorem~\ref{Thm:MCB_Intersect}. 
		\textbf{(a)} Two cycles $C_1,C_2$ that intersect over three paths. Blue nodes belong to $C_1$, red nodes to $C_2$, and purple nodes to both.
		If we traverse $C_1$ clockwise and $C_2$ counterclockwise, then the intersection paths appear in the same order $P_1\,{\to}\,P_2\,{\to}\,P_3$ and with the same orientation.
%		We traverse $C_1$ and $C_2$ starting at the top moving clockwise and counterclockwise respectively. Then, the intersection paths are traversed in the same order $P_1\to P_2\to P_3$ and with the same orientation. 
		The shorter separating paths have equal lengths $\big|Q_1^{(1)}\big|=\big|Q_1^{(2)}\big|$ and $\big|Q_2^{(1)}\big|=\big|Q_2^{(2)}\big|$. \textbf{(b)}~A cycle $C_1'$ that can be exchanged for $C_1$ such that $C_1'$ and $C_2$ intersect over a single path. $C_1'$ is constructed from $C_1$ by swapping the paths $Q_1^{(1)}$ for $Q_1^{(2)}$ and $Q_2^{(1)}$ for $Q_2^{(2)}$}
		\label{fig:Intersect_Thm}
\end{figure}





%Next, we use Lemma~\ref{Lem:cyc_embed} to describe the format in which pairs of cycles in a minimum cycle basis can overlap. In particular, two cycles in a minimum cycle basis can only overlap over a series of paths with common orientation. We direct the reader to Figure~\ref{fig:cyc_pair_overlap} (b) for a visualization of this result. We emphasize, this result does not generalize to arbitrary pairs of cycles. For instance, the (non-relevant) cycles $C_1,C_2$ in Figure~\ref{fig:cyc_pair_overlap} (f) do not have a common orientation, so they appear twisted.




\begin{theorem}
	Let $\calM$ be an MCB and $C_1,C_2\in\calM$, $|C_1|\le|C_2|$, such that $C_1$ and $C_2$ share at least one vertex.
	%	Let $C_1,C_2$ be mutually relevant cycles that share nodes, $|C_1|\le |C_2|$. % and $P_1, P_2, \hdots, P_k$ the paths on which they intersect.
	There exist circulations of $C_1$ and $C_2$ such that: 
	\begin{enumerate}
		\item $C_1$ and $C_2$ intersect over $k$ paths $P_1,P_2,\ldots,P_k$ that appear with the same order and orientation in both cycles. 
		
		\item Let $Q_{i}^{(j)}$ for $i\,{=}\,1{,}...{,}\,k$, $j\,{=}\,1,2,$ be the path in $C_j$ joining $P_i$ to $P_{i+1}$, where $P_{k+1}\,{\coloneq}\, P_1$. These paths are equal in length between $C_1$ and $C_2$, possibly except for the largest path pair which is ordered last:
		\begin{align}
			|Q^{(1)}_i| &\!=\! |Q^{(2)}_i|, &i=1,\!...,k-1,
			\label{eq:sm_path_equal}\\
			|Q^{(j)}_i| &\!\le\! |C_1|/2,\! &i{=}1,\!...,\!k{-}1, j{=}1,\!2,
			\label{eq:sm_path_len}\\
			|Q^{(2)}_k| &\!=\! \mathrlap{|Q^{(1)}_k|+|C_2|-|C_1|,}
			\label{eq:lg_path_equal}\\
			|Q^{(j)}_k| &\!\ge\! |C_j|/2,\! &j=1,2.
			\label{eq:lg_path_len}
		\end{align}
	\end{enumerate}
	\label{Thm:MCB_Intersect}
\end{theorem}

\begin{figure}[t]
	\centering
	\includegraphics{Complex_Intersections.pdf}
	\caption{Two pairs of cycles that do not satisfy the format in Theorem~\ref{Thm:MCB_Intersect}. In both cases, the left cycle $C_1$ is not relevant because it is the sum of shorter cycles. 
	\textbf{(a)} Two cycles $C_1,C_2$ that intersect over two paths. 
	The red and blue arrows depict an example pair of orientations of $C_1$ and $C_2$. 
	There are no orientations of $C_1$ and $C_2$ where the intersection paths have the same orientation in $C_1$ and $C_2$. 
	Hence, $C_1$ appears twisted next to $C_2$.
%		(a) A pair of cycles $C_1,C_2$ that intersect over paths $P_1$ and $P_2$ that do not appear in the same orientation in $C_1$ and $C_2$. Hence, $C_1$ appears twisted next to $C_2$.
	\textbf{(b)}~Two cycles $C_1,C_2$ that are separated by shorter paths $Q_1^{(1)}$ and $Q_1^{(2)}$ with different lengths $|Q_1^{(1)}|\ne |Q_1^{(2)}|$.}
	\label{fig:non-rel_overlap}
\end{figure}

In summary, Theorem~\ref{Thm:MCB_Intersect} states that a pair of cycles belonging to some MCB may intersect over a long path broken up by shorter cycles. 
See Figure~\ref{fig:Intersect_Thm} (a). 
In general, we may define the intersection of two cycles using a series of paths, which can be stored as a list of lists. 
We emphasize, that Theorem~\ref{Thm:MCB_Intersect} does not apply to arbitrary pairs of cycles. 
See Figure~\ref{fig:non-rel_overlap}.

Our theory allows us to go further. Specifically, we can stitch together a pair of cycles in an MCB so that they intersect over a single path.


%In summary, Theorem~\ref{Thm:MCB_Intersect} states that a pair of cycles $C_1,C_2$ in a minimum cycle basis are connected by a series paths separated by short cycles. These short cycles are evenly split between $C_1$ and $C_2$. See Figure~\ref{fig:cyc_pair_overlap}~(b). If we add these shorter cycles to $C_1$, then we obtain a new cycle $C_1'$ that intersects with $C_2$ over a single path. We also show $C_1'$ can be exchanged for $C_1$ to obtain a new minimum cycle basis.% the theory for direct- and trivial-interchangeability to show that $C_1'$ can be exchanged for $C_1$ to obtain a new minimum cycle basis.

%{\color{red} Refer to format in which the cycles intersect Figure~\ref{fig:cyc_pair_overlap}~(b). describe it in more detail as a series of paths separated by small cycles evenly split between $C_1$ and $C_2$.
%
%Discuss what adding these small cycles does. Makes them intersect over a single path as in Figure~\ref{fig:cyc_pair_overlap}~(c).}

\begin{theorem}
	Let $\calM$ be an MCB and let cycles $C_1,C_2\in\calM$ share at least two nodes. There exists a cycle $C_1'$ such that 
	\begin{enumerate}
		\item $\calM'=(\calM\setminus\{C_1\})\cup\{C_1'\}$ is an MCB and,
		\item $C_1'$ and $C_2$ intersect over a single path.
	\end{enumerate}
	\label{Thm:MCB_Pair_Swap}
\end{theorem}

\noindent The cycle $C_1'$ is obtained by exchanging paths in $C_1$ for $C_2$. See Figure~\ref{fig:Intersect_Thm}~(b) for an example of this procedure.

It is desirable that $C_1$ and $C_2$ intersect over a single path because a path can be stored as a single list. We propose to postprocess an MCB so that all pairs of cycles share at most one path, as in Algorithm~\ref{alg:merge_MCB}. Then, the cluster of cycles may be characterized using a set of cycles and the paths that connect them. This is the dual graph representation shown in Figure~\ref{fig:dual_graph}. We remark that the dual graph as shown in Figure~\ref{fig:dual_graph} is not lossless because it only includes path lengths, not the positions of these paths in the cycles.



%{\color{red} Think about a disclaimer here}

\begin{algorithm}[t]
	\KwIn{MCB $\calM$}
	\KwOut{MCB $\calM'$, where all pairs $C_1,C_2\in\calM'$ intersect over a single path}
	Compute the intersection paths for all pairs of cycles $C_i,C_j\in\calM$\;
	\For{$i=1,\hdots,N_{\max}$}{
		Randomly sample a pair $C_i,C_j\in\calM$ that intersect over multiple paths\;
		Exchange $\calM\gets(\calM\setminus\{C_i\})\cup\{C_i'\}$ by Theorem~\ref{Thm:MCB_Pair_Swap} where $C_i',C_j$ intersect over a single path\;
		Recompute intersection paths between $C_i'$ and other cycles\;
	}
	\caption{Postprocess MCB}
	\label{alg:merge_MCB}
\end{algorithm}

At this point, we cannot guarantee that Algorithm~\ref{alg:merge_MCB} always outputs an MCB with every pair of cycles intersecting at most over a single path. It always succeeds on our hydrocarbon data in 10-20 iterations. We leave the investigation of this issue for future work.

%We end with a disclaimer. First, consider $C_1,C_2,C_3$ in a minimum cycle basis such that $C_1$ and $C_3$ intersect over a single path while $C_1$ and $C_2$ do not. Then, we use Theorem~\ref{Thm:MCB_Pair_Swap} to construct $C_1'$ such that $C_1'$ and $C_2$ intersect over a single path. However, it may occur that $C_1'$ and $C_3$ do not intersect over a single path. 

%Thus, it is unclear if Algorithm~\ref{alg:merge_MCB} will always be able to compute a minimum cycle basis with the desired property that all pairs of cycles share at most one path. 
%Further theory is needed to see if Algorithm~\ref{alg:merge_MCB} always succeeds. 
%If so, future studies may also investigate the expected run time of Algorithm~\ref{alg:merge_MCB}, and how much does a postprocessed minimum cycle basis deviate from a random one.
%Follow-up questions include expected run times of Algorithm~\ref{alg:merge_MCB}, and how much the postprocessed minimum cycle basis deviates from the random minimum cycle basis.
%For the $\rm C_{10}H_{16}$ simulations discussed in Section~\ref{sec:results}, Algorithm~\ref{alg:merge_MCB} always succeeds, and it typically takes 10-20 iterations.

%Using Theorem~\ref{Thm:MCB_Pair_Swap}, we modify a minimum cycle basis $\calM\to\calM'$ so that each pair of cycles in $\calM'$ intersect over a single path. We outline this procedure in Algorithm~\ref{alg:merge_MCB}.


\section{Applications to hydrocarbon pyrolysis systems}
\label{sec:results}


\begin{figure*}[t]
	\includegraphics[width=\linewidth]{TiceExs.pdf}
	\caption{Analysis of cycles for a system initialized to adamantane with formula $\rm C_{10}H_{16}$. \textbf{(a)}~Rate distribution of carbon rings obtained via an MCB (blue) and the relevant cycles (orange). The dashed lines represent a power law fit of these distributions for rings of length $15\le k \le 40$. \textbf{(b)} Four examples of {\tt sli} classes: (i) 16 rings of length 45, (ii) 2 rings of length 7, (iii) 4 rings of length 19, and (iv) 3 rings of length 8. \textbf{(c)} Joint distribution of ring lengths and number of rings in {\tt sli} classes. \textbf{(d)} Distribution of the number of rings in {\tt sli} class.}
	\label{fig:TICE_Res}
\end{figure*}

\subsection{Generating data} 
Here, we apply the theory for the {\tt pi} and {\tt sli} relations to analyze loop clusters in hydrocarbon pyrolysis. 
This data comes from a single simulation using a ReaxFF force field~\cite{ashrafExtension2017} in LAMMPS~\cite{plimptonFast1995,thompsonLAMMPS2022,aktulgaParallel2012}. 
The simulation was initialized with 800 adamantane molecules, $\rm{C_{10}H_{16}}$, resulting in $N_{\rm C}=8000$ carbon and $N_{\rm H}=12800$ hydrogen atoms. 
The system was ramped to a temperature of 4000K and pressure of 40.5~GPa. 
This was followed by an NPT simulation for a total of 1.2 ns. 
Bonds were sampled using a length and duration criterion as discussed in previous works~\cite{ruthCyclic2025,dufour-decieuxPredicting2023}. 
The accuracy of the ReaxFF force field~\cite{ashrafExtension2017} at these thermodynamic conditions is questionable, but it provides sufficiently rich data to investigate our ring statistics. %to be a useful for investigating our new experimental techniques.
%The focus of this is to highlight some of the statistics that can be measured from loop clusters.

We started sampling after 0.6 ns to allow the system to reach equilibrium. 
Every 1.2 ps (500 samples total) we decompose the loop cluster into {\tt pi} and {\tt sli} classes as discussed in Section~\ref{sub:DICE_TICE_Comp}. 
From this decomposition we collect the following:
\begin{itemize}
	\item A random MCB.
	\item The length and number of rings within each {\tt sli} class.
	\item The length, number of rings, and number of independent polyhedra within each {\tt pi} class.
\end{itemize}
%{\color{red}
%The number of rings in a }
%For the last bullet point, we interpret a polyhedron of a direct-interchangeability as a left kernel vector of the corresponding block within the matrix $\bm{R}$ defined in Section~\ref{sub:DICE_TICE_Comp}. For more details, see the diagonal blocks of the matrix shown in Figure~\ref{fig:DICE_TICE_alg}~(d). The number of linearly independent polyhedra of the direct-interchangeability class is the nullity of the corresponding block in $\bm{R}$.

%\textbf{Sampling minimum cycle basis.} 
%
%\add{We do use a random MCB.}
%
%\add{Add a figure with a reasonably normalized cycle path overlap matrix or refer to section where the concept is introduced. Mention we do not analyze this in detail.}

\subsection{Relevant cycles versus minimum cycle basis} 
\subsubsection{$\rm C_{10}H_{16}$ simulation data}
In Section~\ref{sub:rand_MCB}, we introduced Algorithm~\ref{alg:rand_MCB} for sampling MCBs uniformly at random. 
To highlight the value of this approach, we compare counts of large carbon rings from random MCBs and from the relevant cycles.
%We argue this approach is better than sampling a larger, unique set of carbon rings.
%We argue that it is more desirable to use a random minimum cycle basis versus a larger, unique set of carbon rings.  
Specifically, we focus on the ring rate distribution, or the number of rings per carbon atom as a distribution over ring length -- see Figure~\ref{fig:TICE_Res}~(a).
%First and foremost, we consider the ring rate distribution $\{\lambda\phi_k\}$, where $\lambda$ is the number of rings per carbon and $\phi_k$ is the fraction of rings with $k$ carbon atoms. See Figure~\ref{fig:TICE_Res}~(a). 
%We sample the ring rate distribution using a random minimum cycle basis and the relevant cycles -- see Figure~\ref{fig:TICE_Res}~(a).
This statistic has also been considered in computational studies of amorphous carbon~\cite{dernovMapping2025,deringerMachine2017}. 

Most rings are short with five or six carbon atoms. 
The amorphous structure of these hydrocarbon networks results in many large rings, as shown by the heavy tail of the ring rate distribution.
%Also, the amorphous structure of the simulated hydrocarbon networks results in many larger loops.
This tail approximately fits to a power law distribution $Ck^{-\alpha}$ with a cutoff.
%The tail of the loop length distribution fits closely to a power law with exponential cutoff $\phi_k\sim C k^{\alpha}e^{-k/k_C}$.

Importantly, the power law constant $\alpha$ varies with the method of sampling loops.
When the loops are sampled from an MCB we obtain $\alpha\approx1.36$. 
In contrast, sampling loops from the relevant cycles results in a power of $\alpha\approx0.47$.

Without the cutoff, the power law tail $k^{-0.47}$ of the loop rate distribution for the relevant cycles is not summable.
%The bond and small ring statistics are fairly consistent between simulations of different sizes.
We observe this cutoff increases with respect to both system size and the carbon-to-hydrogen ratio. 
Larger simulations are needed to see if this cutoff continues to increase with system size. 
If so, then the power law tail $k^{-0.47}$ will cause the relevant cycles to grow asymptotically faster than the dimension of the cycle space $\nu$:
%If the cutoff approaches infinity as the number of carbon atoms $N_{\rm C}\to\infty$, then the power law tail $k^{-0.5}$ will cause the relevant cycles to grow asymptotically faster than the dimension of the cycle space $\nu$
%\begin{equation}
%	|\RelCyc| \gg \nu = O(N_{\rm C}),\quad\text{as } N_{\rm C}\to\infty.
%\end{equation}
\begin{equation}
	\frac{|\RelCyc|}{\nu}\to\infty \quad\text{as}\quad N_{\rm C}\to\infty.
\end{equation}
For our simulation, the time-averaged number of relevant cycles is $|\RelCyc|=1736$ and the average dimension of the cycle space is $\nu=1311$.

\subsubsection{Random geometric graphs}

\begin{figure}[t]
	\centering
	\includegraphics{CycRateDist_RGG.pdf}
	\caption{Analysis of the cycle rate distribution in random geometric graphs as a function of graph size. Cycles are sampled using an MCB $\calM$ (top) and the relevant cycles $\RelCyc$ (bottom). Nodes are assigned random positions in the unit cube $[0,1]^3$ and the cutoff radius for edges is chosen so that the expected degree is 3. The random geometric graph with $10^2$, $10^3$, and $10^4$ nodes is sampled $10^5$, $10^4$ and $10^3$ times respectively.}
	\label{fig:CycRateDist_RGG}
\end{figure}

To further investigate the scaling of $|\RelCyc|$ with system size, we consider random geometric graphs~\cite{penroseRandom2003}. A random geometric graph is a graph with $n$ nodes randomly positioned in the $d$-dimensional cube $[0,1]^d$. A pair of nodes are joined by an edge if their distance is less than a given fixed radius. The cycle statistics from these random graphs share some qualitative characteristics with our simulation because they have a spatial component. We choose a dimension of $d=3$ and periodic boundary conditions for sampling edges. 

We choose the radius so that the average degree is 3. For comparison, the average degree of the carbon skeleton in our ReaxFF simulation is 2.25. We choose a higher mean degree of 3 because it results in more large loops. This reduces the number of samples needed to observe enough large loops.
%The radius for edges is chosen so that the average degree is 3. \add{Why we increased degree}

As in molecular dynamics simulation, the tail of the cycle rate distribution of MCBs in random geometric graphs roughly follows a power law. The cutoff for this power law distribution increases with system size. See Figure~\ref{fig:CycRateDist_RGG}, top. For the random geometric graph with $10^2$ nodes, the cycle rate distribution for the relevant cycles has the same shape as the cycle rate distribution given by MCBs. See the blue curve in Figure~\ref{fig:CycRateDist_RGG}, bottom. However, as system size increases, the number of large relevant cycles increases much faster than the number of large cycles in MCBs. For the random geometric graphs with $10^4$ nodes, we observed a shocking $9900$ relevant cycles per node. See the green curve in Figure~\ref{fig:CycRateDist_RGG}, bottom.

\subsubsection{A potential explosion of the relevant cycle count}

%It is unclear if the cutoff for the ring rate distribution in molecular dynamics will continue to increase with system size. 
It is unclear if the ring rate distribution in molecular dynamics will follow this same trend with system size.
Given this ambiguity, we argue that the relevant cycles should not be used to count rings in large amorphous chemical systems. Our studies of random geometric graphs suggest the number of large rings may explode at scales achievable using multiple GPUs in a molecular dynamics simulation. Beyond the relevant cycles, many other unique sets of rings in the chemistry literature are also worst case exponential in size. See, e.g., the review~\cite{bergerCounterexamples2004}. This is not an issue for our approach, where the size of an MCB grows linearly with the number of edges.

\subsection{Sources of additional cycles} 
The {\tt pi} and {\tt sli} classes can be used to determine where the extra relevant cycles are coming from. For the given initial conditions, we find that the {\tt sli} classes are the most common. Over 99\% of {\tt pi} classes reduce to a single {\tt sli} class. Most {\tt sli} classes consist of a large ring with nested smaller rings -- see Figure~\ref{fig:TICE_Res}~(b.i.). In some cases, there is only a minor distinction between large and small rings, e.g., the two larger 7-rings in Figure~\ref{fig:TICE_Res}~(b.ii.) are separated by a smaller 6-ring. In other cases, the structure of a {\tt sli} class is more complex -- see Figure~\ref{fig:TICE_Res}~(b.iii-iv.).

In Figure~\ref{fig:TICE_Res}~(c), we show the joint distribution of the lengths and number of rings within a {\tt sli} class. 
From this distribution, we observe most (roughly 85\%) of the {\tt sli} classes consist of a single cycle, and the average size of the {\tt sli} classes increases with ring length. 
In Figure~\ref{fig:TICE_Res}~(d), we show the size distribution of the {\tt sli} classes, which roughly follows a power law distribution. 
Interestingly, the oscillations in this size distribution are not noise, and they can be seen as horizontal bands in the joint length size distribution in Figure~\ref{fig:TICE_Res}~(c). 
The number of cycles in a {\tt sli} class tends more often to a composite than a prime number.

In summary, the relevant cycles have an increased number of large cycles forming {\tt sli} classes. At the same time, these large cycles are highly dependent, e.g., the cycles in the {\tt sli} classes in Figure~\ref{fig:TICE_Res}~(b.i-iii.) share many edges. 
%We argue it is better to use a random minimum cycle basis which is bounded in size and puts less weight on large cycles. 
Our codes can be used to isolate and visualize these {\tt sli} classes.% when they are of interest.


\begin{figure*}[t]
	\centering
	\includegraphics[width=\linewidth]{DICE_Res.pdf}
%	{\color{red}Include a more abstract poly}
	\caption{Analysis of polyhedra for a system initialized to adamantane with formula $\rm C_{10}H_{16}$. \textbf{(a)}~Size distribution of the {\tt pi} classes given by the number of polyhedra in the class. Only {\tt pi} classes with at least one polyhedron are considered, otherwise, they reduce to {\tt sli} classes. \textbf{(b)}~Average counts of common {\tt pi} classes per time step and their visualizations. \textbf{(c)}~The largest observed {\tt pi} class composed of five barellene polyhedra glued together by faces. \textbf{(d)}~A {\tt pi} class composed of a polyhedron that has a hexagonal face that is not uniquely defined because it is a {\tt sli} class.}
	\label{fig:DICE_Res}
\end{figure*}

\subsection{Polyhedra and diamond} 
Last, we investigate the {\tt pi} classes that do not reduce to {\tt sli} classes. As discussed in Section~\ref{sub:Poly}, these {\tt pi} classes can be described as polyhedra clusters. These are relatively rare, there are roughly eight of these {\tt pi} classes per time step. Since more theory is needed to describe these structures, we will consider individual polyhedra that are output by these codes in a case-by-case basis.
%Next, we investigate the polyhedron that are present within our system. One goal of this is to see if nanodiamond is present within our simulation. In short, there is not, but we discuss the theoretical tools we use to make this conclusion.
We label the {\tt pi} classes by the small hydrocarbons with the same carbon skeleton. Note, the {\tt pi} classes are not molecules but substructures within the loop cluster.

%While more rare, we also consider the direct-interchangeability classes that do not reduce to trivial-interchangeability classes. 
%These classes can be described as polyhedron clusters, and roughly 75\% of these direct-interchangeability classes consist of cycles of length 6. 
The majority of the {\tt pi} classes are composed of highly organized short rings. 
Roughly 75\% of these are composed of rings of length 6.
See Figure~\ref{fig:DICE_Res}~(b-d) for some examples of these {\tt pi} classes. 
Most of these {\tt pi} classes are small, consisting of a single polyhedron. See Figure~\ref{fig:DICE_Res}~(a). 
Sometimes, the faces of a polyhedron are not uniquely defined. For instance, the observed polyhedron in Figure~\ref{fig:DICE_Res}~(d) has two hexagonal cycles that can be chosen for its right face. However, polyhedra with non-unique faces are more rare.


The most obvious polyhedron to investigate is adamantane. 
Adamantane is the smallest diamondoid, a nanoscale structure with the same network topology as diamond, and it has been observed in petroleum~\cite{fangOrigin2013}.
Even though our simulation was initialized with adamantane molecules, adamantane was not common in our simulation. See Figure~\ref{fig:DICE_Res}~(b). 
 
Instead, the most commonly observed polyhedron is barallene\footnote{
	It is likely that many of these polyhedra more closely match the molecule bicyclo[2,2,2]octane with the same topology as barallene without double bonds. 
%	Some of these polyhedra may more closely match the molecule bicyclo[2,2,2]octane which has the same topology without double bonds. 
	However, we do not store bond order or infer double bonds. Hence, we use the simpler name barallene in both cases.
	}. %bicyclo[2,2,2]octane. 
%Unfortunately, unlike barallene which has the same shape but with double bonds, bicyclo[2,2,2]octane does not have a simpler name like ``barallane''.
%The shape of bicyclo[2,2,2]octane matches that of barallene, but it unfortunately does not have a shorthand name like barallane. 
Barallene can be used to form hexagonal diamond -- also called londsdaleite -- by gluing its faces to iceane. Iceane is relatively rare in our simulations, so it seems unlikely that this crystal structure will form. Instead, the majority of the larger {\tt pi} classes are composed of multiple barallene polyhedra glued together by faces. See twistane in Figure~\ref{fig:DICE_Res}~(b) and the larger structure in Figure~\ref{fig:DICE_Res}~(c).


\section{Discussion}
\label{sec:discussion}

\subsection{Sampling cycles (ring perception)} 
In Section~\ref{sub:rand_MCB}, we introduced a method for sampling MCBs uniformly at random, Algorithm~\ref{alg:rand_MCB}. 
This differs from other approaches because we do not sample a deterministic set of rings from our chemical data. 
Regardless, as a probabilistic structure, the random MCB is well-defined, and can be used to sample summary statistics from a chemical network. 
This approach results in fewer cycles, as discussed in Section~\ref{sec:results}. In particular, the size of the random MCB grows linearly with the number of bonds in the system.

A potential critique of using a random MCB is that it may not uncover all useful cycles.
%Also, a minimum cycle basis is always polynomial in size. 
%In this work, we do not attempt to sample all cycles that are useful.
Unfortunately, it is challenging to define such a set of cycles that is always small in size. 
The set of relevant cycles can be exponential in number, and it only includes cycles from MCBs. 
%The relevant cycles are exponential in number, and the set of relevant cycles is relatively small because it only includes cycles from a minimum cycle basis.
However, all cycles are within the span of a MCB by definition.
Thus, we argue it is better to interpret additional cycles as a consequence of the overlapping structure of a cycle cluster.
%Thus, we also introduce theory for describing how these cycles collect.
%We argue that since all of these cycles are within the span of a minimum cycle basis, it is more useful to describe how these basis cycles merge together. 
%Thus, any cycles outside of our basis are a consequence of this emergent structure.

In Theorem~\ref{Thm:MCB_Intersect} in Section~\ref{sub:pair_intersect}, we fully characterized how a pair of cycles in a given MCB may intersect as a series of paths that appear with the same order and orientation in both cycles. This gives insight on what data structures are needed to properly store these pairwise intersections of cycles. Furthermore, in Algorithm~\ref{alg:merge_MCB} we introduced a method of modifying an MCB so that all pairs of cycles intersect over a single path. This permits a dual graph representation of the loop cluster, where nodes represent cycles in an MCB and edges represent their intersection paths.
These intersection paths are easy to store as a list of nodes.

Theorem~\ref{Thm:MCB_Intersect} is specialized to pairs of cycles in an MCB. Its proof largely relies on the fact that these cycles preserve distances. Because of this, Theorem~\ref{Thm:MCB_Intersect} can be easily adapted to the full set of relevant cycles. The largest set of cycles this proof technique applies to is the set of shortest path rings from Ref.~\cite{franzblau1991computation}.

%Algorithm~\ref{alg:merge_MCB} biases the minimum cycle basis we sample from a graph. 
%Also, Algorithm~\ref{alg:merge_MCB} is a randomized algorithm, so the output minimum cycle basis is still well-defined as a random object. 
%Unfortunately, this algorithm is not guaranteed to complete. 
%For our hydrocarbon pyrolysis data, only a few steps ($\le 20$) are needed to compute a minimum cycle basis with pairs of cycles that intersect over single paths. 
%Further theory is needed to see if such a basis always exists, and if a more sophisticated algorithm is needed to compute it in general.  

%For our data we find this algorithm does not need many steps
%\textbf{Random MCB}
%This addresses lack of uniqueness without 
%
%\textbf{MCB pair intersect}
%
%{\color{red} what algorithms give, limitations go here.}

\subsection{Polyhedra.} In Section~\ref{sub:Poly}, we related lack of uniqueness of MCBs to polyhedral structures. These generalized polyhedra were defined as sets of non-flat faces obtained via a cycle and its expansion with respect to an MCB
\begin{equation*}
	\calP=\{C\}\cup\expand{\calM}{C},\quad C\in\RelCyc\setminus\calM.
\end{equation*}
The {\tt pi} classes were interpreted as clusters of these polyhedra.

A potential application of these polyhedra is in crystallography. A crystal can be represented by a periodic tiling of these polyhedra glued together by faces, e.g., diamond is a tetrahedral tiling of adamantane. This representation has a benefit when compared to the periodic unit cell representation of a crystal; the {\tt pi} relation is defined for general graphs. There is no requirement that the chemical network or its substructures are large. Small {\tt pi} classes might be interpreted as nanocrystals. Also, the input graph does not need to be periodic. In an inhomogeneous crystal (e.g., a mixture of cubic and hexagonal diamond) it may be useful to analyze how these polyhedra vary within the material.

Before investigating these applications, a pair of limitations must be addressed. 
Algorithmically, the polyhedra are left-kernel vectors of the diagonal blocks of the matrix $\bm{R}$ from Section~\ref{sub:DICE_TICE_Comp}. These correspond to {\tt sli} classes that sum to shorter cycles. 
This notion of polyhedra is abstract.
In some cases, this may be ok, where the {\tt sli} classes can be interpreted as generalized cycles and visualized as the union of their cycles (e.g., Figure~\ref{fig:DICE_Res}~(d)). 
If not, then a cycle may be sampled from each {\tt sli} class to represent a polyhedron as a set of face cycles. Polyhedra sampled using this latter approach are not unique.

The second limitation is that these polyhedra are the basis of a vector space. As with cycles, we are only interested in small polyhedra. Ultimately, one would want to properly define a minimum polyhedron basis, and to find an algorithm to compute it. This task is likely complex. A greedy approach may be to bias Algorithm~\ref{alg:rand_MCB} to obtain a matrix $\bm{R}$ that outputs small polyhedra.

\section{Conclusion}
\label{sec:conclusion}

We reviewed and advanced theory and algorithms for identifying and grouping cycles in graphs. 
We proposed a new method for sampling MCBs uniformly at random. The {\tt pi} relation helps to extract generalized polyhedra, that may be start-ups for forming crystalline structures. The {\tt sli} classes may comprise large numbers of relevant cycles as those around the perimeter carbon nanotubes. We propose to include only one representative of each {\tt sli} class in an MCB.
%We investigated algorithms and methods for analyzing cycles from graphs. We propose a new method for sampling cycles using a random MCB. 
We also provide theory for the format in which pairs of cycles in an MCB intersect, and we introduce a postprocessing step that makes these intersections single paths. 
This enables a dual graph representation for a cluster of cycles. We propose using this dual graph approach to model complex material networks. 
We illustrate this theory and algorithms by applying them to a reactive simulation of hydrocarbon pyrolysis MD simulation data initialized with 800 adamantane, $\rm C_{10}H_{16}$, molecules.

%We introduced a new method of sampling cycles using a minimum cycle basis. We also introduced a postproces

%{\color{red} 1 Paragraph.}

\section*{Acknowledgments}
This work was partially supported by ONR MURI Grant No. N000142412547 (M.C.), AFOSR MURI Grant No. FA9550-20-1-0397 (M.C.), and NSF GRFP Grant No. DGE 2236417 (P.R.).

\appendix

\section{Proofs for Section~\ref{sec:Theory}}
\subsection{Mutual relevance and useful lemmas}
\label{app:additional_lems}

To start,
Lemma~\ref{Lem:SmallCycs2MCB} below generalizes Statement 1 of Lemma~\ref{Lem:Relevant_MCB_Test} to sets of cycles. This lemma is useful for many of the results in this appendix.

\begin{lemma}
	Let $\calY$ be a set of cycles bounded in length by $k\in\mathbb{N}$, $|C|\le k$ for $C\in\calY$. Let $\calM$ be a given MCB. There exists a subset $\calY'\subseteq\calM$ also bounded in length by $k$, $|C'|\le k$ for $C'\in\calY'$, such that
	\begin{equation}
		\bigoplus_{C'\in\calY'}C' = \bigoplus_{C\in\calY}C.
		\label{eq:Y_replace_Lem}
	\end{equation}
	\label{Lem:SmallCycs2MCB}
\end{lemma}
\begin{proof}
	Let $\calY$ be a set of cycles bounded in length by $k$ and $\calM$ an MCB. By Definition~\ref{def:ex_op} of the expansion operator, the set of cycles
	\begin{equation}
		\calY' = \expand{\calM}{\bigoplus_{C\in\calY} C}
	\end{equation}
	is the unique subset $\calY'\subseteq\calM$ that satisfies Eq.~\eqref{eq:Y_replace_Lem}. We want to show the cycles $\calY'$ are also bounded in length by $k$. We observe, the cycles $\calY'$ satisfy the inclusion
	\begin{equation}
		\calY' \subseteq \bigcup_{{C\in\calY}} \expand{\calM}{C}.
	\end{equation}
	Thus, the cycles $\calY'$ are bounded in length by
	\begin{equation}
		|C'|\le|C|\le k, \quad C\in\calY,C'\in\expand{\calM}{C}
	\end{equation}
	where the first inequality is Statement 1 of Lemma~\ref{Lem:Relevant_MCB_Test}.
\end{proof}

Next, we introduce mutual relevance, which generalizes the notion of relevance to sets of cycles. We will use mutual relevance in our proofs regarding the {\tt pi} relation.

\begin{definition}
	A set of cycles $\mathcal{X}$ is \emph{mutually relevant} if there exists an MCB $\calM$ such that $\mathcal{X}\subseteq\calM$.
	\label{def:mut_rel}
\end{definition}

\noindent We emphasize linearly independent, relevant cycles are not necessarily mutually relevant. See, e.g., the cycles $C_L,C_L',C_R$ in Figure~\ref{fig:overlapping_diamonds} in Section~\ref{sub:background_partitions}.% are not a subset of a minimum cycle basis.

Recall, vectors $v_1,\hdots,v_k$ from a vector space are linearly independent if the only linear combination that achieves
\begin{equation*}
	c_1v_1+\hdots +c_kv_k=0
\end{equation*}
is when $c_1=\hdots=c_k=0$. In Lemma~\ref{Lem:mut_rel_2} below, we define mutual relevance in a way that resembles this notion of linear independence. 


\begin{lemma}
	A set of cycles $\calX\subseteq\RelCyc$ of equal length are not mutually relevant if and only if there exist cycles $C_1,C_2,\hdots,C_k\in\calX$ such that
	\begin{equation}
		C_1\oplus C_2 \oplus\hdots\oplus C_k = \smCyc
		\label{eq:mut_rel_ind}
	\end{equation}
	where $\smCyc$ is the sum of a set of smaller cycles $\calY$, $|C|<|C_1|$ for $C\in\calY$.
	\label{Lem:mut_rel_2}
\end{lemma}
\begin{proof}
	($\Rightarrow$) Let $\calX\subseteq\RelCyc$ not be mutually relevant. The individual cycles in $\calX$ are relevant. Hence, there exists a mutually relevant subset $\calY\subset\calX$ and a cycle $C\in\calX\setminus\calY$ such that $\calY\cup\{C\}$ is not mutually relevant. By the definition of mutual relevance, let $\calM$ be an MCB where $\calY\subseteq\calM$.
	
	We want to construct a subset of $\calX$ that sums to smaller cycles. First, we expand $C$ with respect to $\calM$,
	\begin{equation}
		C = \enspace\bigoplus_{\mathclap{C'\in\expand{\calM}{C}}} C' = \enspace\bigoplus_{\mathclap{\substack{C'\in\expand{\calM}{C}\\|C'|=|C|}}} C' \enspace\oplus\enspace \bigoplus_{\mathclap{\substack{C''\in\expand{\calM}{C}\\|C''|<|C|}}}C''.
	\end{equation}
%	Next, we move the cycles of length $|C'|=|C|$ to the left-hand side
	We move the cycles of length $|C|=|C'|$ to the left-hand side, so that the right-hand side contains only smaller cycles:
	\begin{equation}
		C\oplus\enspace\bigoplus_{\mathclap{\substack{C'\in\expand{\calM}{C}\\|C'|=|C|}}}C'\enspace =  \enspace\bigoplus_{\mathclap{\substack{C''\in\expand{\calM}{C}\\|C''|<|C|}}}C''. %= \smCyc.
		\label{eq:C_exp_toSmCyc}
	\end{equation}
	We argue the cycles in the left-hand side of Eq.~\eqref{eq:C_exp_toSmCyc} belong to $\calY\cup\{C\}\subseteq\calX$.
	
	Assume towards a contradiction that there exists $C'\in\expand{\calM}{C}$ of length $|C'|=|C|$ such that $C'\notin\calY$. 
	By Statement 2 of Lemma~\ref{Lem:Relevant_MCB_Test}, we can exchange $C$ for $C'$ to obtain a new minimum cycle basis $\calM'=(\calM\setminus\{C'\})\cup\{C\}$.
	Indeed, this contradicts our initial assumptions, as it implies $\calY\cup\{C\}\subseteq\calM'$ is mutually relevant.
	Thus, the cycles in the left-hand side of Eq.~\eqref{eq:C_exp_toSmCyc} form a subset of $\calX$ that sums to smaller cycles as desired.
	
	($\Leftarrow$) Let $C_1,C_2,\hdots,C_k$ satisfy Eq.~\eqref{eq:mut_rel_ind}. We want to show that these cycles are not mutually relevant. We assume $C_2,\hdots,C_k$ are mutually relevant, otherwise our proof is complete. Let $\calM$ be an MCB such that $C_2,\hdots,C_k\in\calM$. By Lemma~\ref{Lem:SmallCycs2MCB} we can choose smaller cycles $\calY'\subseteq\calM$ satisfying Eq.~\eqref{eq:mut_rel_ind}. Re-arranging Eq.~\eqref{eq:mut_rel_ind} and substituting in $\calY'$ gives
	\begin{equation}
		C_1 = C_2\oplus\hdots\oplus C_k\oplus\bigoplus_{C'\in\calY'}C'.
		\label{eq:C_exp_all_mut_rel}
	\end{equation}
	Hence, $C_1$ is in the span of cycles in $\calM$, and $C_1$ does not belong to $\calM$. 
%	The cycles in the right-hand side of Eq.~\eqref{eq:C_exp_all_mut_rel} are contained in $\calM$. As such, $C_1$ does not belong to $\calM$. 
	Here, $\calM$ is an arbitrary MCB containing $C_2,\hdots,C_k$, so there is no MCB containing $C_1,C_2,\hdots,C_k$. In other terms, $C_1,C_2,\hdots,C_k$ are not mutually relevant.
\end{proof}


\subsection{Polyhedron-interchangeability}
\label{app:DICE_Proofs}

%{\color{red} Acknowledge this is a weaker condition than Lemma~\ref{Lem:dir_intchg_2}.}

Here we prove the results for the {\tt pi} relation from Section~\ref{sub:DICE_Props}. These proofs rely on the lemmas and definitions from Appendix~\ref{app:additional_lems}.

First, we prove an equivalent definition of the {\tt pi} relation. Instead of proving Lemma~\ref{Lem:dir_intchg_2}, we prove Lemma~\ref{Lem:dir_intchg_3} below. Lemma~\ref{Lem:dir_intchg_3} has weaker requirements than Lemma~\ref{Lem:dir_intchg_2}. In particular, the shorter cycles used to exchange $C_1$ for $C_2$ no longer need to belong to an MCB. It is straightforward to prove Lemma~\ref{Lem:dir_intchg_2} from Lemma~\ref{Lem:dir_intchg_3}. Lemma~\ref{Lem:dir_intchg_3} is often easier to verify in proofs.
%Lemma~\ref{Lem:dir_intchg_3} is a weaker

\begin{lemma}
	Let $C_1,C_2\,{\in}\,\RelCyc$ where $|C_1|\,{=}\,|C_2|$. $C_1{\pInt}\, C_2$ if and only if there exists equal length cycles $\calX$, $|C|\,{=}\,|C_1|$ for $C\,{\in}\,\calX$ where $C_1,C_2\,{\notin}\,\calX$, and shorter cycles $\calY$, $|C|{<}|C_1|$ for $C{\in}\calY$, such that
	\begin{enumerate}
		\item (mutual relevance) there exists an MCB $\calM$ where $\calX\cup\{C_2\}\subseteq\calM$,
		\item (expansion) $C_1$ is given by the sum 
		\begin{equation}
			C_1 = C_2\oplus \bigoplus_{C'\in\calX} C' \oplus \bigoplus_{C''\in\calY} C''.
			\label{eq:dir_exchg_split}
		\end{equation}
%		where the equal length cycles in the expansion of $C_1$ in $\calM$ are $\calX=\{C\in\expand{\calM}{C_1}:|C|=|C_1|\}$.
	\end{enumerate}
	The set of equal length cycles in the expansion of $C_1$ in $\calM$ is $\calX$.
	\label{Lem:dir_intchg_3}
\end{lemma}

\begin{proof}
	($\Leftarrow$) Let $\calX$ and $\calY$ satisfy the given hypothesis. We want to show $C_1\pInt C_2$. By mutual relevance, let $\calM_2$ be an MCB such that $\calX\cup\{C_2\}\subseteq\calM$. By Lemma~\ref{Lem:SmallCycs2MCB}, there exists shorter cycles $\calY'\subseteq\calM_2$ with the same sum as $\calY$. We replace $\calY'$ for $\calY$ in Eq.~\eqref{eq:dir_exchg_split}
	\begin{equation}
		C_1 = C_2\oplus \bigoplus_{C'\in\calX} C' \oplus \bigoplus_{C''\in\calY'} C'',
		\label{eq:dir_exchg_split2}
	\end{equation}
	The right-hand side of Eq.~\eqref{eq:dir_exchg_split2} is the expansion of $C_1$ in $\calM_2$. By Lemma~\ref{Lem:Relevant_MCB_Test}, Statement 2, we may exchange to a new MCB $\calM_1=(\calM_2\setminus\{C_2\})\cup\{C_1\}$. Thus, $C_1\pInt C_2$.
	
	($\Rightarrow$) Let $C_1{\pInt}\, C_2$ via the MCB $\calM_2$. We want to construct sets $\calX$ and $\calY$ that satisfy the given hypothesis. First, we show that $C_2$ belongs to the expansion of $C_1$ in $\calM_2$, or $C_2\in\expandd{\calM_2}{C_1}$. To do so, consider the MCB $\calM_1=(\calM_2\setminus\{C_2\})\cup\{C_1\}$ given by the {\tt pi} relation. If $C_2\notin\expandd{\calM_2}{C_1}$, then the full expansion $\expandd{\calM_2}{C_1}$ belongs to $\calM_1$. This is not possible, since the set $\{C_1\}\cup\expandd{\calM_2}{C_1}$ is linearly dependent,
	\begin{equation*}
		C_1 = \bigoplus_{C'\in\expandd{\calM_2}{C_1}}C',
	\end{equation*}
	so they cannot all belong to $\calM_1$. 
	Since $C_2$ is in the expansion of $C_1$ in $\calM_2$, we may decompose the expansion as follows
	\begin{equation}
		C_1 = C_2\oplus\enspace
		\bigoplus_{\mathclap{\substack{C'\in\expand{\calM_2}{C_1}\setminus\{C_2\}\\|C'|=|C_1|}}}C'
		\quad\oplus\quad
		\bigoplus_{\mathclap{\substack{C''\in\expand{\calM_2}{C_1}\\|C''|<|C_1|}}}C'.
	\end{equation}
	Thus, the cycles $\calX$ are the cycles of length $|C_1|$ in $\expand{\calM_2}{C_1}$ other than $C_2$, and the cycles $\calY$ are the shorter cycles in $\expand{\calM_2}{C_1}$.
\end{proof}

\begin{figure*}
	\includegraphics{DICE_Transitivity.pdf}
	\caption{Construction of transitivity for the {\tt pi} relation. \textbf{(a)} A graph with relevant cycles $\RelCyc=\calX\cup\calY\cup\{C_1,C_2,C_3\}$, where $C_1\pInt C_2$ and $C_2\pInt C_3$ via $\calX$ and $\calY$ respectively. \textbf{(b)}~MCBs $\calM_j$ such that the mutually relevant sets $\calZ_j=\calX\cup\{C_2\}\cup\{Y_1,\hdots,Y_j\}$ satisfy $\calZ_j\subseteq\calM_j$. The relevant cycles not in $\calM_j$ are shaded with gray stripes. \textbf{(c)} A final MCB $\calM$ obtained by exchanging $C_3$ for $C_2$ in $\calM_4$. $C_3$ is in the expansion of $C_1$ in $\calM$ so $C_1\pInt C_3$.}
	\label{fig:DICE_trans_ex}
\end{figure*}

Next, we prove that {\tt pi} is an equivalence relation, Theorem~\ref{Thm:DirInt_Eq_Rel}. 
First, consider the following property about the expansion operator. Let $\calB$ be a cycle basis. If we sum $C_1$ and $C_2$ and expand in $\calB$ then we obtain
%First, recall the expansion operator $\expand{\calB}{C}$ gives the unique set of cycles in the basis $B$ that sums to $C$. If we sum $C_1$ and $C_2$ and expand in $\calB$ we obtain
\begin{equation}
	C_1\oplus C_2 = \bigoplus_{\mathclap{C'\in\expand{\calB}{C_1}}}C' 
	\;\oplus\; 
	\bigoplus_{\mathclap{C''\in\expand{\calB}{C_2}}} C'' 
	\quad=\quad
	\bigoplus_{\mathclap{C'\in\expand{\calB}{C_1}\symdiff\expand{\calB}{C_2}}} C'
\end{equation}
where the reduced sum in the final equality is over the symmetric difference $\expand{\calB}{C_1}\symdiff\expand{\calB}{C_2}$ because of cancellations of the form $C\oplus C=\emptyset$. Thus, the expansion operator is linear
\begin{equation}
	\expand{\calB}{C_1\oplus C_2} = \expand{\calB}{C_1}\symdiff\expand{\calB}{C_2}
	\label{eq:expand_linear}
\end{equation}
where addition is given by the symmetric difference operator. We will use linearity of the expansion operator in the following proofs.

%\add{Include linearity of the expansion operator
%	As a final preliminary, we characterize the expansion operator over sums. Let $C_1,C_2$ be cycles and $\calB$ a basis. Then, if we sum $C_1$ and $C_2$ and expand in $\calB$ we obtain
%	\begin{equation}
%		C_1\oplus C_2 = \bigoplus_{\mathclap{C'\in\expand{\calB}{C_1}}}C' 
%		\;\oplus\; 
%		\bigoplus_{\mathclap{C''\in\expand{\calB}{C_2}}} C'' 
%		\quad=\quad
%		\bigoplus_{\mathclap{C'\in\expand{\calB}{C_1}\symdiff\expand{\calB}{C_2}}} C'
%	\end{equation}
%	where the reduced sum in the final equality is over the symmetric difference $\expand{\calB}{C_1}\symdiff\expand{\calB}{C_2}$ is due to cancellations of the form $C\oplus C=\emptyset$. Thus, the expansion operator is linear
%	\begin{equation}
%		\expand{\calB}{C_1\oplus C_2} = \expand{\calB}{C_1}\symdiff\expand{\calB}{C_2}
%		\label{eq:expand_linear}
%	\end{equation}
%	where addition is given by the symmetric difference operator.}

%{\color{red} equivalence relation}

\begin{proof}[Proof of Theorem~\ref{Thm:DirInt_Eq_Rel}]
	We show that the {\tt pi} relation satisfies three properties. The reflexive property, $C\,{\pInt}\, C$, is immediate from an MCB $\calM,C\in\calM$, since $\calM=(\calM\setminus\{C\})\cup\{C\}$ is already an MCB.
	
	
	%	\begin{enumerate}
	(Symmetry) Here we use the definition of {\tt pi} as in Definition~\ref{def:dir_intchg}. Let $C_1\pInt C_2$ via $\calM_2$, $C_2\in\calM_2$, an MCB such that the exchange $\calM_1=(\calM_2\setminus\{C_2\})\cup\{C_1\}$ is an MCB. Then, the reverse exchange $\calM_2=(\calM_1\setminus\{C_1\})\cup\{C_2\}$ also results in a MCB. Hence, $C_2\pInt C_1$.
	
	
	(Transitivity) Transitivity is trickier to verify for the {\tt pi} relation. As an intuition leading up to this result, recall from Section~\ref{sub:Poly} that {\tt pi} relations can be interpreted as polyhedra. For instance, consider the graph depicted in Figure~\ref{fig:DICE_trans_ex} (a). The cycles $C_1$ and $C_2$ satisfy $C_1{\pInt}\, C_2$ since they belong to the same cube $\calP_1=\{C_1,C_2\}\cup\calX$. Similarly, $C_2\,{\pInt}\, C_3$ since they belong to the cube $\calP_2=\{C_2,C_3\}\cup\calY$. The transitive relation can be verified using the rectangular prism $\calP_3=\calP_1\symdiff\calP_2 = \{C_1,C_3\}\cup\calX\cup\calY$. Unfortunately, the symmetric difference $\calP_3$ does not define a {\tt pi} relation in general, as the cycles $\calP_3\setminus\{C_1\}$ may not be mutually relevant. Instead, we reproduce this procedure using mutually relevant sets below. We illustrate these methods in Figure~\ref{fig:DICE_trans_ex} as well.
	%The limitation of this approach is that the cycles $\calP_3\setminus\{C_1\}$ are not necessarily mutually relevant, so the prism $\calP_3$ does not always define .
	
	To verify transitivity, we use the definition of the {\tt pi} relation given by Lemma~\ref{Lem:dir_intchg_3}. 
	Let $C_1{\pInt}\, C_2$ and $C_2\,{\pInt}\, C_3$ via $\calX$ and $\calY$ up to smaller cycles respectively:
	\begin{align}
		C_1 &= C_2 \oplus \bigoplus_{C\in\calX} C \oplus\text{(smaller cycles)}, 
		\label{Eq:C1C2DirEx}
		\\
		C_3 &= C_2 \oplus \bigoplus_{C\in\calY} C\oplus\text{(smaller cycles)},
		\label{Eq:C2C3DirEx}
	\end{align}
	where (i) $|C|=|C_1|$ for $C\in\calX\cup\calY$, (ii) $\calX\cup\{C_2\}$ and $\calY\cup\{C_2\}$ are mutually relevant, and  (iii) $C_1,C_2\notin\calX$ and $C_2,C_3\notin\calY$.
	Note, the set $\calX\cup\{C_1\}$ is mutually relevant, as it belongs to the MCB obtained by the exchange $\calM_1=(\calM_2\setminus\{C_2\})\cup\{C_1\}$, where $\calM_2$ is an MCB such that $\calX\cup\{C_2\}\subseteq\calM_2$.
	$\calY\cup\{C_3\}$ is mutually relevant by a similar argument.
	%	Note, $\calX\cup\{C_1\}$ and $\calY\cup\{C_3\}$ are mutually relevant as they are subsets of minimum cycle bases obtained by exchanging $C_2$ for $C_1$ and $C_3$, respectively, as guaranteed by Lemma~\ref{Lem:DICE_MCB_Ex}.
	
	
	% 	As an auxiliary step, we merge the mutually relevant sets $\calX\cup\{C_2\}$ and $\calY\cup\{C_2\}$.
	
	%	To show transitivity, we will construct a minimum cycle basis $\calM$ such that $C_1\in\expand{\calM}{C_3}$.
	%	Then, $C_1\DirEx C_3$ by Statement 2 of Lemma~\ref{Lem:Relevant_MCB_Test}.
	%		To show the transitive property, we will expand $C_1$ in terms of a mutually relevant set of cycles including $C_3$. 
	%	Before this, consider the degenerate case $C_3\in\calX\cup\{C_2\}$. In this case, $C_1\DirEx C_3$ via $\calX'=(\calX\cup\{C_2\})\setminus\{C_3\}$ by re-arranging Eq.~\eqref{Eq:C1C2DirEx}. For the remainder of this proof, we assume $C_3\notin\calX\cup\calY\cup\{C_2\}$. By a symmetric argument for $C_1$ and Eq.~\eqref{Eq:C2C3DirEx}, we focus on the case $C_1\notin\calX\cup\calY\cup\{C_2\}$.
	
	%		Now, we iteratively construct a mutually relevant set containing $\calX\cup\calY\cup\{C_2\}$ up to smaller cycles. 
	Motivated by the polyhedron $\calP_3$ combining $\calP_1$ and $\calP_2$ in our example, we iteratively merge the mutually relevant sets $\calX\cup\{C_2\}$ and $\calY$. 
	First, initialize $\calZ_0=\calX\cup\{C_2\}$ and arbitrarily order the cycles $\calY=\{Y_1,\hdots,Y_k\}$. Then, for $j=1,\hdots,k$ we do the following
	\begin{enumerate}
		\item if there exists $\tilde{Y}_j\in\Span (\calZ_{j-1})$ such that $Y_j=\tilde{Y}_j\oplus\text{(smaller cycles)}$ then $\calZ_{j}=\calZ_{j-1}$,
		\item otherwise, set $\calZ_{j}=\calZ_{j-1}\cup\{Y_{j}\}$.
	\end{enumerate}
	Recall from Lemma~\ref{Lem:mut_rel_2}, a set of cycles is mutually relevant if it does not have a subset that sums to smaller cycles. Thus, at each step the cycles $\calZ_j$ are mutually relevant. Furthermore, all cycles $C\in\calX\cup\calY\cup\{C_2\}$ are within the span of $\calZ_k\subseteq\calX\cup\calY\cup\{C_2\}$ up to smaller cycles.
	%		{\color{red}From Lemma 8 a set of cycles are mutually relevant if and only if they cannot sum to smaller cycles. By construction, there is not a subset of $\calZ_k$ that sums to smaller cycles. Last bits own sentence...}
	%		Recall, by Lemma~\ref{Lem:mut_rel_2}, that a set of cycles is mutually relevant if it does not have a subset that
	%		The final set $\calZ_k\subseteq\calX\cup\{C_2\}\cup\calY$ is mutually relevant by Lemma~\ref{Lem:mut_rel_2}, where $\calX\cup\{C_2\}\subseteq\calZ_k$ and all $Y_j\in\calY$ are in the span of $\calZ_k$ up to smaller cycles.
	
	By definition of mutual relevance, let $\calM_k$ be an MCB containing $\calZ_k$. Consider the expansion of $C_3$ with respect to $\calM_k$. 
	By Eq.~\eqref{Eq:C2C3DirEx}, $C_3$ is in the span of $\calZ_k$ up to smaller cycles.
	%By Eq.~\eqref{Eq:C2C3DirEx}, the cycles in $\expand{\calB}{C_3}$ of length $|C_1|$ are contained in $\calZ_k$, where we replace $Y_j$ with $\tilde{Y}_j$ if $Y_j\notin\calZ_k$. 
	%		Note, $\calY\cup\{C_3\}$ is mutually relevant by symmetry, so $C_3$ is not in the span of $\calY$ up to smaller cycles by Corollary~\ref{Cor:Mutual_Rel_Extend}. 
	By Lemma~\ref{Lem:mut_rel_2}, the mutually relevant cycles $\calY\cup\{C_3\}$ cannot be combined to sum to smaller cycles. Thus, there exists an additional cycle in the expansion $C^\star\in\expand{\calM_k}{C_3}$, $C^\star\in\calZ_k\setminus\calY$, of length $|C^\star|=|C_3|$. By Statement 2 of Lemma~\ref{Lem:Relevant_MCB_Test} we can exchange $C_3$ for $C^\star$ to obtain a new MCB $\calM=(\calM_k\setminus\{C^\star\})\cup\{C_3\}$.
	%		Our final basis is $\calM=(\calM_k\setminus\{C^\star\})\cup\{C_3\}$, which is minimal by Statement 2 of Lemma~\ref{Lem:Relevant_MCB_Test}.
	%		Thus, there exists $C^\star\in\expand{\calB}{C_3}$ of length $|C_1|$ such that $|C^\star|=|C_1|$ and $C^\star\in\calX\cup\{C_2\}$. Furthermore, $\calB'=(\calB\setminus \{C^\star\})\cup\{C_3\}$ is a minimum cycle basis by Lemma~\ref{Lem:Relevant_MCB_Test} and $\calZ' = (\calZ\setminus\{C^\star\}) \cup\{C_3\}\subseteq\calB'$ is mutually relevant.
	
	Finally, we want to show that $C_1\pInt C_3$ in $\calM$ by showing $C_3\in\expand{\calM}{C_1}, |C_1|=|C_3|$. By linearity of the expansion operator
	\begin{equation}
		\expand{\calM}{C_1} = \expand{\calM}{C_1\oplus C_3}\symdiff \{C_3\}.
		\label{eq:C1C3_Ex_M}
	\end{equation}
	where $C_1=(C_1\oplus C_3)\oplus C_3$ and $\expand{\calM}{C_3}=\{C_3\}$. Next, consider the identity
	\begin{equation}
		\expand{\calM_k}{C_1\oplus C_3} = \expand{\calM_k}{C_1}\symdiff \expand{\calM_k}{C_3}.
	\label{eq:C1C3_Ex_Mk}
	\end{equation}
	By construction, $C^\star$ belongs to both terms in the right-hand side of Eq.~\eqref{eq:C1C3_Ex_Mk}, where the cycles of length $|C_1|$ in $\expand{\calM_k}{C_1}$ are $\calX\cup\{C_2\}$ by Eq.~\eqref{Eq:C1C2DirEx}. Hence, $C^\star$ cancels in the symmetric difference in Eq.~\eqref{eq:C1C3_Ex_Mk} and
	\begin{equation}
		\expand{\calM_k}{C_1\oplus C_3} \subseteq \calM_k\setminus\{C^\star\} \subset \calM.
	\end{equation}
	Thus, $\expand{\calM_k}{C_1\oplus C_3}$ is the expansion of $C_1\oplus C_3$ in $\calM$ as well. We substitute $\expand{\calM_k}{C_1\oplus C_3}$ for $\expand{\calM}{C_1\oplus C_3}$ in Eq.~\eqref{eq:C1C3_Ex_M}:
	\begin{equation}
		\expand{\calM}{C_1} = \expand{\calM_k}{C_1\oplus C_3}\cup\{C_3\}
	\end{equation}
	where $C_3\notin\calM_k$.
%	Finally, we will show that $C_1$ and $C_3$ are directly-interchangeable via the minimum cycle basis $\calM$ as in Lemma~\ref{Lem:dir_intchg_2}. In particular, we want to show that $C_1$ belongs to the expansion $C_3\in\expand{\calM}{C_1}$. 
%	Finally, we want to show $C_1$ is directly-interchangeable for $C_3$ in $\calM$ by showing $C_3\in\expand{\calM}{C_1}, |C_1|=|C_3|$.
%	First, by linearity of the expansion operator we obtain
%	\begin{equation}
%		\expand{\calM_k}{C_1\oplus C_3} = \expand{\calM_k}{C_1}\symdiff \expand{\calM_k}{C_3}.
%		\label{Eq:C1C3_Ex_Mk}
%	\end{equation}
%	By construction, $C^\star$ belongs to both $\calX\cup\{C_2\}$ and $\expand{\calM_k}{C_3}$. 
%	Moreover, the cycles of length $|C_1|$ in the expansion $\expand{\calM_k}{C_1}$ are $\calX\cup\{C_2\}$ by Eq.~\eqref{Eq:C1C2DirEx}, so $C^\star$ belongs to both terms in the right-hand side of Eq.~\eqref{Eq:C1C3_Ex_Mk}.
%%	Thus, $C^\star$ belongs to both terms in the right-hand side of Eq.~\eqref{Eq:C1C3_Ex_Mk}, where the cycles of length $|C|=|C_1|$ in the expansion $\expand{\calM_k}{C_1}$ are $\calX\cup\{C_2\}$ by Eq.~\eqref{Eq:C1C2DirEx}. 
%	Hence, $C^\star$ cancels in the symmetric difference in the right-hand side of Eq.~\eqref{Eq:C1C3_Ex_Mk}, and \begin{equation*}
%		\expand{\calM_k}{C_1\oplus C_3}\subseteq \calM_k\setminus\{C^\star\}\subseteq\calM.
%	\end{equation*} 
%	Therefore, $\expand{\calM_k}{C_1\oplus C_3}$ is also the expansion of $C_1\oplus C_3$ in $\calM$. Once again, by linearity of the expansion operator
%	\begin{equation}
%		\expand{\calM}{C_1\oplus C_3} = \expand{\calM}{C_1}\symdiff \{C_3\},
%	\end{equation}
%	where $\expand{\calM}{C_3}=\{C_3\}$ by $C_3\in\calM$.
%	Equating $\expand{\calM_k}{C_1\oplus C_3}=\expand{\calM}{C_1\oplus C_3}$ gives
%	\begin{equation}
%		\expand{\calM}{C_1}\oplus\{C_3\} = \expand{\calM_k}{C_1\oplus C_3}.
%	\end{equation}
%	Finally, we move $C_3$ to the right-hand side
%	\begin{equation}
%		\expand{\calM}{C_1} = \expand{\calM_k}{C_1\oplus C_3}\cup\{C_3\}.
%	\end{equation}
%	where $C_3\notin\calM_k$.
	To conclude $C_3\in\expand{\calM}{C_1}$ so $C_1\pInt C_3$ in $\calM$ as desired.
\end{proof}

Next, we prove the {\tt pi} classes can be computed using a single MCB.

\begin{proof}[Proof of Theorem~\ref{Thm:DirInt_Comp}]
	Let $\calM$ be a given MCB. 
	We have already shown that if $C_1\in\RelCyc$ and $C_2\in\expand{\calM}{C_1}$, $|C_1|=|C_2|$, then $C_1\pInt C_2$. 
%	We have already shown that a cycle $C_1\in\RelCyc\setminus\calM$ is {\tt pi}-interchangeable for $C_2\in\calM$ if $C_2\in\expand{\calM}{C_1}$, $|C_1|=|C_2|$. 
	See Lemma~\ref{Lem:Relevant_MCB_Test}, Statement 2. We want to show the {\tt pi} classes are obtained from the transitive closure of these relations.
%	The validity of the sufficiency condition in Eq.~\eqref{Eq:DirEx_Alg_Th} is an immediate consequence of Statement 2 of Lemma~\ref{Lem:Relevant_MCB_Test}. Here, we show that all direct-interchangeability relations can be obtained via the transitive closure of Eq.~\eqref{Eq:DirEx_Alg_Th}. 
	
	Let $C_1\pInt C_2$. 
%	We want to obtain this relation via Eq.~\eqref{Eq:DirEx_Alg_Th}. 
	By Lemma~\ref{Lem:dir_intchg_3}, there exists a set of cycles $\calX$, $C_1,C_2\notin\calX$, such that
	\begin{equation}
		C_1 = C_2 \oplus\bigoplus_{C\in\calX}C\oplus\smCyc
		\label{eq:C1C2_4_DirExAlg}
	\end{equation}
	where $|C|=|C_1|$ for $C\in\calX$ and $\calX\cup\{C_2\}$ is mutually relevant. 
	In the following identity, we expand Eq.~\eqref{eq:C1C2_4_DirExAlg} in $\calM$ and use linearity of the expansion operator
%	We expand Eq.~\eqref{eq:C1C2_4_DirExAlg} in $\calM$ and use linearity of the expansion operator to obtain the following identity
%	By linearity of the expansion operator, the expansion of Eq.~\eqref{eq:C1C2_4_DirExAlg} in $\calM$ is given by
%	It is useful to expand Eq.~\eqref{eq:C1C2_4_DirExAlg} in terms of the given minimum cycle basis $\calM$
	\begin{equation}
		\begin{split}
		\expand{\calM}{C_1}\symdiff\expand{\calM}{C_2}\symdiff\left(\symdiffseries_{C\in\calX}\expand{\calM}{C}\right) \\= \smCyc
		\end{split}
		\label{eq:DICE_alg_Exs}
	\end{equation}
%	where the symmetric difference moves outside the expansion operator by linearity, Eq.~\eqref{eq:expand_linear}. 
	In Eq.~\eqref{eq:DICE_alg_Exs}, $\smCyc$ represents a set of smaller cycles $\calY$, $|C'|<|C_1|$ for $C'\in\calY$, which 
	belonging to $\calM$ by Lemma~\ref{Lem:SmallCycs2MCB}.
%	belong to the given minimum cycle basis, $\calY\subseteq\calM$, by Lemma~\ref{Lem:SmallCycs2MCB}. 
	We utilize Eq.~\eqref{eq:DICE_alg_Exs} to obtain {\tt pi} relations from $\calM$ which we can stitch together to show that $C_1\pInt C_2$.% $C_1$ and $C_2$ are {\tt pi}-interchangeable.
%	Observe, a cycle $C\in\{C_1,C_2\}$ is directly-interchangeable for a cycle in the expansion $C'\in\expand{\calM}{C}$ of length $|C'|=|C_1|$ as in Eq.~\eqref{Eq:DirEx_Alg_Th}. We utilize Eq.~\eqref{eq:DICE_alg_Exs} to stitch these observed direct-interchangeability relations together. %By relating the expansions of the cycle $\calP=\{C_1,C_2\}\cup\calX$ as in Eq.
	
	%	We will now justify the feasibility of the following procedure
	Consider the following procedure
	\begin{enumerate}
		\item Initialize $\calX_0=\emptyset$ and index $j=1$.
		\item Let $B_j\in\calM$ be a cycle of length $|B_j|{=}|C_1|$ in the expansion of $\calX_{j-1}\cup\{C_2\}$ in $\calM$,
		\begin{equation*}
			%			M_j \in\expand{\calM}{C},\quad\text{such that }C\in\calX_{j-1}\cup\{C_2\}\text{ and } |M_j|= |C_1|.
			B_j\in
			%			 \expand{\calM}{C_2\oplus\bigoplus_{C\in\calX_{j-1}}C} = 
			\expand{\calM}{C_2}\symdiff\bigg(\symdiffseries_{{C\in\calX_{j-1}}}\expand{\calM}{C}\bigg).
		\end{equation*}
%		such that $|B_j|=|C_1|$
		\item Let $X_j\in (\calX\cup\{C_1\})\setminus\calX_{j-1}$ such that $B_j\in\expand{\calM}{X_j}$ and set $\calX_j = \calX_{j-1}\cup\{X_j\}$.
		\item If $X_j=C_1$ then exit, otherwise repeat steps 2-4 with $j\gets j+1$.
	\end{enumerate}
	The feasibility of this procedure is as follows.
	In step 2, the cycles $\calX_{j-1}\cup\{C_2\}\subseteq\calX\cup\{C_2\}$ are mutually relevant, so they do not sum to shorter cycles by Lemma~\ref{Lem:mut_rel_2}. Thus, there exists an equal length cycle $B_j$, $|B_j|=|C_1|$, in the expansion of $\calX_{j-1}\cup\{C_2\}$ with respect to $\calM$.
	For step 3, we observe that $B_j$ is not contained in the right-hand side of Eq.~\eqref{eq:DICE_alg_Exs} because it is not a smaller cycle. 
	Therefore, there exists an additional cycle $X_j$ such that $B_j\in\expand{\calM}{X_j}$ which cancels with the $B_j$ in the expansion of $\calX_{j-1}\cup\{C_2\}$ in the left-hand side of Eq.~\eqref{eq:DICE_alg_Exs}. 
%	Therefore, there exists an additional cycle $X_j$ where $B_j\in\expand{\calM}{X_j}$ cancels with $B_j\in\expand{\calM}{\calX_{j-1}\cup\{C_2\}}$ in the left-hand side of Eq.~\eqref{eq:DICE_alg_Exs}. 
	
	We show by an inductive argument the cycles $\calX_j\cup\{C_2\}$ are in the same {\tt pi} class via the transitive closure of the {\tt pi} relations given by $\calM$.
	Initially, $\calX_0\cup\{C_2\}$ is a single cycle, so there is nothing to prove.
%	Initially, $\calX_0\cup\{C_2\}=\{C_2\}$ is mutually relevant because $C_2$ is relevant. 
	For the inductive step, let $\calX_{j-1}\cup\{C_2\}$ belong to the same {\tt pi} class by %be {\tt pi}-interchangeable from 
	the transitive closure of {\tt pi} relations from $\calM$. By construction, $B_j$ is an equal length cycle in the expansion of both $\calX_{j-1}\cup\{C_2\}$ and $X_j$ in $\calM$. Thus, $B_j$ is {\tt pi}-related to a cycle in $\calX_{j-1}\cup\{C_2\}$ and to $X_j$ via $\calM$. Therefore, the cycles $\calX_j\cup\{C_2\}=\calX_{j-1}\cup\{C_2\}\cup\{X_j\}$ are in the same {\tt pi} class by %{\tt pi}-interchangeable using 
	the transitive closure of {\tt pi} relations from $\calM$. 
	
	The stopping condition of this procedure is that $C_1\in\calX_j\cup\{C_2\}$. Hence, $C_1{\pInt} C_2$ by the transitive closure of the {\tt pi} relations from $\calM$. The cycles $C_1,C_2$ are arbitrary, so $\calM$ can be used to identify all {\tt pi} relations.%, and the direct-interchangeability classes.
%	Thus, any direct-interchangeability relation is given by the transitive closure those from $\calM$ as desired.
%	We complete our proof by showing (by induction) that for each iteration $j$ the cycles $\calX_j\cup\{C_2\}$ are directly-interchangeable via the transitive closure of Eq.~\eqref{Eq:DirEx_Alg_Th}. 
%	We complete our proof by showing by induction that for each iteration $j$ the cycles $\calX_j\cup\{C_2\}$ are directly-interchangeable via the transitive closure of the direct-interchangeable relations from $\calM$. 
%	As a base case, $\calX_0\cup\{C_2\}=\{C_2\}$ is mutually relevant because $C_2$ is relevant. For the inductive step, we assume the cycles $\calX_{j-1}\cup\{C_2\}$ are directly-interchangeable. $M_j$ belongs to the expansion of $\calX_{j-1}\cup\{C_2\}$ in $\calM$, so $M_j$ is directly interchangeable for the cycles $\calX_{j-1}$ by Eq.~\eqref{Eq:DirEx_Alg_Th}. Similarly, $M_j$ and $X_j$ are directly-interchangeable by Eq.~\eqref{Eq:DirEx_Alg_Th}, since $M_j$ belongs to the expansion of $X_j$ in $\calM$. Thus, using $M_j$ as a bridge, the cycles $\calX_j\cup\{C_2\} = \calX_{j-1}\cup\{C_2,X_j\}$ are directly-interchangeable by the transitive closure of Eq.~\eqref{Eq:DirEx_Alg_Th}. 
%	Finally, using the stopping condition $C_1\in\calX_j$, we conclude $C_1\DirEx C_2$ by the transitive closure of Eq.~\eqref{Eq:DirEx_Alg_Th}.
\end{proof}

Last, we prove we can always exchange cycles between distinct MCBs.
%Last, we prove that for two given minimum cycle bases $\calM_1,\calM_2$ we can always exchange a cycle $C_2\in\calM_2\setminus\calM_1$ for another cycle $C_1\in\calM_1\setminus\calM_2$ in $\calM_1$.

\begin{proof}[Proof of Lemma~\ref{Lem:DICE_Irreducible}]
	Let $\calM_1,\calM_2$ be distinct MCBs and $C_2\in\calM_2\setminus\calM_1$. We want to construct a cycle $C_1\in\calM_1\setminus\calM_2$ such that $(\calM_1\setminus\{C_1\})\cup\{C_2\}$ is an MCB.
	
	We expand $C_2$ in $\calM_1$ as follows
	\begin{equation}
		C_2 = \bigoplus_{{\substack{C'\in\expand{\calM_1}{C_2}\\|C'|=|C_2|}}}C'
		\oplus
		\bigoplus_{{\substack{C''\in\expand{\calM_1}{C_2}\\|C''|<|C_2|}}}C''.
	\end{equation} 
	Next, we move the equal length cycles to the left-hand side.
	\begin{equation}
		C_2 \oplus \bigoplus_{\substack{C'\in\expand{\calM_1}{C_2}\\|C'|=|C_2|}}C'
		=
		\bigoplus_{\substack{C''\in\expand{\calM_1}{C_2}\\|C''|<|C_2|}}C''.
	\end{equation}
	Importantly, the cycles of length $|C_2|$ in $\{C_2\}\cup\expand{\calM_1}{C_2}$ sum to shorter cycles, so they are not mutually relevant by Lemma~\ref{Lem:mut_rel_2}.
	
	By definition, the cycles of length $|C_2|$ in $\calM_2$ are mutually relevant. A subset of a mutually relevant set is mutually relevant, so the cycles of length $|C_2|$ in $\{C_2\}\cup\expand{\calM_1}{C_2}$ are not contained in $\calM_2$. In particular, there exists $C_1\in\expand{\calM_1}{C_2}$ of length $|C_1|=|C_2|$ such that $C_1\in\calM_1\setminus\calM_2$. By Lemma~\ref{Lem:Relevant_MCB_Test}, Statement 2, we may exchange $C_2$ for $C_1$ in $\calM_1$ as desired.
\end{proof}


\subsection{Short loop-interchangeability}
\label{app:TICE_Proofs}
Here, we prove the theoretical results for the {\tt sli} relation from Section~\ref{sub:TICE_def}.
Recall, two equal length cycles $C_1$ and $C_2$ satisfy $C_1\slInt C_2$ if there exist shorter cycles $\calY$, $|C|{<}|C_1|$ for $C{\in}\calY$, such that
\begin{equation}
	C_1 = C_2 \oplus\bigoplus_{C\in\calY}C.\label{eq:TrEx_def2}
\end{equation}
First, we prove {\tt sli} is an equivalence relation.

\begin{proof}[Proof of Lemma~\ref{Lem:TrInt_Eq_Rel}]
	The reflexive property is immediate by $C=C\oplus\emptyset$.
	
	(symmetry) Let $C_1\slInt C_2$ via shorter cycles $\calY$. By adding cycles $\calY$ to both sides of Eq.~\eqref{eq:TrEx_def2} we find that $C_2\slInt C_1$:
	\begin{equation*}
		C_1 = C_2\oplus\bigoplus_{C\in\calY} C \enspace\Rightarrow\enspace 
		C_2 = C_1\oplus\bigoplus_{C\in\calY}C.
	\end{equation*}
	
	(transitivity) Let $C_1\slInt C_2$ and $C_2\slInt C_3$ by sets of smaller cycles $\calY$ and $\calY'$ respectively,
	\begin{align*}
		C_1 = C_2\oplus\bigoplus_{C\in\calY}C,\qquad
		C_2 = C_3\oplus\bigoplus_{\mathclap{C'\in\calY'}}C'.
	\end{align*}
	Substituting the right-hand side of the second equation for $C_2$ in the first equation gives
	\begin{equation*}
		C_1 = \left(C_3\oplus\bigoplus_{\mathclap{C'\in\calY'}}C'\right)\oplus\bigoplus_{C\in\calY}C=C_3\oplus\bigoplus_{\mathclap{C''\in\calY\symdiff\calY'}}C''.
	\end{equation*}
	Hence, $C_1{\slInt}\,C_3$ via shorter cycles $\calY\symdiff\calY'$.% as desired.
\end{proof}

%Next, we prove a pair of cycles are {\tt sli}-interchangeable if and only if their expansions with respect to a given MCB differ by shorter cycles.
Next, we prove $C_1$ and $C_2$ satisfy $C_1\slInt C_2$ if and only if their expansions with respect to a given MCB differ by shorter cycles.

\begin{proof}[Proof of Lemma~\ref{Lem:TrEx_Alg}]
	($\Leftarrow$) Let $\calM$ be an MCB and $C_1,C_2\in\RelCyc$ such that $|C_1|=|C_2|$. Let the expansions of $C_1$ and $C_2$ in $\calM$ differ by shorter cycles, i.e., there exist cycles $\calY\subseteq\calM$, $|C|<|C_1|$ for $C\in\calY$, such that
	\begin{equation}
		\expand{\calM}{C_1}=\expand{\calM}{C_2}\symdiff\calY.
		\label{eq:MCB_TrEx}
	\end{equation}
	By summing over cycles in Eq.~\eqref{eq:MCB_TrEx} we obtain
	\begin{equation*}
		C_1 = C_2\oplus\bigoplus_{C\in\calY} C.
	\end{equation*}
	Thus, $C_1{\slInt}\,C_2$ via shorter cycles $\calY$ as desired.
	
	($\Rightarrow$) Let $C_1\slInt C_2$ by shorter cycles $\calY$ as in Eq.~\eqref{eq:TrEx_def2}.
%	\begin{equation}
%		C_1 = C_2\oplus\bigoplus_{C\in\calY}C.
%	\end{equation}
	Then, by linearity of the expansion operator applied to Eq.~\eqref{eq:TrEx_def2} we obtain
	\begin{equation}
		\begin{split}
			\expand{\calM}{C_1} &= \expand{\calM}{C_2} \symdiff \expand{\calM}{\bigoplus_{C\in\calY} C}\\
			&= \expand{\calM}{C_2} \symdiff \calY'
		\end{split}
	\end{equation}
	where $\calY'\subseteq\calM$ is a set of shorter cycles in $\calM$ by Lemma~\ref{Lem:SmallCycs2MCB}. Thus, the expansions of $C_1$ and $C_2$ in $\calM$ differ by shorter cycles $\calY'$ as desired.	
\end{proof}



\section{Proofs for Section~\ref{sub:pair_intersect}}
\label{app:intersect_proofs}
%{\color{red}Make these proofs more visual}

%\begin{figure*}
%	\centering
%	% to get figure run "pdflatex --shell-escape cyc_overlap.tex" for a png
%	\includegraphics{TikzFigs/cyc_overlap/cyc_overlap.png}
%	\caption{
%		\add{Keep (b), (c) in main text, try to do (d),(e),(f) inline with shorter text. Less text, more figures. Potentially, f into smaller figures.} {\color{red} Length of $C_2$ has changed!}
%		\textbf{(a)} Shared embedding $\varphi(x)=(d(u,x),d(v,x))$ of two cycles $C_1,C_2$ in a minimum cycle basis $\calM$ using a pair of nodes $u,v$ belonging to both $C_1$ and $C_2$ where $d(u,v)<|C_1|/2$. Blue nodes belong to $C_1$, red nodes belong to $C_2$, and purple nodes belong to both. Edges are colored similarly. Mixed blue-red nodes represent a pair of nodes, $x_1\in C_1$, $x_2\in C_2$, that share an embedding $\varphi(x_1)=\varphi(x_2)$. \textbf{(b)}~Another representation where nodes in $C_1$ are placed on the left, nodes in $C_2$ on the right, and their common nodes in the center. \textbf{(c)} Exchange of the cycle $C_1'$ for $C_1$ such that $C_1'$ and $C_2$ intersect over a single path. \textbf{(d)} Degenerate case where $C_1$ and $C_2$ intersect over a single node $u$. \textbf{(e)} Degenerate case where the nodes $u,v$ are antipodal in $C_1$, so they cannot be used to construct an embedding as in {(a)}. \textbf{(f)} A pair of cycles $C_1,C_2$ that do not intersect in the same way as in (b); these cycles are not relevant.}
%	\label{fig:cyc_pair_overlap}
%\end{figure*}



We prove these results using a distance-based approach. 
The distance $d(u,v)$ between nodes of a graph $G$ is the length of a shortest path connecting $u$ and $v$.
Recall, a path $P$ belonging to a relevant cycle $C$ of length $|P|\le |C|/2$ is a shortest path~\cite{vismaraUnion1997} -- see Lemma~\ref{lem:shortest_paths}. 
A pair of nodes $u$ and $v$ in $C$ are connected by two paths in $C$, one for each orientation of $C$. 
The shorter of these paths, $P_{u,v}$, is bounded in length by $|P_{u,v}|\le |C|/2$. 
Thus, every pair of nodes in $C$ are connected by a shortest path of $G$ in $C$. We restate this as a result:

\begin{corollary}[distance preservation]
	Let $C$ be a relevant cycle of a graph $G$ and let $u,v$ be nodes belonging to $C$. The distance between $u$ and $v$ with respect to $C$ is equal to their distance with respect to $G$.
	\label{Cor:dist_pres}
\end{corollary}

Using Corollary~\ref{Cor:dist_pres}, we employ the notion of resolving sets. A resolving set~\cite{tillquist2023getting} is a set of nodes $R=\{r_1,\hdots,r_k\}\subseteq V$ such that each node $x\in V$ is uniquely identified by its distances to the nodes $R$. For general graphs, an efficient resolving set is hard to construct. However, a resolving set of a cycle requires only two nodes.

%\begin{figure*}[t]
%	\begin{tikzpicture}[scale=1]
%		% plot axes
%		\draw[very thick,->] (0,-.5) -- (0,4) node[label=90:{$d(v,x)$}] {};
%		\draw[very thick,->] (-.5,0) -- (4,0) node[label=0:{$d(u,x)$}] {};
%		
%		% blue circle
%		\def\circS{.16}
%		\def\bx{-8.5}
%		\def\by{3.5}
%		\foreach \l in {1,2,3,4,5,6}
%		{\node[namednode,fill=white] (b\l) at ({\bx-6*\circS*cos((\l-1)*60)},{\by+6*\circS*sin((\l-1)*60)}) {\l};}
%		\node[draw,very thick,circle,minimum size=15pt, label=left:{$u$}] () at (b1) {};
%		\node[draw,very thick,circle,minimum size=15pt, label=60:{$v$}] () at (b3) {};
%		\draw[very thick,blue](b1)--(b2)--(b3)--(b4)--(b5)--(b6)--(b1);
%		
%		% green circle
%		\def\gx{-8.5}
%		\def\gy{.5}
%		\foreach \l in {1,2,3,4,5,6,7}
%		{\node[namednode,fill=white] (g\l) at ({\gx-7*\circS*cos((\l-1)/7*360)},{\gy+7*\circS*sin((\l-1)/7*360)}) {\l};}
%		\node[draw,very thick, circle, minimum size=15pt,label=180:{$u$}] () at (g1) {};
%		\node[draw,very thick, circle, minimum size=15pt,label={180-360/7}:{$v$}] () at (g2) {};
%		\draw[very thick,black!40!green,very thick] (g1)--(g2)--(g3)--(g4)--(g5)--(g6)--(g7)--(g1);
%		
%		% red circle
%		\def\circS{.16}
%		\def\rx{-6}
%		\def\ry{2}
%		\foreach \l in {1,2,3,4,5,6}
%		{\node[namednode,fill=white] (r\l) at ({\rx-6*\circS*cos((\l-1)*60)},{\ry+6*\circS*sin((\l-1)*60)}) {\l};}
%		\node[draw,very thick,circle,minimum size=15pt, label=left:{$u$}] () at (r1) {};
%		\node[draw,very thick,circle,minimum size=15pt, label=right:{$v$}] () at (r4) {};
%		\draw[very thick,red](r1)--(r2)--(r3)--(r4)--(r5)--(r6)--(r1);
%		
%		% transformation arrow
%		\draw[arrows={-Stealth[length=12pt, inset=0pt, angle'=25]},line width = 1.75pt] (-3.75,2) -- node[above] {$\varphi:V\rightarrow\mathbb{Z}^2\ \ $} (-1.25,2);
%		
%		% red "rectangle"
%		\draw[thick, red] (0,3) node[regnode,red]{} --
%		(1,2) node[regnode,minimum size=7pt,red]{} --
%		(2,1) node[regnode,minimum size=7pt,red]{} -- 
%		(3,0) node[regnode,red]{};
%		% blue rectangle
%		\draw[thick, blue] (0,2) node[regnode,blue]{} -- 
%		(1,1) node[regnode,blue]{} -- 
%		(2,0) node[regnode,blue]{} -- 
%		(3,1) node[regnode,blue]{} -- 
%		(2,2) node[regnode,blue]{} -- 
%		(1,3) node[regnode,blue]{} -- cycle;
%		% green rectangle
%		\draw[thick, black!40!green] (0,1) node[regnode,black!40!green]{} -- 
%		(1,0) node[regnode,black!40!green]{} -- 
%		(2,1) node[regnode,black!40!green]{} -- 
%		(3,2) node[regnode,black!40!green]{} -- (3.5,2.5) --
%		(3,3) node[regnode,black!40!green]{} -- (2.5,3.5) --
%		(2,3) node[regnode,black!40!green]{} -- 
%		(1,2) node[regnode,black!40!green]{} -- cycle;
%		
%		% axis markings
%		\node[anchor=east] () at (-.5,2) {$u$};
%		\draw[] (-.5,2) -- (-.2,2);
%		\draw[] (-.2,3) -- (-.5,3) -- (-.5,1) -- (-.2,1);
%		\node[anchor=north] () at (2,-.5) {$v$};
%		\draw[] (2,-.5) -- (2,-.2);
%		\draw[] (3,-.2) -- (3,-.5) -- (1,-.5) -- (1,-.2);
%	\end{tikzpicture}
%	\caption{
%		(left) three cycle graphs with labeled nodes $u$ and $v$. (right) Embedding of the cycles using the function $\varphi(x)=(d(u,x),d(v,x))$. The resulting rectangles are color-coded to match the corresponding cycles. 
%		The embedding of the nodes $v$ and $u$ are highlighted as the x and y-intercepts of the rectangle respectively. Neighboring nodes $p,q$ map to consecutive points $\varphi(p)$ and $\varphi(q)$ along the rectangle. $\varphi$ is one-to-one except for the red cycle which does not satisfy $d(u,v)<|C|/2$.
%	}
%	\label{fig:Res_Cycle}
%\end{figure*}

\begin{lemma}
	Let $C$ be a cycle graph and $u,v$ nodes in $C$ such that $d(u,v)<|C|/2$. Then, every node $x$ in $C$ is uniquely identified by its distances to $u$ and $v$, i.e., the function
	\begin{equation}
		\varphi(x) = (d(x,u),d(x,v)),\quad x\in V(C)\label{eq:phi_def}
	\end{equation}
	is one-to-one.
	\label{Lem:cyc_embed}
\end{lemma}

\begin{proof}
%	{\color{red} ``x axis is distance to u. y axis is distance to v. $\varphi$ will map the if u and v are antipodal, then due to symmetry then the loop will collapse.'' Say less, talk more to diagram. Don't talk about corners on nodes... in Figure...}
	Let $u,v$ be nodes belonging to a cycle graph $C$ and $\varphi$ given by Eq.~\eqref{eq:phi_def}. We want to show $\varphi$ is bijective. We also want to describe the shape given by $\varphi(C)$ which will be useful in the proof of Theorem~\ref{Thm:MCB_Intersect} as well. When plotting $\varphi$, the x and y coordinates correspond to the distance of a given node to $u$ and $v$ respectively. This is shown for an even cycle in Figure~\ref{fig:Res_Set_Lem1}.
	\begin{figure}[h]
		\includegraphics{Res_Set_Lem1.pdf}
		\caption{Even cycle for Lemma~\ref{Lem:cyc_embed}.}
		\label{fig:Res_Set_Lem1}
	\end{figure}
	\noindent We observe that $\varphi$ maps this cycle to the integer coordinates of a rectangle.
	
	In contrast, $\varphi$ maps an odd cycle to a rectangle with two of its corners cut off as shown in Figure~\ref{fig:Res_Set_Lem2} below.
	\begin{figure}[h]
		\includegraphics{Res_Set_Lem2.pdf}
		\caption{Odd cycle for Lemma~\ref{Lem:cyc_embed}.}
		\label{fig:Res_Set_Lem2}
	\end{figure}
	
	In either case, the function $\varphi$ maps our cycle to the integer coordinates of a polygon. 
	This shape is not self-intersecting, so $\varphi$ is one-to-one.
	The only exception to this is if the nodes $u$ and $v$ are antipodal ($d(u,v)=|C|/2$) as shown in Figure~\ref{fig:Res_Set_Lem3} below.
	\begin{figure}[h]
		\includegraphics{Res_Set_Lem3.pdf}
		\caption{Degenerate case for Lemma~\ref{Lem:cyc_embed}.}
		\label{fig:Res_Set_Lem3}
	\end{figure}
	\noindent In this case, the symmetries of $u$ and $v$ cause the mapping $\varphi(C)$ to collapse from a rectangle into a line. 
	Because this line doubles back on itself the function $\varphi$ is not one-to-one when $d(u,v)=|C|/2$.
\end{proof}

Next, we use the resolving set representation of cycles to describe the intersection format of two cycles in an MCB.

%The function $\varphi$ in Lemma~\ref{Lem:cyc_embed} maps a cycle to the integer coordinates of a rectangle. 
%See Figure~\ref{fig:Res_Cycle}.
%The condition that the nodes $u,v$ are not antipodal, $d(u,v)<|C|/2$, is there to ensure this rectangle has non-zero width.

%Consider a pair of cycles $C_1,C_2$ in a minimum cycle basis $\calM$ that share two nodes $u,v$. 
%By Lemma~\ref{Lem:cyc_embed}, we construct a map $\varphi$ using $u$ and $v$ to create a shared embedding of $C_1$ and $C_2$. 
%See Figure~\ref{fig:cyc_pair_overlap}~(a). 
%Most notably, this shared embedding provides a common orientation of $C_1$ and $C_2$, so that their paths of intersection appear with the same orientation and in the same order. 
%More details about the format in which cycles in a minimum cycle basis can intersect are given in Theorem~\ref{Thm:MCB_Intersect}.
%This result does not extend to arbitrary pairs of cycles. 
%The (non-relevant) cycles in Figure~\ref{fig:cyc_pair_overlap}~(f) lack a common orientation, so they appear twisted.

\begin{proof}[Proof of Theorem~\ref{Thm:MCB_Intersect}]
	%	Without loss of generality, let $|C_1|\le|C_2|$. 
	Let $C_1,C_2,|C_1|\le|C_2|,$ be cycles in an MCB $\calM$. 
	We split this proof into three cases, a main case where the methods in Lemma~\ref{Lem:cyc_embed} apply, and two degenerate cases.
	
	We remark that a path can be a single node, e.g., the path $P_3$ in Figure~\ref{fig:intersect_pf_2}. 
	The length of a path is the number of edges it has. A path with one node has length 0.
%	Here, we include single nodes as paths, e.g., the path $P_3$ in Figure~\ref{fig:cyc_pair_overlap} (b). The length of a path is the number of edges it has; a single node is a path of length 0. First, we address the main case where the methods in Lemma~\ref{Lem:cyc_embed} apply, then we address two degenerate cases.
	
	
	%	We traverse $C_1$ and $C_2$ according to counter-clockwise orientation of the map $\varphi$. The paths $P_i$ are maximal consecutive sets of nodes shared by $C_1$ and $C_2$ in this orientation. The paths $P_i$ appear with the same order and orientation because of the shared embedding of $C_1$ and $C_2$. What remains to show is that the paths $Q_i^j$ satisfy the length requirements in 
	
	
	
	\emph{Main case.} Let $u$ and $v$ be nodes belonging to both $C_1$ and $C_2$ such that $d(u,v)\le |C_1|/2$. 
	Then, we may use the function $\varphi(x)=(d(u,x),d(v,x))$ to embed $C_1$ and $C_2$ as overlapping rectangles $\varphi(C_1)$ and $\varphi(C_2)$ (or rectangles with cut corners as in Figure~\ref{fig:Res_Set_Lem2}). 
	By Lemma~\ref{Lem:cyc_embed}, the mapping $\varphi$ is one-to-one when its domain is restricted to $C_1$ or to $C_2$.
	This shared embedding is shown in Figure~\ref{fig:intersect_pf_1}.
%	Then, we may embed both of these cycles using the function $\varphi(x)=(d(u,x),d(v,x))$. 
%	By Lemma~\ref{Lem:cyc_embed}, $\varphi$ is one-to-one when its domain is restricted to either $C_1$ or $C_2$. 
%	See Figure~\ref{fig:intersect_pf_1} for a visualization of this embedding.
	\begin{figure}[h]
		\includegraphics[]{Intersection_Proof1.pdf}
		\caption{Embedding of $C_1$ and $C_2$ for the main case of Theorem~\ref{Thm:MCB_Intersect}.}
		\label{fig:intersect_pf_1}
	\end{figure}
	In Figure~\ref{fig:intersect_pf_1}
	blue nodes belong to $C_1$, red nodes to $C_2$, and purple nodes to both. 
	Dual blue and red points denote a pair of nodes $x_1\in V(C_1)$, $x_2\in V(C_2)$ that share the same embedding $\varphi(x_1)=\varphi(x_2)$.
	Edges are colored similarly.
	
	The mapping $\varphi$ couples the cycles $C_1$ and $C_2$ by stacking them on top of each other. 
%	In particular, we traverse $C_1$ and $C_2$ counter-clockwise in the mapping $\varphi(C_1)$ and $\varphi(C_2)$. 
	We orient $C_1$ and $C_2$ by traversing $\varphi(C_1)$ and $\varphi(C_2)$ in the counter-clockwise direction.
	If $C_1$ and $C_2$ are equal in length, then the rectangles $\varphi(C_1)$ and $\varphi(C_2)$ overlap completely. 
	If $C_2$ is longer than $C_1$, then the rectangle $\varphi(C_2)$ extends past $\varphi(C_1)$, as in Figure~\ref{fig:intersect_pf_1}.
	
%	See the solid blue and red lines in the top right of Figure~\ref{fig:intersect_pf_1}.
%	These trajectories overlap, possibly except at the northeast edge of each rectangle. See the solid blue and red lines in Figure~\ref{fig:intersect_pf_1}.
	
	We use this shared embedding to decompose $C_1$ and $C_2$ into paths, starting with their intersection paths $P_1,\hdots,P_k$.
%	First, we discuss the structure of the intersection paths $P_1,\hdots,P_k$ between $C_1$ and $C_2$.
	Since we traverse both $\varphi(C_1)$ and $\varphi(C_2)$ in the counter-clockwise direction, these intersection paths have the same order and orientation in $C_1$ and $C_2$.
%	These intersection paths appear with the same order and orientation between $C_1$ and $C_2$. 
	See Figure~\ref{fig:intersect_pf_2}.
	\begin{figure}[h]
	\includegraphics[]{Intersection_Proof2.pdf}
	\caption{Intersection paths of $C_1$ and $C_2$ and their directions in the main case of Theorem~\ref{Thm:MCB_Intersect}.}
	\label{fig:intersect_pf_2}
	\end{figure}
	In Figure~\ref{fig:intersect_pf_2},
	we mark the orientation of $P_1$ and $P_2$ using purple arrows. The orientation of $P_3$ is arbitrary, because it is a single point.
	We remark that we have not selected a starting point when traversing $\varphi(C_1)$ and $\varphi(C_2)$. Because of this, the order of the intersection paths is circular, e.g., in Figure~\ref{fig:intersect_pf_2} $P_2$ follows $P_1$, $P_3$ follows $P_2$, and $P_1$ follows $P_3$.
%	The orientations of the paths $P_1$ and $P_2$ are shown using purple arrows. As a single node, the orientation of $P_3$ is arbitrary. 
%	Since the starting point is not fixed, the order of these paths is cyclical:
%	$P_2$ follows $P_1$, $P_3$ follows $P_2$, and $P_1$ follows $P_3$.
	
	Next, we discuss the paths $Q_i^{(j)}$ that separate $C_1$ and $C_2$. All of the path pairs $Q_i^{(1)},Q_i^{(2)}$ fully overlap -- as in Figure~\ref{fig:intersect_pf_3} (a) -- unless $|C_1|<|C_2|$. If $C_1$ is shorter than $C_2$, then the pair that includes the northeast lines of the rectangles $\varphi(C_1),\varphi(C_2)$ do not overlap, as in Figure~\ref{fig:intersect_pf_3} (b).
	We prove three identities regarding the lengths of these pairs.
	\begin{figure}[h]
		\includegraphics[]{Intersection_Proof3a.pdf}
		\includegraphics[]{Intersection_Proof3b.pdf}
		\caption{Separated path pairs $Q_i^{(1)},Q_i^{(2)}$ for the main case of Theorem~\ref{Thm:MCB_Intersect}.}
		\label{fig:intersect_pf_3}
	\end{figure}
	
	\begin{enumerate}
		\item Let $|C_1|<|C_2|$, and order the path pair with different lengths last: $|Q_k^{(1)}| \ne |Q_k^{(2)}|$. This is the pair shown in Figure~\ref{fig:intersect_pf_3} (b). We will show this pair satisfies
		\begin{align}
			|Q_k^{(2)}| &= |Q_k^{(1)}|+|C_2|-|C_1|\label{eq:Int_Thm_QId1a}\\
			|Q_k^{(j)}| &\ge |C_j|/2,\quad j=1,2.\label{eq:Int_Thm_QId1b}
		\end{align}
		Eq.~\eqref{eq:Int_Thm_QId1a} is obvious because $Q_k^{(1)}$ and $Q_k^{(2)}$ account for the difference in lengths between $C_1$ and $C_2$:
		\begin{equation*}
			|Q_k^{(2)}| - |Q_k^{(1)}| = |C_2|-|C_1|
		\end{equation*}
		As for Eq.~\eqref{eq:Int_Thm_QId1b}, consider the cycle $C_2'$ obtained by swapping the path $Q_k^{(2)}$ in $C_2$ for $Q_k^{(1)}$,
		\begin{equation*}
			C_2'= C_2\oplus C_k^e,\quad C_k^e = Q_k^{(1)}\cup Q_k^{(2)}.
		\end{equation*}
		$C_2$ cannot be broken into shorter cycles because it is relevant. $C_2'$ is shorter than $C_2$ because $Q_k^{(1)}$ is shorter than $Q_k^{(2)}$. Hence, the cycle $C_k^e$ used to exchange $C_2$ for $C_2'$ satisfies
		\begin{equation}
			|C_k^e| = |Q_k^{(1)}|+|Q_k^{(2)}| \ge |C_2|.
			\label{eq:Int_Thm_Exk_Ineq}
		\end{equation}
		Eq.~\eqref{eq:Int_Thm_QId1b} is obtained by combining Eqs.~\eqref{eq:Int_Thm_QId1a} and~\eqref{eq:Int_Thm_Exk_Ineq}.
%		\add{Second identity}
		
		\item Let $|C_1|=|C_2|$. We will show there exists a unique large path pair satisfying
		\begin{equation}
			|Q_k^{(1)}|,|Q_k^{(2)}|\ge |C_1|/2.
		\end{equation}
		This pair is ordered last to match the previous item. 
		
		Suppose towards a contradiction that each path pair is shorter,
		\begin{equation*}
			|Q_i^{(1)}|=|Q_i^{(2)}| < |C_1|/2,
		\end{equation*}
		for $i=1,\hdots,k$. Then $C_2$ can be obtained from $C_1$ by swapping each path $Q_i^{(1)}$ for $Q_i^{(2)},$
		\begin{equation*}
			\begin{split}
			C_2 &= C_1\oplus(C_1^e\oplus\hdots\oplus C_k^e)\\
			C_i^e &= Q_i^{(1)}\cup Q_i^{(2)},\quad i=1,\hdots,k.
			\end{split}
		\end{equation*}
		The cycles $C_i^e$ are shorter
		\begin{equation*}
			|C_i^e| = |Q_i^{(1)}|+|Q_i^{(2)}|<|C_1|.
		\end{equation*}
		Thus, $C_1\slInt C_2$. This is a contradiction because at most one cycle in a {\tt sli} class may belong to an MCB by Corollary~\ref{Cor:TrInt_Exchanges}.
		
		\item Last, we show the remaining path pairs are bounded above in length by
		\begin{equation}
			|Q_i^{(1)}|,|Q_i^{(2)}|<|C_1|/2, 
			\label{eq:Int_Pf_Short}
		\end{equation}
		for  $i=1,\hdots,k-1$. 
		Consider a $Q_i^{(j)}$ that lies on a shortest path from $u$ to $v$, such as $Q_1^{(1)}$ or $Q_1^{(2)}$ in Figure~\ref{fig:intersect_pf_3}~(a). Such a path is bounded in length by
%		If a path $Q_i^{(j)}$ lies on a shortest path from $u$ to $v$, then its length is bounded by
		\begin{equation*}
			|Q_i^{(j)}|\le d(u,v) < |C_1|/2.
		\end{equation*}
		Next, consider a $Q_i^{(j)}$ that does not lie on a shortest path from $u$ to $v$, e.g., $Q_2^{(1)}$ and $Q_2^{(2)}$ in Figure~\ref{fig:intersect_pf_3}~(a). Then, $Q_i^{(j)}, i\neq k$ is bounded in length by
		\begin{equation*}
			\begin{split}
			|Q_i^{(j)}| = |Q_i^{(1)}| &\le |C_1| - |Q_k^{(1)}|-d(u,v) \\&< |C_1|/2
			\end{split}
		\end{equation*}
		because (i) $|Q_i^{(1)}|=|Q_i^{(2)}|$ for $i\neq k$, (ii) the path $Q_k^{(1)},|Q_k^{(1)}|>|C_1|/2$ also does not lie on a shortest path from $u$ to $v$, and (iii) $d(u,v)>0$ because $u\neq v$.
		In both cases $Q_i^{(j)}$ satisfies the length bound in Eq.~\eqref{eq:Int_Pf_Short}.
%		where $|Q_k^{(1)}|\ge|C_1|/2$ by Eq.~\eqref{eq:Int_Thm_QId1b} and $d(u,v)>0$ because $u\ne v$.
	\end{enumerate}
	
	In summary, these paths satisfy the following:
	\begin{align}
		|Q^{(1)}_i| &\!=\! |Q^{(2)}_i|, &i=1,\!...,k-1,
		\label{eq:sm_path_equal_pf}\\
		|Q^{(j)}_i| &\!<\! |C_1|/2,\! &i{=}1,\!...,\!k{-}1, j{=}1,\!2,
		\label{eq:sm_path_len_pf}\\
		|Q^{(2)}_k| &\!=\! \mathrlap{|Q^{(1)}_k|+|C_2|-|C_1|,}
		\label{eq:lg_path_equal_pf}\\
		|Q^{(j)}_k| &\!\ge\! |C_j|/2,\! &j=1,2.
		\label{eq:lg_path_len_pf}
	\end{align}
	We note that, Eq.~\eqref{eq:sm_path_len_pf} is a strict inequality, which is a stronger result than we intended to show. 
	This observation will be useful in the proof of Theorem~\ref{Thm:MCB_Pair_Swap}. 
	This inequality does not extend to degenerate case 2 below.
	
	
	\emph{Degenerate case 1.} Let $C_1$ and $C_2$ share only a single node $u$. See Figure~\ref{fig:intersect_pf_deg1}. 
	\begin{figure}[h]
		\includegraphics{Intersection_ProofDeg1.pdf}
		\caption{Degenerate case 1 for Theorem~\ref{Thm:MCB_Intersect}.}
		\label{fig:intersect_pf_deg1}
	\end{figure}
	Then they overlap over a single path $P_1=(u)$. They are separated by a pair of paths $Q_1^{(1)}$ and $Q_1^{(2)}$ connecting $u$ to itself. These paths are of length $|Q_1^{(1)}|=|C_1|$ and $|Q_1^{(2)}|=|C_2$ which satisfies the required length requirements of the larger path pairs $Q_k^{(1)}$, $Q_k^{(2)}$.
%	, the path $Q_1^{(1)}$ joining $P_1$ into a loop is a traversal of $C_1$ starting and ending at $u$. For $C_2$, we define the path $Q_1^{(2)}$ similarly. Thus, $|Q_1^{(1)}|=|C_1|$ and $|Q_1^{(2)}|=|C_2|$, satisfying Eqs.~\eqref{eq:lg_path_equal} and~\eqref{eq:lg_path_len}.
	
	
	\emph{Degenerate case 2.} Let $C_1$ and $C_2$ share two nodes $u$ and $v$ which are antipodal in $C_1$, i.e., $d(u,v)=|C_1|/2$. %These nodes do not satisfy the assumption $d(u,v)<|C_1|/2$ in Lemma~\ref{Cor:dist_pres}. Nevertheless, $C_1$ and $C_2$ still satisfy the desired format. 
	See Figure~\ref{fig:intersect_pf_deg2}. 
	\begin{figure}[h]
		\includegraphics{Intersection_ProofDeg2.pdf}
		\caption{Degenerate case 2 for Theorem~\ref{Thm:MCB_Intersect}.}
		\label{fig:intersect_pf_deg2}
	\end{figure}
	In this case, $C_1$ and $C_2$ overlap over two paths consisting of single nodes $P_1=(u)$ and $P_2=(v)$. $C_1$ is composed of two shortest paths from $u$ to $v$. Label these paths $Q_1^{(1)}$ and $Q_2^{(1)}$ arbitrarily. Also, $C_2$ has a shortest path connecting $u$ and $v$. Label this path $Q_1^{(2)}$ where $|Q_1^{(1)}|=|Q_1^{(2)}|=|C_1|/2$ satisfying Eqs.~\eqref{eq:sm_path_equal} and~\eqref{eq:sm_path_len}. Let $Q_2^{(2)}$ be the other path connecting $u$ and $v$ in $C_2$. This path has length $|Q_2^{(2)}| = |C_2| - |C_1|/2$ which satisfies Eqs.~\eqref{eq:lg_path_equal} and~\eqref{eq:lg_path_len}.
\end{proof}

Next, we prove that if $C_1$ and $C_2$ belong to an MCB and they intersect over multiple paths, then we may exchange $C_1'$ for $C_1$ such that $C_1'$ and $C_2$ intersect over a single path.

\begin{proof}[Proof of Theorem~\ref{Thm:MCB_Pair_Swap}]
	
	Here, we assume the cycles $C_1,C_2$ intersect over multiple paths, otherwise, no modification is needed and we may set $C_1'=C_1$. By Theorem~\ref{Thm:MCB_Intersect}, let $C_1$ and $C_2$ intersect over paths $P_1,\hdots,P_k$ and deviate over path pairs $Q_i^{(j)}$, $j=1,2$, for $i=1,\hdots,k$. We split the proof into two cases. Note, we do not assume $|C_1|<|C_2|$.%As before, we split the proof into cases.
	
	\emph{Main case.} 
	Let $C_1$ and $C_2$ belonging to an MCB $\calM$ share two nodes $u$ and $v$ of distance $d(u,v)<|C_1|/2$. As mentioned in the main case in the proof of Theorem~\ref{Thm:MCB_Intersect}, the paths $Q_i^{(1)}$ separating $C_1$ from $C_2$ satisfy the length bound $|Q_i^{(1)}|<|C_1|/2$ for $i=1,\hdots,k-1$. Also by Theorem~\ref{Thm:MCB_Intersect}, the paths $Q_i^{(2)}$ satisfy $|Q_i^{(2)}|=|Q_i^{(1)}|$ for $i=1,\hdots,k-1$.
	Thus, we may exchange the paths in $C_1$ for those in $C_2$ by adding cycles
	\begin{equation}
		C_i^e = {\color{blue}Q_i^{(1)}}\cup {\color{red}Q_i^{(2)}},\quad i=1,\hdots,k-1,
	\end{equation}
	resulting in the cycle $C_1'$ shown in Figure~\ref{fig:MCB_Pair_Swap_Main} below.
	
	\begin{figure}[h]
		\includegraphics{postprocess_pfMain.pdf}
		\caption{Main case of Theorem~\ref{Thm:MCB_Pair_Swap}.}
		\label{fig:MCB_Pair_Swap_Main}
	\end{figure}
	
	Importantly, the exchange cycles are shorter than $C_1$
	\begin{equation}
		|C^e_i| = |Q_i^{(1)}|+|Q_i^{(2)}|<|C_1|,
	\end{equation}
	Thus, the cycle
	\begin{equation}
		C_1' = C_1\oplus(C_1^e\oplus\hdots\oplus C_{k-1}^e)
	\end{equation}
	satisfies $C_1'\slInt C_1$.
%	is {\tt sli}-interchangeable for $C_1$. 
	By Corollary~\ref{Cor:TrInt_Exchanges}, we may exchange $C_1'$ for $C_1$ to obtain a new MCB $\calM'=(\calM\setminus\{C_1\})\cup\{C_1'\}$. Moreover, $C_1'$ and $C_2$ intersect over a single larger path.
	
	\emph{Degenerate case.} 
	Let $C_1$ and $C_2$ belonging to an MCB $\calM$ intersect at two nodes $u$ and $v$ of distance $d(u,v)=|C_1|/2$. In this case, $C_1$ is composed of two shortest paths $Q_1^{(1)}$ and $Q_2^{(1)}$ joining $u$ to $v$. Also, $C_2$ has a shortest path $Q_1^{(2)}$ connecting $u$ to $v$. We consider two candidate cycles $C_1'$ and $C_1''$ obtained by switching one of the paths from $u$ to $v$ in $C_1$ for $Q_2^{(1)}$ in $C_2$. The cycles $C_1'$ and $C_1''$ and their relation to $C_1$ are shown in Figure~\ref{fig:MCB_Pair_Swap_Deg} below.
	
	\begin{figure}[h]
		\includegraphics{postprocess_pfDeg.pdf}
		\caption{Degenerate case of Theorem~\ref{Thm:MCB_Pair_Swap}.}
		\label{fig:MCB_Pair_Swap_Deg}
	\end{figure}
	
	Algebraically, $C_1'$ and $C_1''$ are given by
	\begin{align}
		C_1' &= C_1\oplus({\color{red}Q_1^{(2)}}\cup{\color{blue}Q_2^{(1)}}) ={\color{blue}Q_1^{(1)}}\cup {\color{red}Q_1^{(2)}},\\
		C_1'' &= C_1\oplus({\color{blue}Q_1^{(1)}}\cup {\color{red}Q_1^{(2)}})={\color{red}Q_1^{(2)}}\cup {\color{blue}Q_2^{(1)}},
	\end{align}
	and $C_1$ is equal to their sum
	\begin{equation}
		C_1 = C_1'\oplus C_1'' = {\color{blue}Q_1^{(1)}}\cup{\color{blue}Q_2^{(1)}}.
	\end{equation}
	
	We remark that, $C_1'$ and $C_1''$ deviate from $C_1$ by equal length cycles, so we can not exchange them for $C_1$ using the {\tt sli} relation as in the main case.
	However, by linearity of the expansion operator, we observe $C_1$ is in the expansion of $C_1'$ or $C_1''$
	\begin{equation*}
		\expandd{\calM}{C_1'}\symdiff\expandd{\calM}{C_1''} = \mathcal{E}_{\calM}(\underbrace{C_1'\oplus C_1''}_{C_1}) = \{C_1\}.
	\end{equation*}
	Without loss of generality, let $C_1\in\expandd{\calM}{C_1'}$. Then, by Lemma~\ref{Lem:Relevant_MCB_Test}, $\calM'=(\calM\setminus\{C_1\})\cup\{C_1'\}$ is an MCB. Furthermore, $C_1'$ and $C_2$ intersect over a single path $Q_2^{(1)}$, completing our proof.
\end{proof}







	\bibliographystyle{ieeetr}
	\bibliography{bibliography.bib}
	
	
	

\end{document}
